\documentclass[11pt]{ctexart}
\usepackage[a4paper,margin=1in]{geometry}
\usepackage{graphicx}
\usepackage{xcolor}
\usepackage{hyperref}
\definecolor{refgray}{gray}{0.45}
\title{\textbf{市场最新观点汇总}\\[0.4em]{\large 全球资金正从追逐风险转向重新定价“安全”，AI与能源超级周期重塑产业格局，但宏观利率与地缘政治风险仍构成核心约束。}}
\author{KC桌面 自动生成}
\date{260613}
\begin{document}
\maketitle
\tableofcontents
\newpage
\section{一页摘要}
\begin{itemize}
  \item 全球资金流向显示，资金正同时涌入美国科技股和短期债券，却在撤出中国A股和欧洲股票，市场并非risk-on，而是在重新定义“安全资产”。
  \item 亚洲工业超级周期（AI、数据中心、电网）正驱动全球资本开支结构性集中，2026年预期增速13\%，但增长几乎完全由少数领域贡献。
  \item 美联储6月会议大概率按兵不动，但将删除所有降息暗示，点阵图中位数显示利率可能维持至2027年底，市场对降息的定价需要根本性修正。
  \item 中国信贷总量超预期但结构疲弱，企业中长期贷款净减少，票据融资飙升，实体融资需求并未真实回暖。
  \item 中国电动车渗透率完成结构性跃迁，5月BEV渗透率42\%，总EV渗透率63\%，即使在乘用车市场同比降22\%的背景下仍逆势增长。
  \item 美元多头共识已达历史极端区间，NOM回溯20年数据发现，当美国经济惊喜指数突破60后，美元未来三个月下跌概率高达75\%。
  \item 韩国居民资产正从房地产向权益市场结构性迁移，2025年权益持有量同比增长48\%，KOSPI目标被上调至9000点。
  \item 欧洲石油巨头估值隐含的长期油价仅60美元/桶，远低于市场共识的78美元和远期期货的79美元，提供了不对称的风险收益窗口。
  \item AI在放射科的渗透正从“可选项”变为“必需品”，老龄化叠加医生供给缺口，结构性需求缺口只能靠AI填补。
  \item 厄尔尼诺已正式形成，极强事件概率从37\%跃升至63\%，糖是最敏感的品种，大豆和玉米的影响取决于时间窗口。
\end{itemize}
\section{宏观与利率：美联储按兵不动，利率曲线重新定价路径被低估}
\textbf{市场对短期通胀的定价已较充分，但真正未被消化的是利率曲线在多重路径下的重新定价——美国曲线陡化概率被系统性低估，而欧洲和英国的前端利率存在过度定价的加息风险。}

\begin{itemize}
  \item GS认为，美国前端利率市场定价了30-35个基点的加息风险，但更可能的情景是“长期按兵不动”，曲线陡化概率被低估。
  \item NOM预期6月FOMC将维持利率不变，但会后声明将删除所有关于“宽松倾向”的前瞻指引，点阵图中位数上调至3.625\%，暗示利率维持至2027年底。
  \item NOM指出，美联储新主席Warsh可能利用就业疲软信号打压市场加息预期，其沟通风格比点阵图更重要。
  \item 欧洲通胀的持续性而非短期油价飙升才是决定ECB政策路径的真正变量，欧元前端利率对油价敏感度已下降。
  \item 英国利率对油价敏感度最高（t-stat高达8.5），若油价缓和，英镑前端利率有更大下行空间。
  \item 加拿大加息溢价会慢慢消退，BoC按兵不动，增长与政策因素大致对冲。
  \item 中国利率市场方面，NOM指出央行已在悄悄收紧中期流动性，3月以来累计净回笼持续，政府债供给即将放量，流动性正常化传导加速。
  \item NOM维持做空3年期NDIRS的头寸，目标看向1.60\%，认为债市极度宽松的流动性环境正在系统性逆转。
\end{itemize}
\begin{figure}[htbp]
\centering
\includegraphics[width=0.88\linewidth]{figures/fig_015.jpg}
\caption{Exhibit 1 - Exhibit 1: As inflation has repriced lower, we see value to being long 10y real rates on beta vs nominals}
\end{figure}
\begin{figure}[htbp]
\centering
\includegraphics[width=0.88\linewidth]{figures/fig_016.jpg}
\caption{Exhibit 2 - Exhibit 2: CAD growth and policy pricing roughly offset each other currently Contributions from Macro PCA factors over past month, to Canadian government bond yields}
\end{figure}
\begin{figure}[htbp]
\centering
\includegraphics[width=0.88\linewidth]{figures/fig_017.jpg}
\caption{Exhibit 3 - Exhibit 3: UK rates have been the most sensitive to oil moves Volatility adjusted beta (t-stat) of daily changes in 2y1y swap rates for a \$1\textbackslash{}\%\$ change in oil prices (in local currency)}
\end{figure}
\begin{figure}[htbp]
\centering
\includegraphics[width=0.88\linewidth]{figures/fig_019.jpg}
\caption{Exhibit 4 - Exhibit 5: Divergence in policy outlooks should drive the widening further}
\end{figure}
\begin{figure}[htbp]
\centering
\includegraphics[width=0.88\linewidth]{figures/fig_020.jpg}
\caption{Exhibit 6 - Exhibit 6: Belly yields have continued to move higher in tandem with traded inflation}
\end{figure}
\begin{figure}[htbp]
\centering
\includegraphics[width=0.88\linewidth]{figures/fig_077.jpg}
\caption{Figure 5 - \#\# 4. Our view on liquidity in coming weeks: We continue to expect the PBoC to be more flexible with 7d OMOs to inject short-term liquidity amid higher government bond supply, as well as tax payment and month-end funding needs. As for medium- to long-term liqu}
\end{figure}
\begin{figure}[htbp]
\centering
\includegraphics[width=0.88\linewidth]{figures/fig_078.jpg}
\caption{Figure 10 - As shown in Figure 10, after accumulating massive (ultra) long-end CGBs/policy bank bonds in April and May, onshore funds reduced bond holdings in June. In the first week of June, funds were still adding 10y and above tenors while net selling the front-end and}
\end{figure}
{\small\color{refgray}\textbf{References:} [R003] GS：市场真正低估的不是通胀本身，而是利率曲线重新定价的路径; [R006] NOM：中国利率市场真正被低估的风险不是降息预期，而是流动性正常化的传导加速; [R010] NOM：美联储的“按兵不动”比降息或加息更值得市场重新定价}

\section{外汇与美元：美元多头共识已达峰值，历史规律指向未来三个月走弱}
\textbf{美元多头共识逼近危险临界点，当美国经济惊喜指数突破60后，美元未来三个月下跌概率高达75\%，市场可能站在美元情绪周期的拐点。}

\begin{itemize}
  \item NOM回溯23年数据发现，美国经济惊喜指数突破60共出现41次，随后三个月美元兑G10等权重指数平均回报-1.8\%，走弱概率73.2\%。
  \item 三个可能触发美元逆转的风险点：7月美国就业数据的季节性残留效应、新任美联储主席Warsh可能释放鸽派信号、AI驱动的科技股叙事面临成本质疑。
  \item GS认为，美元并非被高估，而是被错误定价——能源冲击的“项圈效应”收窄美元波动区间，而非决定方向。
  \item 全球外汇市场的定价权正在从宏观风险偏好（AI叙事、商品价格）向国内利率差异（货币政策分化、实际利差）转移。
  \item 人民币逆势走强，GS预计12个月内美元/人民币目标价6.50，中国贸易顺差回升和出口韧性超预期是支撑。
  \item 英镑政治风险溢价不高，真正决定走势的是能源冲击演变；日元受制于美日利差，日本央行政策空间有限。
  \item 亚洲央行正在打一场“汇率防御战”，政策工具箱从单纯消耗储备转向加息、资本流入管理和行政措施的组合拳。
\end{itemize}
\begin{figure}[htbp]
\centering
\includegraphics[width=0.88\linewidth]{figures/fig_010.jpg}
\caption{Exhibit 1 - Exhibit 1: Global factors like equity and commodity prices could explain the majority of FX moves in March and April}
\end{figure}
\begin{figure}[htbp]
\centering
\includegraphics[width=0.88\linewidth]{figures/fig_011.jpg}
\caption{Exhibit 2 - Exhibit 2: ...but lately more domestic factors—like rate differentials—are starting to become more important once again}
\end{figure}
\begin{figure}[htbp]
\centering
\includegraphics[width=0.88\linewidth]{figures/fig_012.jpg}
\caption{Exhibit 3 - Exhibit 3: Even as the pace of total returns has flattened out over the past month, the past week has marked another year-to-date low in the USD/CNY fix}
\end{figure}
\begin{figure}[htbp]
\centering
\includegraphics[width=0.88\linewidth]{figures/fig_013.jpg}
\caption{Exhibit 4 - Exhibit 4: Modest Sterling reactions to rising odds of a Burnham led government}
\end{figure}
\begin{figure}[htbp]
\centering
\includegraphics[width=0.88\linewidth]{figures/fig_014.jpg}
\caption{Exhibit 5 - Exhibit 5: SNB CHF selling was fairly small in March and April—consistent with other periods where inflation is comfortably within their 0-2\% target range}
\end{figure}
\begin{figure}[htbp]
\centering
\includegraphics[width=0.88\linewidth]{figures/fig_001.jpg}
\caption{Exhibit 1 - Exhibit 1: FX reserves have dropped in most Asian economies since February}
\end{figure}
\begin{figure}[htbp]
\centering
\includegraphics[width=0.88\linewidth]{figures/fig_002.jpg}
\caption{Exhibit 2 - Exhibit 2: CNH continued to strengthen vs. USD despite rangebound DXY}
\end{figure}
\begin{figure}[htbp]
\centering
\includegraphics[width=0.88\linewidth]{figures/fig_003.jpg}
\caption{Exhibit 3 - Exhibit 3: FX conversion ratio for goods trade balance has generally been rising}
\end{figure}
{\small\color{refgray}\textbf{References:} [R001] GS：亚洲央行正在打一场“汇率防御战”，但真正决定胜负的变量不在利率; [R002] GS：美元并非被高估，而是被错误定价——能源冲击的“项圈效应”才是核心; [R009] NOM：美元多头共识已达峰值，历史规律指向未来三个月大概率走弱}

\section{AI与科技：资本开支结构性集中，光互连与PCB成为下一波赢家}
\textbf{全球资本开支增长高度集中于数据中心、半导体和电力公用事业，AI光互连的下一波赢家不在光模块，而在电路板，光收发器PCB市场正经历量价齐升式重估。}

\begin{itemize}
  \item GS Capex Tracker显示，2026年全球资本开支预期增速13\%，但增长几乎完全由数据中心（35\% CAGR）、半导体（23\%）、电力公用事业驱动，三者占比从2022年的22\%跃升至2026年的39\%。
  \item 摩根斯坦利预测，AI光收发器PCB市场规模将从2025年的6.2亿美元增长至2028年的37.7亿美元，CAGR高达83\%，远超光收发器出货量的60\%增速。
  \item 到2028年，AI光收发器PCB市场将超越苹果成为全球最大的HDI PCB需求方，重塑PCB行业竞争格局。
  \item 摩根斯坦利认为，臻鼎（ZDT）是最大赢家，其800G市场份额有望从7.5\%扩至约20\%，1.6T市场份额从5\%扩至约25\%。
  \item 日本线缆企业方面，摩根斯坦利指出超大规模云厂商2026年资本开支同比增速预期已上修至87\%，约75\%投向AI基础设施。
  \item 古河电工被摩根斯坦利列为首选，预计F3/27-29年营业利润CAGR达42.1\%，液冷模块是核心增量。
  \item 住友电工的连续波激光二极管（CW-LD）在磷化铟平台垂直整合，市占率高，CPO产品普及将提升高附加值器件占比。
  \item 藤仓产能是瓶颈，光纤生产受氢气短缺影响，光缆销量增速预计放缓至+10\%。
\end{itemize}
\begin{figure}[htbp]
\centering
\includegraphics[width=0.88\linewidth]{figures/fig_281.jpg}
\caption{Exhibit 3 - Exhibit 3: We see Renewables, Transmission, Datacenters and Semiconductors as structural growth areas GS Capex Tracker by end market}
\end{figure}
\begin{figure}[htbp]
\centering
\includegraphics[width=0.88\linewidth]{figures/fig_284.jpg}
\caption{Exhibit 8 - Exhibit 8: We see a concentration of growth in a small number of areas all related to datacenters, semis and utilities, now representing c.39\% of total capex... End-market size implied by our Capex Tracker, 2026}
\end{figure}
\begin{figure}[htbp]
\centering
\includegraphics[width=0.88\linewidth]{figures/fig_285.jpg}
\caption{Exhibit 9 - Exhibit 9: ....versus only 22\% just 4 years ago End-Market Size implied by our Capex Tracker, 2022}
\end{figure}
\begin{figure}[htbp]
\centering
\includegraphics[width=0.88\linewidth]{figures/fig_286.jpg}
\caption{Exhibit 10 - Exhibit 10: We continue to see capex growth in 2026/27 at historical highs GS General Industrial Capex Tracker (c.4k companies, €3.4 tn of capex in 2026) (adjusted Throughcycle)\textbackslash{}*}
\end{figure}
\begin{figure}[htbp]
\centering
\includegraphics[width=0.88\linewidth]{figures/fig_571.jpg}
\caption{Exhibit 2 - Exhibit 2: Our Global team raised forecasts for 800G and 1.6T optical transceiver shipments Exhibit 3: Our Global team estimates AI optical transceivers growing at \textbackslash{}\textasciitilde{}60\% CAGR from 2025-28E We estimate AI transceivers growing at 60\%}
\end{figure}
\begin{figure}[htbp]
\centering
\includegraphics[width=0.88\linewidth]{figures/fig_572.jpg}
\caption{Exhibit 4 - Exhibit 4: Optical transceiver PCB specs Exhibit 5: We estimate Global AI Transceiver PCB TAM to grow at 83\% CAGR from 2025-28E AI transceiver PCB TAM is growing at 83\% CAR from 2025-28E (US\$ M)}
\end{figure}
\begin{figure}[htbp]
\centering
\includegraphics[width=0.88\linewidth]{figures/fig_573.jpg}
\caption{Exhibit 6 - Exhibit 6: TAM comparison: AI transceiver vs. AI transceiver PCB}
\end{figure}
\begin{figure}[htbp]
\centering
\includegraphics[width=0.88\linewidth]{figures/fig_574.jpg}
\caption{Exhibit 7 - Exhibit 7: AI transceiver PCB value contribution to suppliers - we expect suppliers to see 50-176\% CAGR over CY25-28E AI Transceiver PCB value by supplier (US\textbackslash{}\$M)}
\end{figure}
\begin{figure}[htbp]
\centering
\includegraphics[width=0.88\linewidth]{figures/fig_581.jpg}
\caption{Exhibit 5 - Exhibit 5: Hyperscaler CAPEX YoY}
\end{figure}
\begin{figure}[htbp]
\centering
\includegraphics[width=0.88\linewidth]{figures/fig_582.jpg}
\caption{Exhibit 6 - Exhibit 6: Hyperscaler capital intensity Top 14 Cloud Providers: Capital Intensity}
\end{figure}
\begin{figure}[htbp]
\centering
\includegraphics[width=0.88\linewidth]{figures/fig_583.jpg}
\caption{Exhibit 7 - Exhibit 7: Hyperscaler CAPEX}
\end{figure}
{\small\color{refgray}\textbf{References:} [R023] GS：资本开支的“结构性集中”正在重塑全球工业投资逻辑; [R032] 摩根斯坦利：AI光互连的下一波赢家不在光模块，而在电路板; [R033] 摩根斯坦利：市场低估的不是AI投资放缓，而是日本线缆企业供给侧的差异化定价权}

\section{能源与大宗：欧洲石油巨头隐含油价远低于共识，厄尔尼诺精准打击作物窗口}
\textbf{欧洲石油巨头估值隐含的长期油价仅60美元/桶，远低于市场共识，提供不对称风险收益窗口；厄尔尼诺已确认，极强事件概率跃升至63\%，糖是最敏感的品种。}

\begin{itemize}
  \item GS求解欧洲石油巨头在实现8\%自由现金流收益率时对应的2026-2029年平均布伦特油价仅为60美元/桶，而彭博共识预期为78美元，远期期货为79美元。
  \item 在70美元的中性情景下，欧洲石油巨头的股东总回报率（分红+回购）仍接近7\%，不算差。
  \item GS看好的三家欧洲石油公司：BP（财务去杠杆空间大）、Repsol（受益于欧盟航空燃料市场偏紧）、Galp（油气储量寿命最长）。
  \item NOAA确认厄尔尼诺已正式形成，极强事件概率从37\%跃升至63\%，摩根斯坦利认为市场已定价“厄尔尼诺会发生”，但远未定价其精准打击作物窗口的时间和强度。
  \item 糖是最敏感的品种，厄尔尼诺带来的干旱可能收紧亚洲供应，印度作为全球第二大产糖国，季风雨量面临下行风险。
  \item 大豆和玉米的影响取决于时间窗口：巴西中西部和MATOPIBA产区面临干旱风险，但南里奥格兰德州和阿根廷反而可能受益。
  \item 巴拿马运河可能再次面临水位限制，推高运费成本；秘鲁鳀鱼捕捞因海水变暖受阻，鱼粉供应紧张会传导至养殖成本。
\end{itemize}
\begin{figure}[htbp]
\centering
\includegraphics[width=0.88\linewidth]{figures/fig_177.jpg}
\caption{Exhibit 3 - Exhibit 3: On our estimates, EU Oils are pricing a long-term Brent of \textbackslash{}\$60/bbl solving for an 8\% FCF yield Implied 2026-29 oil price for the EU majors solving by a 8\% FCF yield (Brent, \textbackslash{}\$/bl)}
\end{figure}
\begin{figure}[htbp]
\centering
\includegraphics[width=0.88\linewidth]{figures/fig_178.jpg}
\caption{Exhibit 4 - Exhibit 4: On consensus estimates, EU Oils are pricing a long-term Brent of \textbackslash{}\$60/bbl solving for an 8\% FCF yield Implied 2026-29 oil price for the EU majors by a 8\% FCF yield (Brent, \textbackslash{}\$/bl)}
\end{figure}
\begin{figure}[htbp]
\centering
\includegraphics[width=0.88\linewidth]{figures/fig_179.jpg}
\caption{Exhibit 5 - Exhibit 5: The market is currently pricing 2030E Brent at \textbackslash{}\$70/\textbackslash{}\$72/bl on Bloomberg consensus/forwards (vs GSe of \textbackslash{}\$75/bl).}
\end{figure}
\begin{figure}[htbp]
\centering
\includegraphics[width=0.88\linewidth]{figures/fig_544.jpg}
\caption{Exhibit 1 - Exhibit 1: Central/eastern Pacific SST anomalies are now clearly positive, confirming the ocean-side setup for a likely very strong El Niño Relative SST Anomalies (°C) 03 JUN 2026}
\end{figure}
\begin{figure}[htbp]
\centering
\includegraphics[width=0.88\linewidth]{figures/fig_545.jpg}
\caption{Exhibit 2 - Exhibit 2: NOAA/CPC now assigns a 63\% (up from 37\%) probability to a very strong El Niño in Nov-Jan, making intensity the key watchpoint NOAA CPC ENSO Strength Probabilities (issued June 2026) based on thresholds in ERSSTv5 Relat}
\end{figure}
\begin{figure}[htbp]
\centering
\includegraphics[width=0.88\linewidth]{figures/fig_546.jpg}
\caption{Exhibit 4 - Exhibit 4: El Niño is now confirmed, while very-strong odds increased to 63\% in Nov-Jan from 37\% in the prior update Exhibit 5: El Niño strength distribution remains wide}
\end{figure}
\begin{figure}[htbp]
\centering
\includegraphics[width=0.88\linewidth]{figures/fig_551.jpg}
\caption{Exhibit 12 - Exhibit 12: India usually receives below-normal rainfall during El Niño years. India is the world second largest sugar producer and below normal monsoon may reduce its sugar output potential Correlation between El Niño and Indian}
\end{figure}
\begin{figure}[htbp]
\centering
\includegraphics[width=0.88\linewidth]{figures/fig_552.jpg}
\caption{Exhibit 14 - Exhibit 14: Strong El Niño historically leads to soybean yield losses in MT, which is the largest producing state with 29\% of national output}
\end{figure}
\begin{figure}[htbp]
\centering
\includegraphics[width=0.88\linewidth]{figures/fig_553.jpg}
\caption{Exhibit 13 - Exhibit 13: Sugar remains the cleanest bullish El Niño trade}
\end{figure}
\begin{figure}[htbp]
\centering
\includegraphics[width=0.88\linewidth]{figures/fig_560.jpg}
\caption{Exhibit 23 - Exhibit 23: Sugar remains the cleanest bullish El Niño trade.}
\end{figure}
{\small\color{refgray}\textbf{References:} [R015] GS：欧洲石油巨头隐含的长期油价远低于市场共识，这才是真正的价值锚点; [R030] 摩根斯坦利：市场真正低估的不是厄尔尼诺本身，而是它对作物窗口的精准打击}

\section{地产与消费：中国楼市K型分化，香港写字楼迎结构性拐点}
\textbf{中国房地产并非整体回暖，而是K型分化下的资产再定价——高端住宅需求旺盛，中间价位段持续承压；香港写字楼市场正经历供给端驱动的结构性拐点。}

\begin{itemize}
  \item JPM实地调研显示，深圳和上海地产市场正经历K型复苏：高端住宅（2000万+）靠科技新贵撑场，入门级住房（约300万）回暖，中间价位段（500-1000万）最弱。
  \item 摩根斯坦利高频数据显示，25城二手房实时成交同比增速已从4月的30\%骤降至6月初的9\%，低能级城市已率先转弱。
  \item 摩根斯坦利判断，若政策效果和积压需求继续衰减，3季度二手房销售可能转负，2026-2027年房价大概率继续承压。
  \item 香港住宅库存从2025年Q1的28,000套峰值降至20,000套，开发商定价权正在从买方市场回归。
  \item 摩根斯坦利预测香港2026年全年住宅价格上涨7\%-10\%，年初至今已涨约6\%。
  \item 香港甲级写字楼市场正经历供给端结构性拐点，未来4年新增供应仅260万平方英尺，中环和西九龙是复苏主力。
  \item 香港学生公寓正从边缘话题变成可规模化的新资产类别，八大学生与床位比3.4:1，缺口值得关注。
  \item 中国酒店业内部定价逻辑分化：中端市场用降价换入住率，高端市场靠供给侧约束和入境需求维持价格刚性。
\end{itemize}
\begin{figure}[htbp]
\centering
\includegraphics[width=0.88\linewidth]{figures/fig_454.jpg}
\caption{Figure 1 - Figure 1: Strong sell-through at Marivista (观潮) in Bao'an, Shenzhen 深圳观潮房源销控表 MARIVISTA 深圳观潮房源销控表 MARIVISTA}
\end{figure}
\begin{figure}[htbp]
\centering
\includegraphics[width=0.88\linewidth]{figures/fig_455.jpg}
\caption{Figure 2 - Figure 2: Neighborhood and school network near Marivista (观潮) in Bao'an, Shenzhen 西安云游集团海洋学校 西安中宇集团海洋学校}
\end{figure}
\begin{figure}[htbp]
\centering
\includegraphics[width=0.88\linewidth]{figures/fig_458.jpg}
\caption{Figure 5 - Figure 5: COLI's Anlan Shanghai (left) and HKL's Westbund Central project (right) in Xuhui, Shanghai 3D architectural model of a modern city with skyscrapers and road networks, no visible text or symbols}
\end{figure}
\begin{figure}[htbp]
\centering
\includegraphics[width=0.88\linewidth]{figures/fig_459.jpg}
\caption{Figure 6 - Figure 6: Shui On's Riverville villa project in Yangpu, Shanghai Exterior view of a modern brick building with a decorative roof and glass balcony, surrounded by greenery and urban buildings in the background (no visible text or}
\end{figure}
\begin{figure}[htbp]
\centering
\includegraphics[width=0.88\linewidth]{figures/fig_463.jpg}
\caption{Exhibit 1 - Exhibit 1: Bingshan 25-city secondary real-time home sales continued to decelerate, to \$+9\textbackslash{}\%\$ y-y in MTD June}
\end{figure}
\begin{figure}[htbp]
\centering
\includegraphics[width=0.88\linewidth]{figures/fig_464.jpg}
\caption{Exhibit 2 - Exhibit 2: Chinese New Year-adjusted cumulative secondary real-time home sales fell \$\textbackslash{}sim 1\textbackslash{}\%\$ y-y in those cities}
\end{figure}
\begin{figure}[htbp]
\centering
\includegraphics[width=0.88\linewidth]{figures/fig_466.jpg}
\caption{Exhibit 1 - Exhibit 1: Official R\&V Home Price Rose 5.7\% YTD}
\end{figure}
\begin{figure}[htbp]
\centering
\includegraphics[width=0.88\linewidth]{figures/fig_467.jpg}
\caption{Exhibit 2 - Exhibit 2: Office - 1Q26 New Lettings}
\end{figure}
\begin{figure}[htbp]
\centering
\includegraphics[width=0.88\linewidth]{figures/fig_565.jpg}
\caption{Figure 1 - Figure 1: Mainland China RevPAR YoY change}
\end{figure}
\begin{figure}[htbp]
\centering
\includegraphics[width=0.88\linewidth]{figures/fig_567.jpg}
\caption{Figure 3 - Figure 3: RevPAR YoY performance by segment (May 31-Jun 6) RevPAR (YoY change \%)}
\end{figure}
\begin{figure}[htbp]
\centering
\includegraphics[width=0.88\linewidth]{figures/fig_568.jpg}
\caption{Figure 4 - Figure 4: ADR YoY performance by segment (May 31-Jun 6) ADR (YoY change \%)}
\end{figure}
\begin{figure}[htbp]
\centering
\includegraphics[width=0.88\linewidth]{figures/fig_569.jpg}
\caption{Figure 5 - Figure 5: OCC YoY performance by segment (May 31-Jun 6) Occ (YoY change ppt)}
\end{figure}
{\small\color{refgray}\textbf{References:} [R004] GS：高频数据正在揭示一个被低估的“成本传导型”经济周期; [R024] JPM：中国房地产并非整体回暖，而是一场K型分化下的资产再定价; [R025] 摩根斯坦利：市场对楼市“小阳春”的可持续性存在显著误判; [R027] 摩根斯坦利：香港地产市场真正被低估的不是住宅，而是写字楼的结构性拐点; [R031] UBS：酒店业真正的拐点不在RevPAR，而在定价策略的深层分化}

\section{中国信贷与金融：总量超预期但结构疲弱，流动性正常化传导加速}
\textbf{五月信贷超预期不是复苏信号，而是结构性疲弱的又一次确认——总量靠银行供给端发力，实体融资需求并未真实回暖。}

\begin{itemize}
  \item GS指出，5月新增人民币贷款5200亿，高于预期，但企业中长期贷款净减少200亿（去年同期净增3300亿），票据融资贡献5570亿。
  \item 票据融资飙升而中长期贷款收缩，说明企业没有投资意愿，银行只能用低成本票据融资来“填满”信贷额度。
  \item 居民端5月住户贷款存量减少1410亿，而去年同期是增加540亿，居民仍在去杠杆。
  \item M1同比从5.0\%微升至5.5\%，但财政存款5月增加了7100亿，M1回升可能靠财政“输血”。
  \item NOM指出，央行已在悄悄收紧中期流动性，3月以来通过MLF、买断式逆回购等工具持续净回笼，累计净投放降至仅140亿元。
  \item 政府债供给即将放量，前5个月发行进度仅35\%，Q3供给压力会显著加大。
  \item 7天回购利率已回到1.4\%以上，DR007和1年NCD收益率自6月5日起明显上行。
  \item 机构行为方面，5月基金大买长债后，6月开始减仓，保险则在30年国债下跌时开始加仓。
\end{itemize}
\begin{figure}[htbp]
\centering
\includegraphics[width=0.88\linewidth]{figures/fig_140.jpg}
\caption{Exhibit 2 - Exhibit 1: TSF flows rose from April to May after seasonal adjustment}
\end{figure}
\begin{figure}[htbp]
\centering
\includegraphics[width=0.88\linewidth]{figures/fig_141.jpg}
\caption{Exhibit 2 - Exhibit 2: M1 growth edged up in May}
\end{figure}
\begin{figure}[htbp]
\centering
\includegraphics[width=0.88\linewidth]{figures/fig_077.jpg}
\caption{Figure 5 - \#\# 4. Our view on liquidity in coming weeks: We continue to expect the PBoC to be more flexible with 7d OMOs to inject short-term liquidity amid higher government bond supply, as well as tax payment and month-end funding needs. As for medium- to long-term liqu}
\end{figure}
\begin{figure}[htbp]
\centering
\includegraphics[width=0.88\linewidth]{figures/fig_078.jpg}
\caption{Figure 10 - As shown in Figure 10, after accumulating massive (ultra) long-end CGBs/policy bank bonds in April and May, onshore funds reduced bond holdings in June. In the first week of June, funds were still adding 10y and above tenors while net selling the front-end and}
\end{figure}
{\small\color{refgray}\textbf{References:} [R006] NOM：中国利率市场真正被低估的风险不是降息预期，而是流动性正常化的传导加速; [R008] GS：五月信贷超预期，但真正的信号藏在结构里}

\section{电动车与供应链：中国渗透率结构性跃迁，电力设备回调是重新定价}
\textbf{中国电动车渗透率完成结构性跃迁，即使在宏观逆风中也能自主增长；亚洲电力设备板块20-30\%回调不是风险出清，而是结构性机会的重新定价。}

\begin{itemize}
  \item JPM数据显示，5月中国BEV渗透率42\%，总EV渗透率63\%，创2025年8月以来新高，在乘用车市场同比降22\%背景下BEV销量逆势增长5\%。
  \item 全球三大市场中，只有中国在完成真正的渗透率跃迁，欧盟渗透率24\%，美国连续五个月停在6\%。
  \item JPM判断锂辉石价格有温和上行压力，储能市场持续强劲+供应端滞后提供支撑。
  \item JPM认为亚洲电力设备股从4月高点回撤20-30\%不是基本面恶化，而是估值、资金流和催化剂真空三重叠加。
  \item 回调后的估值水平为订单结构偏向公用事业和电网的公司提供了更有吸引力的入场窗口。
  \item 中国电力设备在美渗透拐点比预期更快到来，变压器交付周期6-12个月 vs 西方同行2-3年，美国CSP已开始参观中国工厂并下订单。
  \item 欧盟对中国电力设备加征关税的可能性较低，中国产品在欧盟市占率不到10\%，对本地厂商构不成威胁。
\end{itemize}
\begin{figure}[htbp]
\centering
\includegraphics[width=0.88\linewidth]{figures/fig_155.jpg}
\caption{Figure 1 - Figure 1: China vehicle sales vs. BEV sales}
\end{figure}
\begin{figure}[htbp]
\centering
\includegraphics[width=0.88\linewidth]{figures/fig_156.jpg}
\caption{Figure 2 - Figure 2: China EV penetration rates}
\end{figure}
\begin{figure}[htbp]
\centering
\includegraphics[width=0.88\linewidth]{figures/fig_157.jpg}
\caption{Figure 3 - Figure 3: Battery electric vehicle penetration rates by region}
\end{figure}
\begin{figure}[htbp]
\centering
\includegraphics[width=0.88\linewidth]{figures/fig_158.jpg}
\caption{Figure 4 - Figure 4: Battery electric vehicle sales – monthly (China + EU + US)}
\end{figure}
\begin{figure}[htbp]
\centering
\includegraphics[width=0.88\linewidth]{figures/fig_167.jpg}
\caption{Figure 14 - Figure 14: Lithium spodumene prices}
\end{figure}
\begin{figure}[htbp]
\centering
\includegraphics[width=0.88\linewidth]{figures/fig_168.jpg}
\caption{Figure 15 - Figure 15: Lithium hydroxide prices}
\end{figure}
\begin{figure}[htbp]
\centering
\includegraphics[width=0.88\linewidth]{figures/fig_169.jpg}
\caption{Figure 16 - Figure 16: Lithium carbonate prices}
\end{figure}
\begin{figure}[htbp]
\centering
\includegraphics[width=0.88\linewidth]{figures/fig_541.jpg}
\caption{Figure 1 - Figure 1: Korea Power Equipment Companies' Valuation}
\end{figure}
\begin{figure}[htbp]
\centering
\includegraphics[width=0.88\linewidth]{figures/fig_542.jpg}
\caption{Figure 2 - Figure 2: Fund Rotations from China Power Equipment into China AI-Themed Stocks}
\end{figure}
\begin{figure}[htbp]
\centering
\includegraphics[width=0.88\linewidth]{figures/fig_543.jpg}
\caption{Figure 3 - Figure 3: Share Price Weaknesses Across Global Power Equipment Companies Since May}
\end{figure}
\begin{figure}[htbp]
\centering
\includegraphics[width=0.88\linewidth]{figures/fig_252.jpg}
\caption{Figure 1 - Figure 1: Eaton's MV SST 2.0 launched in IDCE 2026 Exterior view of a large industrial control room with multiple panels and control panels (no visible text or symbols)}
\end{figure}
\begin{figure}[htbp]
\centering
\includegraphics[width=0.88\linewidth]{figures/fig_253.jpg}
\caption{Figure 2 - Figure 2: TGOOD's AI PowerHouse Exterior view of a modern industrial building with green and white facade signage, featuring a large 'Xingtong' logo and Chinese text on the facade (no other signage visible)}
\end{figure}
{\small\color{refgray}\textbf{References:} [R012] JPM：全球电动车市场的真正拐点不在销量增速，而在中国渗透率的结构性跃迁; [R020] JPM：中国电力设备在美渗透的拐点比市场预期的更快到来; [R029] JPM：电力设备板块的20-30\%回调，不是风险出清，而是结构性机会的重新定价}

\section{区域市场：韩国居民资产结构性迁移，日本治理改善信号被低估}
\textbf{韩国居民资产从房地产向权益市场的迁移具有结构性而非周期性，日本市场真正值得关注的不是AI反弹，而是股东会投票率背后的治理信号。}

\begin{itemize}
  \item 摩根斯坦利指出，2025年韩国居民权益资产持有量同比增长48\%，此前十年平均增速仅7.4\%，房地产政策收紧与股市改革是核心驱动力。
  \item 韩国散户投资者数量从2019年的620万增至2025年的1450万，占成年人50\%，40岁以下投资者占43\%。
  \item 摩根斯坦利上调KOSPI 12个月目标至9000点（此前8500点），牛市情景上调至10500点。
  \item GS发现，日本去年股东大会低赞成票率的公司正在系统性地改善治理表现，12家中有10家今年赞成率明显提升。
  \item 低赞成票率组合在去年股东大会到今年伊朗冲突爆发前跑赢TOPIX指数13个百分点，是预测企业未来股东友好行为的最强先行指标。
  \item 日本“股东优待+分红”组合在市场回调时抗跌，6月3日以来相对既无优待也无分红的股票跑赢12个百分点。
  \item 韩国KOSDAQ正吸引结构性外资流入，而KOSPI遭遇大规模外资流出，GS认为KOSDAQ存在战术性alpha机会。
  \item KOSDAQ估值虽比KOSPI贵（28.8倍 vs 7.1倍远期PE），但盈利修正趋势在改善，外资流入集中在信息技术和医疗保健板块。
\end{itemize}
\begin{figure}[htbp]
\centering
\includegraphics[width=0.88\linewidth]{figures/fig_468.jpg}
\caption{Exhibit 1 - Exhibit 1: Historically, Korean households have shown a heavy bias toward real assets over financial assets – marginal asset growth led mainly by real assets Korean households annual change in asset composition}
\end{figure}
\begin{figure}[htbp]
\centering
\includegraphics[width=0.88\linewidth]{figures/fig_469.jpg}
\caption{Exhibit 2 - Exhibit 2: Real estate assets across main and other residences composed 70\% of total household assets while the size of total assets nearly doubled in the past decade Household asset breakdown by sub-category (2025)}
\end{figure}
\begin{figure}[htbp]
\centering
\includegraphics[width=0.88\linewidth]{figures/fig_470.jpg}
\caption{Exhibit 3 - Exhibit 3: Korean household assets nearly doubled during 2000-24; fastest among DMs}
\end{figure}
\begin{figure}[htbp]
\centering
\includegraphics[width=0.88\linewidth]{figures/fig_471.jpg}
\caption{Exhibit 5 - Exhibit 5: The number of retail investors increased to 50\% of total adult population vs. 21\% in 2019 Number of retail investors and \% to adult population}
\end{figure}
\begin{figure}[htbp]
\centering
\includegraphics[width=0.88\linewidth]{figures/fig_472.jpg}
\caption{Exhibit 4 - Exhibit 4: Korea registered a rapid expansion in financial assets of +94\% vs. 31\% in Japan during 2015-25}
\end{figure}
\begin{figure}[htbp]
\centering
\includegraphics[width=0.88\linewidth]{figures/fig_473.jpg}
\caption{Exhibit 6 - Exhibit 6: 43\% of retail investors are aged 40s or below, groups with a higher marginal propensity to consume \% share of retail investors by age group (2025)}
\end{figure}
\begin{figure}[htbp]
\centering
\includegraphics[width=0.88\linewidth]{figures/fig_478.jpg}
\caption{Exhibit 12 - Exhibit 12: We raise our KOSPI target to 9,000 (from 8,500) with our new bull case being 10,500 and our bear case set at 6,500. We set our three- to six-month KOSPI trading band at 7,000-10,000.}
\end{figure}
\begin{figure}[htbp]
\centering
\includegraphics[width=0.88\linewidth]{figures/fig_043.jpg}
\caption{Exhibit 1 - Exhibit 1: KOSDAQ has significantly underperformed since August 2023}
\end{figure}
\begin{figure}[htbp]
\centering
\includegraphics[width=0.88\linewidth]{figures/fig_044.jpg}
\caption{Exhibit 2 - Exhibit 2: KOSPI earnings upgrades have been materially stronger than KOSDAQ earnings revisions}
\end{figure}
\begin{figure}[htbp]
\centering
\includegraphics[width=0.88\linewidth]{figures/fig_045.jpg}
\caption{Exhibit 3 - Exhibit 3: KOSDAQ valuations remain more elevated than KOSPI's 12-month forward P/E}
\end{figure}
\begin{figure}[htbp]
\centering
\includegraphics[width=0.88\linewidth]{figures/fig_046.jpg}
\caption{Exhibit 4 - Exhibit 4: However, KOSDAQ has seen stronger foreign inflows year-to-date, while KOSPI has experienced significant foreign outflows}
\end{figure}
\begin{figure}[htbp]
\centering
\includegraphics[width=0.88\linewidth]{figures/fig_047.jpg}
\caption{Exhibit 5 - Exhibit 5: Foreign outflows from KOSPI have been largely driven by the two largest semiconductor stocks}
\end{figure}
\begin{figure}[htbp]
\centering
\includegraphics[width=0.88\linewidth]{figures/fig_048.jpg}
\caption{Exhibit 6 - Exhibit 6: Foreign inflows into KOSDAQ have been concentrated in the information technology and health care sectors}
\end{figure}
\begin{figure}[htbp]
\centering
\includegraphics[width=0.88\linewidth]{figures/fig_180.jpg}
\caption{Exhibit 1 - Exhibit 1: Mar FY-end companies AGM schedule by number of companies}
\end{figure}
\begin{figure}[htbp]
\centering
\includegraphics[width=0.88\linewidth]{figures/fig_183.jpg}
\caption{Exhibit 4 - Exhibit 4: Since the 2025 AGM season, low approval screen outperformed until the start of the Iran war, but remains a laggard as the market focused on AI themes}
\end{figure}
\begin{figure}[htbp]
\centering
\includegraphics[width=0.88\linewidth]{figures/fig_184.jpg}
\caption{Exhibit 6 - Exhibit 6: Y+D screen has outperformed NYND screen by 12ppt since June 3rd ...}
\end{figure}
\begin{figure}[htbp]
\centering
\includegraphics[width=0.88\linewidth]{figures/fig_185.jpg}
\caption{Exhibit 7 - Exhibit 7: ... and SMID Y+D has outperformed SMID NYND by 7ppts in the same period}
\end{figure}
{\small\color{refgray}\textbf{References:} [R005] GS：韩国股市的真正机会不在大盘，而在被低估的KOSDAQ; [R016] GS：日本市场真正值得关注的不是AI反弹，而是股东会投票率背后的治理信号; [R028] 摩根斯坦利：韩国居民资产迁移到股票市场的趋势具有结构性，而非周期性}

\section{医疗与教育：AI放射科从可选项变必需品，教育市场结构分化加速}
\textbf{AI在放射科的渗透正从“可选项”转变为“必需品”，老龄化叠加医生供给缺口创造结构性需求；中国教育市场最值得跟踪的不是增长，而是AI原生产品、硬件定价权和海外需求三条逻辑线的分化加速。}

\begin{itemize}
  \item Bernstein指出，美国65岁以上人口将以每年约3\%速度增长至2030年，而放射科医生供给缺口将在2034年达到10\%，AI是目前唯一能填补缺口的方案。
  \item 英国NHS数据显示，等待超过六周进行CT或MRI检查的患者从疫情前约8000人飙升至2020年约80000人，2025年仍高达约75000人。
  \item McKinsey数据显示，AI能够释放医疗设备准备人员高达48\%的工作时间。
  \item GS数据显示，字节跳动的“豆包爱学”DAU同比增速在5月底加速至146\%，稳居中国第6大AI原生应用。
  \item Gauth（好未来旗下出海AI产品）5月月流水约110万美元，同比增长24\%，但增长斜率正在放缓。
  \item AI教育应用竞争已从“谁先推出产品”进入“谁能持续获取用户注意力”阶段。
  \item 中国学习机市场量跌价稳，6大品牌线上GMV同比-35\%，但均价同比降幅收窄至-4\%，价格战趋缓。
  \item 留学签证回暖，中国学生赴加拿大、澳洲、英国签证量1季度同比-9\%，相比去年-22\%的降幅明显收窄。
\end{itemize}
\begin{figure}[htbp]
\centering
\includegraphics[width=0.88\linewidth]{figures/fig_240.jpg}
\caption{Exhibit 2 - EXHIBIT 2: The U.S. senior population will grow at a 3\% CAGR through 2030}
\end{figure}
\begin{figure}[htbp]
\centering
\includegraphics[width=0.88\linewidth]{figures/fig_241.jpg}
\caption{EXHIBIT 3 - EXHIBIT 3: The European senior population will grow at a 2\% CAGR through 2030}
\end{figure}
\begin{figure}[htbp]
\centering
\includegraphics[width=0.88\linewidth]{figures/fig_242.jpg}
\caption{Exhibit 4 - EXHIBIT 4: The gap between U.S. demand and supply of physicians in radiology is expected to reach -10\% by 2034}
\end{figure}
\begin{figure}[htbp]
\centering
\includegraphics[width=0.88\linewidth]{figures/fig_243.jpg}
\caption{Exhibit 5 - EXHIBIT 5: AI could free up a significant amount of time for hospital staff}
\end{figure}
\begin{figure}[htbp]
\centering
\includegraphics[width=0.88\linewidth]{figures/fig_244.jpg}
\caption{Exhibit 6 - EXHIBIT 6: The number of patients who had to wait at least six weeks for a CT or an MRI more than quadrupled between 2019 and 2020, and continues to be elevated compared to pre-pandemic Patients waiting six weeks or more for a MRI}
\end{figure}
\begin{figure}[htbp]
\centering
\includegraphics[width=0.88\linewidth]{figures/fig_254.jpg}
\caption{Exhibit 1 - Exhibit 1: Doubao Study DAU growth accelerated to +146\% yoy as of May-end and sustained the position of the 6th largest AI native app in China in terms of DAUs DAUs of major AI apps as of Apr-end}
\end{figure}
\begin{figure}[htbp]
\centering
\includegraphics[width=0.88\linewidth]{figures/fig_255.jpg}
\caption{Exhibit 3 - Exhibit 3: Gauth has outperformed global AI-native education peers in terms of MAUs MAUs by AI education apps (mn)}
\end{figure}
\begin{figure}[htbp]
\centering
\includegraphics[width=0.88\linewidth]{figures/fig_256.jpg}
\caption{Exhibit 2 - Exhibit 2: Doubao Study continued to expand time spent market share among major K-12 learning tool/study companion apps in May Time spent share among major K-12 learning tool/study companion apps Exhibit 4: Gauth achieved c.US\textbackslash{}\$1}
\end{figure}
\begin{figure}[htbp]
\centering
\includegraphics[width=0.88\linewidth]{figures/fig_258.jpg}
\caption{Exhibit 5 - Exhibit 5: Gauth monthly billings per avg. DAU grew +39\% yoy to US\textbackslash{}\$0.4 in May (vs. +49\% yoy in Apr) Gauth monthly billings (US\textbackslash{}\$ mn) and monthly billings per avg. DAU (US\textbackslash{}\$)}
\end{figure}
\begin{figure}[htbp]
\centering
\includegraphics[width=0.88\linewidth]{figures/fig_263.jpg}
\caption{Exhibit 10 - GMV: Combined GMV of 6 major learning device brands dropped -35\% yoy or increased +46\% mom in May (vs. -2\% yoy/-47\% mom in Apr) on Douyin/Taobao/Tmall/JD. By brand, GMV of (+) Yuanfudao increased +31\% yoy in May, while (-) iFlytek/TAL/Xiaodu/BBK/Zuoyebang drop}
\end{figure}
\begin{figure}[htbp]
\centering
\includegraphics[width=0.88\linewidth]{figures/fig_269.jpg}
\caption{Exhibit 20 - Exhibit 20: Student visas \# granted to Chinese for Canada, Australia, and UK combined dropped -9\% yoy in 1Q26}
\end{figure}
\begin{figure}[htbp]
\centering
\includegraphics[width=0.88\linewidth]{figures/fig_275.jpg}
\caption{Exhibit 26 - Exhibit 26: East Buy monthly GMV on Douyin grew +38\% yoy/+7\% mom to c.Rmb1.03bn in May (vs. +42\% yoy/-10\% mom in Apr) East Buy monthly GMV on Douyin (Rmb mn)}
\end{figure}
{\small\color{refgray}\textbf{References:} [R018] Bernstein：AI在放射科的落地，真正被低估的不是技术，而是结构性供需缺口; [R022] GS：中国教育市场最值得跟踪的不是增长，而是结构分化的加速}

\section{风险偏好与资金流向：全球资金重新定价“安全”，而非追逐风险}
\textbf{全球资金流向传递的信号比表面数据更复杂——资金同时涌入美国科技股和短期债券，却在撤出中国A股和欧洲股票，市场在重新定义什么是当前环境下真正的“安全资产”。}

\begin{itemize}
  \item GS数据显示，截至6月10日当周，全球股票基金净流入310亿美元，固收基金净流入177亿美元，但驱动力量几乎完全来自美国市场。
  \item 欧洲股票基金持续净流出，新兴市场内部出现显著分化：台湾和韩国股票基金是EM净流入主要来源，中国大陆股票基金则录得净流出。
  \item 行业层面，科技基金是最大赢家，消费品基金继续失血，工业、基建类基金年初至今累计流入最多。
  \item 债券基金全面获支撑，短期债和通胀保护债尤其受欢迎，新兴市场硬通货债券基金出现净流出。
  \item 匈牙利债券流入显著增加，GS认为匈牙利经济政策转向将推动其收益率向欧元区收敛。
  \item 跨境外汇方面，美元和韩元需求最强，印度卢比、巴西雷亚尔和人民币流出最多。
  \item 从区域累计数据看，韩国、巴西、台湾是今年资金流入最多的市场，中国内地和印度流出明显。
\end{itemize}
\begin{figure}[htbp]
\centering
\includegraphics[width=0.88\linewidth]{figures/fig_049.jpg}
\caption{Exhibit 8 - Up (↑) = Up wow vs. the previous week Asterisk (\textbackslash{}*) = Expressed in standard deviation of 1-wk change in 1-year Year-to-date Foreign Inflows to Korea}
\end{figure}
\begin{figure}[htbp]
\centering
\includegraphics[width=0.88\linewidth]{figures/fig_050.jpg}
\caption{Exhibit 9 - Exhibit 9: Equity inflows to 5 AEJ markets 4-week rolling sum}
\end{figure}
\begin{figure}[htbp]
\centering
\includegraphics[width=0.88\linewidth]{figures/fig_051.jpg}
\caption{Exhibit 10 - Exhibit 10: Equity inflows to 5 AEJ markets}
\end{figure}
\begin{figure}[htbp]
\centering
\includegraphics[width=0.88\linewidth]{figures/fig_052.jpg}
\caption{Exhibit 11 - Exhibit 11: Bond inflows to 4 AEJ markets 4-week rolling sum}
\end{figure}
\begin{figure}[htbp]
\centering
\includegraphics[width=0.88\linewidth]{figures/fig_053.jpg}
\caption{Exhibit 12 - Exhibit 12: Bond inflows to 4 AEJ markets}
\end{figure}
{\small\color{refgray}\textbf{References:} [R013] GS：全球资金正在重新定价“安全”，而非追逐风险}

\section{汽车半导体与人形机器人：复苏已至，但结构性分化是关键}
\textbf{汽车半导体复苏已确认，但真正的分歧在中国放缓与工业超预期之间；人形机器人供应链的“确定性”正从概念走向订单，但市场仍低估了传统业务的反转弹性。}

\begin{itemize}
  \item UBS确认，2026年Q1汽车半导体收入同比增速达7.8\%，模拟芯片收入同比增长20\%，预计Q2也将有21\%增长。
  \item UBS维持2026年汽车半导体市场10\%同比增长预测，但提醒中国1-5月零售数据已出现15\%-20\%同比下滑，若引发库存修正，全球增长面临下行风险。
  \item 工业半导体需求被AI相关产品拉动，UBS将2026年工业半导体收入增速预测从16\%上调至25\%。
  \item 模拟芯片板块12个月前瞻PE已达30倍，十年均值仅19倍，安全边际在收窄。
  \item UBS实地调研显示，双环传动Q2 NEV齿轮出货环比显著增长，产能利用率已超90\%，验证国内NEV需求企稳回升。
  \item 人形机器人方面，双环传动与北美头部人形机器人公司联合研发新型减速器，荣泰电工的灵巧手精密螺丝已获生产批准，6月订单爬坡。
  \item 传统汽车供应链公司凭借精密制造能力，正在打开第二增长曲线，人形机器人提供估值弹性，传统业务复苏提供安全边际。
\end{itemize}
\begin{figure}[htbp]
\centering
\includegraphics[width=0.88\linewidth]{figures/fig_079.jpg}
\caption{Figure 1 - Figure 1: UBS leading indicators for Autos/Industrial semis - we see a general positive trend in terms of cycle Figure 2: We expect next quarter to grow by 9\% q-o-q after an increase of 4\% in Q1'26...}
\end{figure}
\begin{figure}[htbp]
\centering
\includegraphics[width=0.88\linewidth]{figures/fig_080.jpg}
\caption{Figure 3 - Figure 3: ...Q2'26E likely to deliver 21\% y-o-y growth after 20\% in Q1'26}
\end{figure}
\begin{figure}[htbp]
\centering
\includegraphics[width=0.88\linewidth]{figures/fig_081.jpg}
\caption{Figure 4 - Figure 4: CQ2 26E Guidance Comp. Y-o-Y \% Analog Semis}
\end{figure}
\begin{figure}[htbp]
\centering
\includegraphics[width=0.88\linewidth]{figures/fig_082.jpg}
\caption{Figure 5 - Figure 5: CQ2 26E Guidance Comp. Q-o-Q \% Analog Semis}
\end{figure}
\begin{figure}[htbp]
\centering
\includegraphics[width=0.88\linewidth]{figures/fig_083.jpg}
\caption{Figure 6 - Figure 6: We estimate 10\% y-o-y auto semis revenue growth in 2026E (unchanged) and +c14\% in 2027E vs (unchanged) based on our bottom-up estimates. Figure 7: The analog semis within our coverage are trading on a -17\% discount to se}
\end{figure}
\begin{figure}[htbp]
\centering
\includegraphics[width=0.88\linewidth]{figures/fig_084.jpg}
\caption{Figure 8 - Figure 8: China \% of revenue for main analog players China \% of revenue}
\end{figure}
\begin{figure}[htbp]
\centering
\includegraphics[width=0.88\linewidth]{figures/fig_085.jpg}
\caption{Figure 9 - Figure 9: China growth y-o-y by companies shows significant difference Figure 10: China grew c7\% y-o-y in 2025 but is set to grow more in line with the market going forward}
\end{figure}
\begin{figure}[htbp]
\centering
\includegraphics[width=0.88\linewidth]{figures/fig_086.jpg}
\caption{Figure 11 - Figure 11: When accounting for share loss, we estimate incumbents' revenue growth should be stronger outside China 2025-29E}
\end{figure}
\begin{figure}[htbp]
\centering
\includegraphics[width=0.88\linewidth]{figures/fig_087.jpg}
\caption{Figure 14 - Figure 14: When accounting for content growth, we estimate inventory digestion started in Q4'23 and we expect an improving trend through the year}
\end{figure}
\begin{figure}[htbp]
\centering
\includegraphics[width=0.88\linewidth]{figures/fig_088.jpg}
\caption{Figure 15 - Figure 15: Auto/industrial OEM inventory days stabilizing}
\end{figure}
\begin{figure}[htbp]
\centering
\includegraphics[width=0.88\linewidth]{figures/fig_089.jpg}
\caption{Figure 16 - Figure 16: Total inventories '000 units: inventories were up 3\% during 1Q26 vs 1Q25 levels. This is back to a relatively high level for Stellantis for example}
\end{figure}
\begin{figure}[htbp]
\centering
\includegraphics[width=0.88\linewidth]{figures/fig_090.jpg}
\caption{Figure 17 - Figure 17: We expect next quarter to grow by 9\% q-o-q after an increase of 4\% in Q1'26...}
\end{figure}
\begin{figure}[htbp]
\centering
\includegraphics[width=0.88\linewidth]{figures/fig_091.jpg}
\caption{Figure 18 - Figure 18: ...Q2'26E likely to deliver 21\% y-o-y growth after 20\% in Q1'26}
\end{figure}
\begin{figure}[htbp]
\centering
\includegraphics[width=0.88\linewidth]{figures/fig_092.jpg}
\caption{Figure 21 - Figure 21: China's PV inventory changes (monthly)}
\end{figure}
{\small\color{refgray}\textbf{References:} [R007] UBS：汽车半导体复苏已至，但真正的分歧在“中国放缓”与“工业超预期”之间; [R021] UBS：人形机器人供应链的“确定性”正在从概念走向订单，但市场仍低估了传统业务的反转弹性}

\section{医疗政策与邮轮：2026中期选举是医疗政策拐点，奢华邮轮品牌定位是真正护城河}
\textbf{医疗保健政策的下一个拐点不在2028年总统大选，而在2026年中期选举；奢华邮轮市场的真正稀缺资产不是船，而是品牌定位的不可替代性。}

\begin{itemize}
  \item Bernstein指出，若民主党在2026年拿下众议院（预测市场给出82\%概率），核心议题将是恢复Medicaid和ACA市场 enrollment。
  \item Medicaid导向的保险公司（Centene、Molina等）已有所反映，但医院股尚未充分price in政策转向。
  \item Bernstein认为，奢华邮轮正站在邮轮行业最强的结构性需求风口上，美国55岁以上人群持有近75\%家庭财富。
  \item Viking的差异化定位在于“思考者”而非奢华感，2025年回头客率达54\%，客户评分行业领先。
  \item Viking占豪华邮轮新增运力的50\%以上，新船成本控制出色，每舱位成本与主流品牌相当但收益率更高。
  \item 奢华酒店新增供应正在放缓，高端邮轮正好承接溢出需求，供给约束为定价权提供支撑。
\end{itemize}
\begin{figure}[htbp]
\centering
\includegraphics[width=0.88\linewidth]{figures/fig_142.jpg}
\caption{EXHIBIT 1 - EXHIBIT 1: Trump's approval rating}
\end{figure}
\begin{figure}[htbp]
\centering
\includegraphics[width=0.88\linewidth]{figures/fig_143.jpg}
\caption{Exhibit 5 - EXHIBIT 2: Change in House seats in mid-term elections for President's party Changes in House seats during Mid-term elections for President's party}
\end{figure}
\begin{figure}[htbp]
\centering
\includegraphics[width=0.88\linewidth]{figures/fig_144.jpg}
\caption{EXHIBIT 3 - EXHIBIT 3: Chance to Win House Majority}
\end{figure}
\begin{figure}[htbp]
\centering
\includegraphics[width=0.88\linewidth]{figures/fig_145.jpg}
\caption{EXHIBIT 4 - EXHIBIT 4: Midterm House Seats Projection}
\end{figure}
\begin{figure}[htbp]
\centering
\includegraphics[width=0.88\linewidth]{figures/fig_146.jpg}
\caption{EXHIBIT 5 - EXHIBIT 5: Prediction market on who will will the US House?}
\end{figure}
\begin{figure}[htbp]
\centering
\includegraphics[width=0.88\linewidth]{figures/fig_147.jpg}
\caption{EXHIBIT 6 - EXHIBIT 6: Senate - Key races for 2026 midterms EXHIBIT 7: Chance to Win Senate Majority}
\end{figure}
\begin{figure}[htbp]
\centering
\includegraphics[width=0.88\linewidth]{figures/fig_148.jpg}
\caption{EXHIBIT 8 - EXHIBIT 8: Midterm Senate Seats Projection}
\end{figure}
\begin{figure}[htbp]
\centering
\includegraphics[width=0.88\linewidth]{figures/fig_149.jpg}
\caption{EXHIBIT 9 - EXHIBIT 9: Prediction market on who will win the US Senate}
\end{figure}
\begin{figure}[htbp]
\centering
\includegraphics[width=0.88\linewidth]{figures/fig_150.jpg}
\caption{EXHIBIT 10 - EXHIBIT 10: ACA favorability as per KFF Health Tracking Polls}
\end{figure}
\begin{figure}[htbp]
\centering
\includegraphics[width=0.88\linewidth]{figures/fig_151.jpg}
\caption{EXHIBIT 12 - EXHIBIT 12: Uninsured rate in the US US Uninsured rate}
\end{figure}
\begin{figure}[htbp]
\centering
\includegraphics[width=0.88\linewidth]{figures/fig_207.jpg}
\caption{Exhibit 4 - EXHIBIT 1: Luxury cruise industry ```mermaid graph LR}
\end{figure}
\begin{figure}[htbp]
\centering
\includegraphics[width=0.88\linewidth]{figures/fig_208.jpg}
\caption{EXHIBIT 2 - EXHIBIT 2: Viking is priced at the lower end of the luxury cruise market, behind smaller-ship luxury brands such as Silversea, Seabourn and Regent. Nightly prices for sailings aboard yachts range well into the \textbackslash{}\$000s, with an Aman}
\end{figure}
\begin{figure}[htbp]
\centering
\includegraphics[width=0.88\linewidth]{figures/fig_209.jpg}
\caption{EXHIBIT 3 - EXHIBIT 3: Viking and Oceania have the largest average ship capacity in the luxury cruise segment at \textbackslash{}\textasciitilde{}950 berths, over double the size of the largest yacht operated by Ritz-Carlton Yacht Collection Average Ship Capacity}
\end{figure}
\begin{figure}[htbp]
\centering
\includegraphics[width=0.88\linewidth]{figures/fig_210.jpg}
\caption{EXHIBIT 4 - EXHIBIT 4: Customers are willing to pay significantly more for cruises aboard smaller ships, but Viking commands a price premium relative to its average ship size Cruise Price vs Ship Capacity}
\end{figure}
\begin{figure}[htbp]
\centering
\includegraphics[width=0.88\linewidth]{figures/fig_211.jpg}
\caption{Exhibit 5 - EXHIBIT 6: They also make a point of being differentiated from the mainstream cruise lines (“no casinos”) and also other luxury cruise lines (“no butlers or white gloves”)}
\end{figure}
\begin{figure}[htbp]
\centering
\includegraphics[width=0.88\linewidth]{figures/fig_212.jpg}
\caption{EXHIBIT 7 - EXHIBIT 7: Viking's unique offering is well-loved by their customers, with an average review score ahead of the rest of the luxury segment Average guest review score - Luxury Ocean}
\end{figure}
\begin{figure}[htbp]
\centering
\includegraphics[width=0.88\linewidth]{figures/fig_218.jpg}
\caption{Exhibit 16 - EXHIBIT 13: The average age of luxury cruise goers is roughly ten years older than guests aboard mainstream cruise brands Average Guest Age by Cruise Brand}
\end{figure}
\begin{figure}[htbp]
\centering
\includegraphics[width=0.88\linewidth]{figures/fig_219.jpg}
\caption{EXHIBIT 14 - EXHIBIT 14: Over the next decade, the US population is set to add \textbackslash{}\textasciitilde{}22m over 45s but see virtually no growth in under 45s}
\end{figure}
\begin{figure}[htbp]
\centering
\includegraphics[width=0.88\linewidth]{figures/fig_220.jpg}
\caption{EXHIBIT 15 - EXHIBIT 15: Almost 75\% of wealth owned by those over 55, with this share up more than 15\% since 1990}
\end{figure}
\begin{figure}[htbp]
\centering
\includegraphics[width=0.88\linewidth]{figures/fig_221.jpg}
\caption{EXHIBIT 16 - EXHIBIT 16: Over the last three decades, the wealthiest Americans have become wealthier with the top quintile's share of wealth increasing from 61\% to 72\% Share of Household Wealth (Net Worth) by Income Quintile}
\end{figure}
\begin{figure}[htbp]
\centering
\includegraphics[width=0.88\linewidth]{figures/fig_222.jpg}
\caption{EXHIBIT 17 - EXHIBIT 17: The rich are getting richer - he highest income groups have also seen the greatest increase in income US income by quintile, indexed to 2004}
\end{figure}
\begin{figure}[htbp]
\centering
\includegraphics[width=0.88\linewidth]{figures/fig_223.jpg}
\caption{EXHIBIT 18 - EXHIBIT 18: The highest income groups have seen the largest reduction in hours worked - with more time for leisure Change in average hours worked on an average day (by income group)}
\end{figure}
\begin{figure}[htbp]
\centering
\includegraphics[width=0.88\linewidth]{figures/fig_224.jpg}
\caption{EXHIBIT 19 - EXHIBIT 19: Travel agents expect Luxury cruise to deliver the strongest growth Travel advisor expected booking growth compared to one year ago}
\end{figure}
\begin{figure}[htbp]
\centering
\includegraphics[width=0.88\linewidth]{figures/fig_225.jpg}
\caption{Exhibit 28 - EXHIBIT 20: US hotel supply growth is firmly below 1\%}
\end{figure}
\begin{figure}[htbp]
\centering
\includegraphics[width=0.88\linewidth]{figures/fig_226.jpg}
\caption{EXHIBIT 21 - EXHIBIT 21: Luxury hotel rooms under construction have tumbled, with owner returns under pressure}
\end{figure}
\begin{figure}[htbp]
\centering
\includegraphics[width=0.88\linewidth]{figures/fig_227.jpg}
\caption{EXHIBIT 22 - EXHIBIT 22: Luxury hotel rooms under construction relative to supply has plummeted vs pre pandemic}
\end{figure}
\begin{figure}[htbp]
\centering
\includegraphics[width=0.88\linewidth]{figures/fig_228.jpg}
\caption{EXHIBIT 23 - EXHIBIT 23: Luxury hotels are already close to maximum occupancy US hotels - occupancy}
\end{figure}
\begin{figure}[htbp]
\centering
\includegraphics[width=0.88\linewidth]{figures/fig_233.jpg}
\caption{EXHIBIT 28 - EXHIBIT 28: Luxury RevPAR has been significantly outperforming the rest of the industry}
\end{figure}
\begin{figure}[htbp]
\centering
\includegraphics[width=0.88\linewidth]{figures/fig_234.jpg}
\caption{EXHIBIT 29 - EXHIBIT 29: Luxury hotels have also had significantly more pricing power than any other chain scale US 12 month ADR by Chain Scale vs CPI (indexed Jan 2000)}
\end{figure}
{\small\color{refgray}\textbf{References:} [R011] Bernstein：医疗保健政策的下一个拐点不在2028，而在2026中期选举; [R017] Bernstein：邮轮市场的真正稀缺资产，不是船，是品牌定位的不可替代性}

\section{结语}
综合来看，当前市场处于多重结构性力量交汇的节点：AI与能源超级周期重塑产业格局，但宏观利率环境、地缘政治风险和资金流向分化构成核心约束。投资者需要超越总量叙事，聚焦于结构性分化中的定价偏差——无论是欧洲石油巨头隐含的低油价、韩国居民资产迁移的持续性，还是AI光互连PCB的量价齐升，这些被低估的变量可能定义未来6-12个月的资产回报格局。
\newpage
\section{报告覆盖清单}
\begin{itemize}
  \item [R001] 外汇与美元：美元多头共识已达峰值，历史规律指向未来三个月走弱 - GS：亚洲央行正在打一场“汇率防御战”，但真正决定胜负的变量不在利率
  \item [R002] 外汇与美元：美元多头共识已达峰值，历史规律指向未来三个月走弱 - GS：美元并非被高估，而是被错误定价——能源冲击的“项圈效应”才是核心
  \item [R003] 宏观与利率：美联储按兵不动，利率曲线重新定价路径被低估 - GS：市场真正低估的不是通胀本身，而是利率曲线重新定价的路径
  \item [R004] 地产与消费：中国楼市K型分化，香港写字楼迎结构性拐点 - GS：高频数据正在揭示一个被低估的“成本传导型”经济周期
  \item [R005] 区域市场：韩国居民资产结构性迁移，日本治理改善信号被低估 - GS：韩国股市的真正机会不在大盘，而在被低估的KOSDAQ
  \item [R006] 宏观与利率：美联储按兵不动，利率曲线重新定价路径被低估 / 中国信贷与金融：总量超预期但结构疲弱，流动性正常化传导加速 - NOM：中国利率市场真正被低估的风险不是降息预期，而是流动性正常化的传导加速
  \item [R007] 汽车半导体与人形机器人：复苏已至，但结构性分化是关键 - UBS：汽车半导体复苏已至，但真正的分歧在“中国放缓”与“工业超预期”之间
  \item [R008] 中国信贷与金融：总量超预期但结构疲弱，流动性正常化传导加速 - GS：五月信贷超预期，但真正的信号藏在结构里
  \item [R009] 外汇与美元：美元多头共识已达峰值，历史规律指向未来三个月走弱 - NOM：美元多头共识已达峰值，历史规律指向未来三个月大概率走弱
  \item [R010] 宏观与利率：美联储按兵不动，利率曲线重新定价路径被低估 - NOM：美联储的“按兵不动”比降息或加息更值得市场重新定价
  \item [R011] 医疗政策与邮轮：2026中期选举是医疗政策拐点，奢华邮轮品牌定位是真正护城河 - Bernstein：医疗保健政策的下一个拐点不在2028，而在2026中期选举
  \item [R012] 电动车与供应链：中国渗透率结构性跃迁，电力设备回调是重新定价 - JPM：全球电动车市场的真正拐点不在销量增速，而在中国渗透率的结构性跃迁
  \item [R013] 风险偏好与资金流向：全球资金重新定价“安全”，而非追逐风险 - GS：全球资金正在重新定价“安全”，而非追逐风险
  \item [R014] 未被主题摘要引用 - JPM：2026世界杯的真正价值不在门票，而在北美服务生态的定价权重估
  \item [R015] 能源与大宗：欧洲石油巨头隐含油价远低于共识，厄尔尼诺精准打击作物窗口 - GS：欧洲石油巨头隐含的长期油价远低于市场共识，这才是真正的价值锚点
  \item [R016] 区域市场：韩国居民资产结构性迁移，日本治理改善信号被低估 - GS：日本市场真正值得关注的不是AI反弹，而是股东会投票率背后的治理信号
  \item [R017] 医疗政策与邮轮：2026中期选举是医疗政策拐点，奢华邮轮品牌定位是真正护城河 - Bernstein：邮轮市场的真正稀缺资产，不是船，是品牌定位的不可替代性
  \item [R018] 医疗与教育：AI放射科从可选项变必需品，教育市场结构分化加速 - Bernstein：AI在放射科的落地，真正被低估的不是技术，而是结构性供需缺口
  \item [R019] 未被主题摘要引用 - 摩根斯坦利：市场低估的不是消费复苏，而是亚洲工业超级周期对中国出口的锚定效应
  \item [R020] 电动车与供应链：中国渗透率结构性跃迁，电力设备回调是重新定价 - JPM：中国电力设备在美渗透的拐点比市场预期的更快到来
  \item [R021] 汽车半导体与人形机器人：复苏已至，但结构性分化是关键 - UBS：人形机器人供应链的“确定性”正在从概念走向订单，但市场仍低估了传统业务的反转弹性
  \item [R022] 医疗与教育：AI放射科从可选项变必需品，教育市场结构分化加速 - GS：中国教育市场最值得跟踪的不是增长，而是结构分化的加速
  \item [R023] AI与科技：资本开支结构性集中，光互连与PCB成为下一波赢家 - GS：资本开支的“结构性集中”正在重塑全球工业投资逻辑
  \item [R024] 地产与消费：中国楼市K型分化，香港写字楼迎结构性拐点 - JPM：中国房地产并非整体回暖，而是一场K型分化下的资产再定价
  \item [R025] 地产与消费：中国楼市K型分化，香港写字楼迎结构性拐点 - 摩根斯坦利：市场对楼市“小阳春”的可持续性存在显著误判
  \item [R026] 未被主题摘要引用 - 摩根斯坦利：重卡销量超预期背后，真正的分化信号才刚刚开始
  \item [R027] 地产与消费：中国楼市K型分化，香港写字楼迎结构性拐点 - 摩根斯坦利：香港地产市场真正被低估的不是住宅，而是写字楼的结构性拐点
  \item [R028] 区域市场：韩国居民资产结构性迁移，日本治理改善信号被低估 - 摩根斯坦利：韩国居民资产迁移到股票市场的趋势具有结构性，而非周期性
  \item [R029] 电动车与供应链：中国渗透率结构性跃迁，电力设备回调是重新定价 - JPM：电力设备板块的20-30\%回调，不是风险出清，而是结构性机会的重新定价
  \item [R030] 能源与大宗：欧洲石油巨头隐含油价远低于共识，厄尔尼诺精准打击作物窗口 - 摩根斯坦利：市场真正低估的不是厄尔尼诺本身，而是它对作物窗口的精准打击
  \item [R031] 地产与消费：中国楼市K型分化，香港写字楼迎结构性拐点 - UBS：酒店业真正的拐点不在RevPAR，而在定价策略的深层分化
  \item [R032] AI与科技：资本开支结构性集中，光互连与PCB成为下一波赢家 - 摩根斯坦利：AI光互连的下一波赢家不在光模块，而在电路板
  \item [R033] AI与科技：资本开支结构性集中，光互连与PCB成为下一波赢家 - 摩根斯坦利：市场低估的不是AI投资放缓，而是日本线缆企业供给侧的差异化定价权
  \item [R034] 未被主题摘要引用 - Bernstein：AI购物离“自主决策”还差一个支付环节
\end{itemize}
\section{逐篇报告摘录}
\subsection{[R001] GS：亚洲央行正在打一场“汇率防御战”，但真正决定胜负的变量不在利率}
{\small\color{refgray}归类：外汇与美元：美元多头共识已达峰值，历史规律指向未来三个月走弱}

GS：亚洲央行正在打一场“汇率防御战”，但真正决定胜负的变量不在利率

当一场地缘冲突持续进入第四个月，最初被市场视为“短期冲击”的能源价格上涨，已经演变为对亚洲货币的系统性压力测试。过去三个月，除人民币外的所有亚洲货币对美元均出现贬值，韩元、印尼盾、菲律宾比索和印度卢比相继突破前期低点。外汇储备的普遍消耗，是这场防御战最直观的成本。

大多数市场讨论仍停留在“央行是否加息”的表层叙事上。但这份来自GS的研报揭示了一个更深层的判断：亚洲央行的政策工具箱正在发生结构性切换——从单纯消耗储备的“干预模式”，转向加息、资本流入管理和行政措施的组合拳。这不仅仅是应对方式的升级，更意味着各国对“汇率弱势容忍度”的重新定价。

真正值得关注的，不是哪家央行会加息，而是这场防御战之后，亚洲货币的相对强弱格局将如何被重塑。GS的核心结论是：在能源冲击得到明确解决之前，亚洲货币难以对美元出现显著升值。但在这个大背景下，不同经济体的分化路径已经清晰可见——而这恰恰是资产配置的关键。

理解亚洲央行当前的政策逻辑，不能只看利率决议。GS的报告指出，这轮政策响应的广度远超市场预期。以印尼央行为例，其在5月加息50个基点、6月非会议时间再加息25个基点之外，还同步推出了三项关键措施：提高SRBI各期限收益率、将外资对冲掉期利率降低10\%，以及强化“三重干预”政策（即干预即期、远期和境内债券市场）。

菲律宾央行则是加息政策最为激进的先行者——4月率先加息25个基点至4.50\%，GS预计还将连续四次加息。韩国央行虽然按兵不动，但转向了更鹰派的措辞，同时鼓励国民年金公团进行汇率对冲，这是典型的资本流动管理手段。

印度储备银行则采取了“利率不动、规则松绑”的策略：豁免外国投资者投资政府证券的资本利得税和利息所得税，将外资准入范围扩展至超长期限债券，并豁免银行发行外币债券和存款的相关限制。

这些措施的共同逻辑是：在无法直接控制能源价格和美联储政策的前提下，通过提高本币资产的吸引力来“截留”资本外流。从效果看，GS认为这些措施“在一定程度上有效”——至少为亚洲货币的下跌踩了刹车。

但这里有一个未被完全回答的问题：这些措施能持续多久？外汇储备的消耗数据显示，菲律宾、印度和印尼的储备在2月以来分别下降了7.5\%、4.8\%和4.0\%。当储备消耗到一定程度，央行加息的“可信度”本身也会受到质疑。这份报告没有展开讨论的是，如果能源价格继续维持高位，部分经济体的防御空间还剩多少。

在所有亚洲货币中，人民币的走势最为特殊——在整个中东冲突期间，人民币对美元持续升值，与DXY指数的区间震荡形成鲜明对比。GS明确指出，中国央行对“有序、渐进的升值”感到舒适。

这个判断背后有三个支撑点。第一，出口的持续超预期表现提供了基本面支撑。第二，央行估算的“企业超额美元囤积”峰值约为5000亿美元，而自去年12月以来，出口商已超额结汇约1500亿美元，这意味着升值预期已经转化为实际行为。第三，人民币国际化战略需要汇率保持吸引力，而年化约4\%的升值速度，恰好足以抵消外资的套利成本，又不会对出口竞争力造成过大拖累。

GS给出的预测路径是：人民币对美元在6个月内升至6.70，12个月内升至6.50。这个预测隐含的前提是，央行能够维持当前这种“半管制、半市场化”的节奏控制。

但报告没有充分展开的一个风险点是：如果国内需求持续疲软，而出口增速开始回落（例如由于全球需求放缓），央行是否还能维持当前的升值路径？GS经济学家认为“全面宽松的门槛很高”，但门槛高不等于不会触发。这是一个需要持续跟踪的关键假设。

这是报告中最为有趣的对比之一。韩国和台湾都受益于AI驱动的半导体出口繁荣，但汇率表现却完全相反：台币是亚洲货币中的相对强势者，而韩元则意外地跑输了所有亚洲同行。

问题出

[... middle omitted ...]

也净卖出约140亿美元台股，但最近一个季度转为净买入50亿美元，资本流动方向更为混合。更重要的是，台湾的经常账户盈余占GDP比例预计将从2025年创纪录的20\%升至2026年的25\%，这为台币提供了坚实的缓冲垫。

GS对韩元的基本面看法是“建设性的”，但认为在股票资金流出缓解之前，韩元难以跑赢亚洲同行。而台币则将继续相对强势，尽管央行偏好“渐进”节奏。

这个对比的启示是：在AI驱动的出口繁荣周期中，汇率表现取决于“出口收入能否转化为本币需求”。如果出口企业将收入留在海外或用于海外投资，再强的出口也无法支撑汇率。韩国的情况恰恰说明了这一点。

印尼是这份报告中风险信号最为密集的经济体。央行加息力度远超预期，但GS依然对印尼盾持看空态度。原因在于，印尼面临的不是单纯的汇率问题，而是三重压力的叠加。

第一重压力来自政策不确定性。政府计划将自然资源出口收入集中至国有银行，商业团体正在寻求政策澄清。虽然旨在提高税收和防止出口低开发票，但投资者担心增加的官僚程序会延迟出口。

\subsection{[R002] GS：美元并非被高估，而是被错误定价——能源冲击的“项圈效应”才是核心}
{\small\color{refgray}归类：外汇与美元：美元多头共识已达峰值，历史规律指向未来三个月走弱}

GS：美元并非被高估，而是被错误定价——能源冲击的“项圈效应”才是核心

这份最新发布的GS外汇策略报告，标题看似平淡——“项圈中的美元”——但如果你只看到它对美元短期走势的判断，就错过了真正的信号。报告的核心主张不在于预测美元涨跌，而在于揭示一个正在发生的、市场尚未充分定价的机制切换：全球外汇市场的定价权，正在从宏观风险偏好（AI叙事、商品价格）向国内利率差异（货币政策分化、实际利差）转移。这个转移，将重新定义未来6-12个月所有主要货币对的交易逻辑。

报告发布于一个微妙的时点。市场仍在消化中东冲突对能源供应的冲击，AI驱动的资本开支周期正经历第一轮预期修正，而美联储即将召开会议。在这样的背景下，GS团队没有给出简单的“买美元”或“卖美元”建议，而是提供了一个分析框架：理解当前外汇市场，需要同时把握两个维度——能源冲击的持续时间和国内政策的分化程度。前者决定美元的底部，后者决定货币对的相对强弱。

这份报告最有价值的判断是：市场真正低估的不是美元的方向，而是能源冲击正在从“脉冲式扰动”演变为“结构性约束”。这意味着，传统上依赖风险偏好和增长差异的外汇交易模型，可能需要重新校准。GS用两个图表清晰地展示了这种切换——从3-4月全球因素（股票、商品）主导汇率波动，到5月后利率差异重新成为关键驱动。这个信号，对于任何持有跨境资产或关注汇率风险敞口的决策者，都值得认真对待。

1. 能源冲击的“项圈效应”：收窄美元波动区间，而非决定方向

理解这份报告的关键，在于区分“能源冲击的规模”和“能源冲击的持续性”。GS的观点是，市场可能过于关注前者——即油价具体涨多少——而低估了后者带来的结构性影响。报告原文指出：“过热的油价正在降温，这让美元被困在项圈中。”这个判断的底层逻辑是：如果能源短缺是短暂且可控的，那么美元会因美国相对能源独立而获得避险溢价；但如果短缺变成持久战，反而会通过侵蚀全球增长预期，反过来压制美元。

这正是报告所说的“项圈效应”——能源冲击同时设定了美元的上限和下限。上限来自：持续的能源短缺拖累全球增长，最终反噬美国经济前景，限制美元进一步走强。下限来自：只要中东局势未彻底解决，能源供应的不确定性就为美元提供支撑。GS认为，当前市场对“能源缺口小于预期”的认知正在帮助稳定增长预期，这意味着美元短期不会大幅贬值，但也缺乏向上突破的催化剂。

这个框架对投资者的含义是：短期美元交易应关注波动率而非方向。报告特别指出，尽管宏观不确定性高企，但美元波动率（vol）并未同步上升——这本身就是一个信号，表明市场在等待新的催化剂，无论是美联储政策路径的重新定价，还是能源局势的实质性变化。对于套利交易者而言，当前环境更适合卖出波动率；对于对冲需求者而言，则需要警惕“项圈突然收紧”的风险。

2. 人民币的“结构性强弱”：贸易顺差与汇率低估的双重支撑，但动能正在切换

这份报告对人民币的判断，可能是对亚洲投资者最有价值的部分。GS明确看多人民币，目标价6.50（12个月），但更值得关注的是其论证逻辑——它没有依赖传统的利差或资本流动分析，而是从中国贸易竞争力的结构性优势出发。

核心论据有三层。第一层是贸易顺差的韧性。5月中国贸易顺差重回1000亿美元以上，这在全球能源冲击背景下显得尤为突出。GS特别指出，尽管对欧洲和拉美的出口有所下降，但对其他地区的名义出口仍在增长，且主要受科技产品价格上涨推动。这暗示中国出口结构正在向高附加值领域迁移，而非简单的数量扩张。第二层是汇率的深度低估。GS估算人民币低估超过20\%（基于GSDEER和GSFEER模型），这意味着即使出现一定程度的实际升值，也不会损害竞争力。第三层是政策容忍度。报告认为，中国决策层对“渐进式升值”的容忍度正在提高，这从近期中间价和市场汇率的同步走低可以看出。

[... middle omitted ...]

期驱动”切换到“基本面驱动”。贸易顺差和汇率低估提供了硬支撑，但短期节奏取决于能源价格和全球需求的变化。这是一个适合逐步建立多头头寸而非追逐突破的窗口。

3. 英镑的“政治溢价陷阱”：选举事件的风险被高估，能源冲击才是真正变量

GS对英镑的分析提供了一个很好的案例，展示如何看待政治风险事件对汇率的影响。下周的Makefield补选被视为工党领袖Andy Burnham潜在领导力挑战的关键节点，但报告认为，这对英镑的冲击将是“相当狭窄的”。

为什么？因为政治溢价已经在市场中有所定价，但程度有限。GS构建的英镑政治风险溢价指标显示，尽管Burnham领导的政府概率在预测市场中已升至70\%以上，但英镑的反应一直相当温和。报告认为，要让政治溢价显著扩大并实质打压英镑，需要满足两个条件之一：要么Burnham公布具体的政策计划，让市场担忧英国财政轨迹；要么政治不确定性持续更长时间。但即使在这种情况下，GS认为其他更基本的因素——尤其是能源冲击的演变——将继续主导英镑走势。

\subsection{[R003] GS：市场真正低估的不是通胀本身，而是利率曲线重新定价的路径}
{\small\color{refgray}归类：宏观与利率：美联储按兵不动，利率曲线重新定价路径被低估}

GS：市场真正低估的不是通胀本身，而是利率曲线重新定价的路径

当全球投资者还在为地缘政治冲击下的油价波动焦虑时，一份来自GS全球利率交易团队的研报提出了一个更值得深思的框架：市场对短期通胀的定价已经相当充分，但真正未被消化的是利率曲线在多重路径下的重新定价。这份报告的核心判断是——美国利率曲线陡化的概率被系统性低估，而欧洲和英国的前端利率则存在被过度定价的加息风险。

这份报告发布于美联储新主席Warsh首次议息会议前夕，时间点本身就具有信号意义。全球利率市场正在经历一个罕见的“多路径均衡”阶段：油价在供给端冲击与需求端疲软之间拉扯，各国央行在“抗通胀的尾声”与“经济增长的脆弱”之间摇摆，而新一任美联储主席的沟通风格本身就是一个不确定变量。

GS分析师在报告中明确指出，市场定价隐含的加息概率已经包含了过多的鹰派预期。对于美国，30-35个基点的加息风险定价意味着市场已经计入了多次加息的可能性，而GS认为更可能的情景是“长期按兵不动”。对于欧元区和英国，情况则更加复杂——通胀的持续性而非短期油价飙升，才是决定央行政策路径的真正变量。

这篇文章将沿着GS的分析框架，拆解三个核心问题：为什么美国利率曲线陡化被低估、欧洲和英国的利率定价存在什么结构性偏差、以及新一任美联储主席的沟通方式可能如何改变市场游戏规则。最后，我们将指出报告尚未完全回答的几个关键问题，这些将决定投资者能否真正利用这份研报的洞见。

1. 美国利率市场最大的误判，是以为“加息风险”就是全部风险

GS研报的第一个核心洞察，直指美国利率市场当前定价的结构性偏差。报告指出，美国前端利率的加息定价已经从高点回落，曲线也从近期低点有所陡化，但市场仍然在接下来一年内定价了30-35个基点的加息风险。这个数字意味着市场已经隐含了多次加息的概率。

但GS的判断恰恰相反：长期按兵不动的可能性远高于加息，而中期降息的可能性也不应被完全排除。这个判断基于两个逻辑。第一，核心通胀的构成数据虽然比预期略强，但GS经济学家对中期核心通胀的前景仍然保持温和判断。第二，油价的下行风险——尤其是美伊协议的可能性——如果兑现，将显著改变通胀预期。

这意味着什么？如果GS的判断正确，当前市场定价中的加息溢价将在未来逐步消退，1y1y利率的下行风险显著大于上行风险。而最受益于这种情景的，是利率曲线的“红绿部分”——即中期和长期债券。GS建议，对鹰派风险和油价上行风险的最佳战术对冲，是建立那些能够从“加息风险提前兑现”中受益的头寸，而非简单做空利率。

这个判断的关键假设在于：美联储是否会真正转向“平衡”的沟通指引。报告提到，4月会议已经出现了转向平衡指引的支持声音，而新主席Warsh的首次会议将决定市场如何解读这个信号。如果新主席强调劳动力市场缺乏通胀压力，或者强调剔除波动后的通胀指标，市场将解读为鸽派信号，支持曲线陡化。反之，如果强调通胀风险的警惕信号，则会被视为鹰派。

2. 真实利率的“beta”被低估了，这是当前最值得关注的相对价值交易

GS的第二个核心洞察，聚焦于一个被多数投资者忽视的维度：真实利率与名义利率之间的beta关系。报告指出，美国利率的宽松定价主要来自通胀预期的下降，而长期真实收益率仅略低于2月高点。这意味着，如果通胀预期进一步下降，真实收益率的下行空间将非常可观。

GS在这里提出了一个具体的交易框架：做多10年期通胀互换，同时以0.25倍的名义权重做空10年期SOFR。这个交易的逻辑在于，真实利率的波动性远高于名义利率，而当前市场定价中，大部分风险溢价都集中在真实曲线中。如果通胀预期继续下行，真实利率的下降幅度将超过名义利率；如果通胀预期回升，名义利率的上行也会被真实利率的上升部分对冲。

这个交易的核心假设是：美联储的反应函数不会发生根本性变化。

[... middle omitted ...]

”

GS的欧洲利率分析，揭示了一个容易被市场忽略的结构性问题。报告指出，欧洲央行成为第一个在伊朗冲突后加息的全球主要央行，但这背后的驱动力不是短期油价飙升，而是通胀的持续性。欧洲央行的最新预测显示，通胀的“持续性”正在成为更棘手的问题。

这里有一个关键的市场信号：尽管前端利率曾因油价上涨而大幅重新定价，但最近油价下跌后，利率几乎没有回落。GS分析师发现，更好的参考指标是远期的z6油价——这个指标最近一直持平，反映的是更持久的通胀风险。这意味着，即使短期油价回落，欧洲利率的下行空间也可能有限。

但报告同时指出，欧洲央行的“主动收紧”策略与其他G4央行形成对比，这意味着欧元利率曲线（尤其是1年以上部分）在抛售中可能表现更好。GS维持做多5y5y真实OIS的头寸，预期名义利率下行和通胀曲线陡化将同时发生。

对于欧元区投资者，这个判断的含义是：不要因为短期油价下跌就盲目做多欧洲前端利率。通胀的持续性意味着欧洲央行可能比市场预期的更鹰派，而前端利率的下行空间可能被高估。

\subsection{[R004] GS：高频数据正在揭示一个被低估的“成本传导型”经济周期}
{\small\color{refgray}归类：地产与消费：中国楼市K型分化，香港写字楼迎结构性拐点}

当前市场对中国经济的讨论，大多仍停留在“需求复苏有多强”这个框架里。房地产销售是否企稳、消费信心是否回升、出口订单是否可持续——这些问题的答案当然重要，但它们可能正在让人忽略一个更深层的结构性变化。

GS在6月12日发布的这份中国经济高频追踪周报中，释放了一个值得反复推敲的信号：中国经济的短期波动，正在从“需求驱动”向“供给成本驱动”切换。 进口价格、出口价格、工业品价格之间的联动关系，正在形成一个过去几年罕见的传导链条。这一链条的走向，将直接决定2025年下半年到2026年的企业盈利分配、通胀路径以及政策空间。

这份报告的核心价值，不在于它提供了多少数据点，而在于它用一种“高频+多维”的视角，捕捉到了传统月度或季度数据容易平滑掉的拐点。以下是我们从这份报告提炼出的四个核心洞察，以及一个尚未被充分回答的关键问题。

报告中的高频消费数据，呈现出一种罕见的“背离”状态。一方面，Morning Consult的消费者信心指数在5月攀升至多年新高，甚至超过了2021年的水平。另一方面，真实世界的消费活动，至少在几个关键维度上，并没有同步走强。

最值得关注的是房地产市场。报告显示，30城新房日均成交面积在6月已跌至低于去年同期水平。而16城二手房成交量虽然仍高于去年，但环比已出现明显下滑。新房与二手房之间的“温差”，暗示市场正在经历一个“以价换量”的尾声阶段——二手房的放量，更多依赖业主降价，而非真实需求的持续涌入。

与此同时，国内航班量虽然有所回升，但依然低于去年同期；航班取消率虽有下降，但绝对水平仍然偏高。这说明居民的跨区域流动意愿，并未完全恢复到疫情前的活跃程度。汽车销售数据也印证了这一点：5月新能源汽车销量环比增长，但与去年同期相比仍为负值，整体汽车销量同样低于去年水平。

这些数据合在一起，指向一个判断：消费的“韧性”更多体现在高收入群体的信心指标上，而大众消费的“弹性”正在被压缩。 房地产市场的流动性改善，尚未转化为价格的企稳或新开工的扩张。对于依赖地产链的行业而言，下半年的需求端支撑可能比市场预期的要弱。

报告的工业生产指标，呈现出一种“量稳价升”的微妙状态。钢铁需求环比小幅下滑，周度产量基本持平，但同比仍低于2019年水平。沿海省份的煤炭日耗从近期高点回落，这通常被视为工业活动边际放缓的信号。

然而，与“量”的平淡形成对比的，是“价”的剧烈波动。报告重点追踪了进口价格和出口价格的变化。5月数据显示，以美元计价的进口价格全面上涨：煤炭、原油、成品油、天然气、初级塑料、集成电路的进口价格指数，均出现了显著跳升。其中，原油进口价格指数从3月的140跃升至5月的155，成品油从145升至170，集成电路从140升至165。

更值得警惕的是，出口价格也开始跟进。成品油出口价格指数从3月的110飙升至5月的165，集成电路从148升至210。铝的出口价格也在温和上涨。

这意味着什么？中国正在经历一轮“输入型成本推动”向“出口价格传导”的加速过程。 上游能源和原材料的涨价，已经不再是企业可以内部消化的成本压力，而是开始通过出口价格向外转移。这种传导，对于下游制造企业的毛利率意味着什么，报告虽然没有直接计算，但逻辑上是清晰的：那些无法将成本完全转嫁的企业，将面临利润挤压；而那些具备议价能力的企业，则可能受益于价差扩大。

报告将能源价格冲击作为当前追踪的核心变量，这一点值得高度关注。GS之所以将这一追踪频率提升至每周，本身就说明他们认为这是一个“非常规”的变化窗口。

从报告展示的国内汽柴油价格来看，尽管国际能源价格大幅波动，但国内成品油定价机制在短期内保持了相对稳定。然而，化工品价格已经开始反应。聚丙烯价格在过去一周小幅上涨，硫酸、甲醇等基础化工品的价格也在走高。这些化工品是塑料、化肥、纺织等下游行业

[... middle omitted ...]

”，出口价格往往滞后于进口价格，且涨幅较小。但这一次，集成电路的出口价格涨幅（从3月的148到5月的210，涨幅约42\%）已经超过了进口价格涨幅（从140到165，涨幅约18\%）。

这是一个非常规的信号。它可能意味着，中国在集成电路等高端制造领域的出口，正在从“数量型增长”转向“价值型增长”。或者说，中国出口的产品结构，正在从低附加值向高附加值迁移。这种迁移，如果能够持续，将显著改善中国的贸易条件——即用更少的出口量，换取等量的进口额。

但报告也提示了另一面的风险：成品油的出口价格涨幅（从110到165，涨幅约50\%）几乎与进口价格同步，说明炼化环节的利润空间并未显著扩大。而铝的出口价格涨幅有限，暗示传统大宗商品的定价权仍然不强。

*对于投资者而言，真正的机会可能不在于“中国出口总量是否增长”，而在于“哪些行业的出口价格能够跑赢进口成本”。** 那些具备技术壁垒、品牌溢价或供应链优势的行业，将在这一轮成本传导中脱颖而出。而那些依赖价格竞争的行业，压力将显著加大。

\subsection{[R005] GS：韩国股市的真正机会不在大盘，而在被低估的KOSDAQ}
{\small\color{refgray}归类：区域市场：韩国居民资产结构性迁移，日本治理改善信号被低估}

当全球投资者都在盯着KOSPI的波动和外资流出时，一个更值得关注的信号正在被忽视：KOSDAQ正在吸引结构性的外资流入，而其背后的驱动因素——盈利修正的背离与估值的错位——恰恰构成了一个战术性alpha机会。这是GS最新一期韩国市场周报中，最值得认真对待的一个判断。

这份报告发布于市场持续波动的背景下。KOSPI上周小幅下跌0.5\%，外资净卖出4,400亿韩元，主要集中在科技和造船板块。与此同时，KOSDAQ逆势上涨2.7\%，外资净流入67亿韩元。韩元兑美元单周升值2.8\%，韩国股票风险指标进一步滑入风险厌恶区域。这些数字本身并不惊人，但将它们放在一起看，一个清晰的图景浮现出来：市场正在重新定价韩国股市的内部结构，而大多数投资者可能还停留在对大盘的单一关注中。

GS的核心论点并不复杂，但需要认真拆解。自2023年8月以来，KOSDAQ相对于KOSPI持续跑输，累计差距已相当显著。但与此同时，KOSDAQ的盈利修正却远弱于KOSPI——KOSPI的12个月远期每股盈利上修幅度显著强于KOSDAQ。更关键的是，尽管KOSDAQ大幅跑输，其估值却仍然比KOSPI更贵：KOSDAQ的12个月远期市盈率远高于KOSPI，而其盈利修正却弱得多。这种“估值溢价+盈利弱势”的组合，通常意味着市场对KOSDAQ的定价存在某种结构性偏差。

然而，外资的行为却在讲述另一个故事。今年以来，KOSPI遭遇了大规模外资流出，但KOSDAQ却吸引了持续的外资流入。GS分析认为，KOSPI的外资流出很大程度上与投资组合再平衡有关，尤其是两只最大的半导体股票贡献了大部分卖压。而流入KOSDAQ的外资则高度集中在信息技术和医疗保健板块。这种资金流向的分化，与盈利修正和估值的背离形成了有趣的对照。

这意味着什么？简单说，市场可能正在对KOSPI的某些大盘股进行被动减持，同时主动加仓KOSDAQ中那些与AI、半导体设备和医疗创新主题相关的标的。这并非简单的“避险”或“追逐小盘”，而是一种有主题逻辑的战术轮动。

GS报告中最值得关注的数据之一，是KOSPI与KOSDAQ之间盈利修正的巨大差异。从Exhibit 2可以看出，KOSPI的盈利上修轨迹自2023年以来持续走强，而KOSDAQ的盈利修正则明显疲弱，甚至在某些时期出现下修。这种背离并非短期现象，而是持续了超过一年半的结构性趋势。

从逻辑上讲，盈利修正的强弱直接反映企业基本面的边际变化。KOSPI的盈利上修主要来自半导体、化工等周期性板块的改善，而KOSDAQ的盈利弱势则可能与其更依赖国内消费和中小企业景气度有关。韩国宏观经济的复苏并不均匀，出口导向的大企业受益于全球半导体周期和汇率优势，而内需相关的中小企业仍在承压。

但这里有一个关键问题：盈利修正的背离是否已经充分反映在价格中？GS的数据显示，KOSDAQ的估值溢价并未因盈利弱势而显著收窄。这意味着市场对KOSDAQ的定价中，可能包含了对未来盈利改善的预期，或者对其增长属性的溢价。如果未来盈利修正无法跟上，这种估值溢价就可能面临压缩风险。反之，如果KOSDAQ的盈利开始改善，当前的估值水平反而可能提供安全边际。

对于产业决策者和投资者而言，这意味着需要区分两类标的：一类是KOSDAQ中真正受益于结构性趋势、盈利有望改善的公司；另一类是纯粹因小盘流动性或主题炒作而被高估的标的。GS的报告暗示，前者可能正是外资持续买入的对象。

KOSPI的外资流出规模令人瞩目。上周单周净卖出4,400亿韩元，今年以来累计流出更为显著。但GS的分析指出，这些流出并非均匀分布，而是高度集中在两只最大的半导体股票上。Exhibit 5清晰地展示了这一现象：KOSPI的外资流出曲线与三星电子和SK海力士的流出曲线几乎重合。

这背后的逻辑值得深思。GS认

[... middle omitted ...]

1，进一步滑入风险厌恶区域。这通常意味着市场情绪已相当悲观，从历史经验看，极端风险厌恶往往对应着阶段性底部。

与此同时，KOSDAQ的外资流入虽然绝对规模不大，但方向性意义更重要。在KOSPI整体流出的背景下，KOSDAQ能够吸引净流入，说明有主动资金在寻找结构性机会。这些资金集中在信息技术和医疗保健板块，与全球AI产业链和生物科技投资主题高度吻合。

GS报告中最令人困惑但也最有洞察力的图表，是Exhibit 3。它显示，尽管KOSDAQ自2023年以来大幅跑输KOSPI，但其12个月远期市盈率仍然高于KOSPI，且高于2010年以来的历史均值。换句话说，KOSDAQ的估值并没有因为跑输而变得便宜，反而仍然维持着溢价。

这种“跑输但估值不便宜”的组合，通常意味着两种可能性：一是KOSDAQ的盈利预期被过度下调，导致估值被动抬升；二是市场对KOSDAQ的成长性给予了持续溢价，即使盈利表现不佳。GS倾向于前一种解释，因为KOSDAQ的盈利修正确实明显弱于KOSPI。

\subsection{[R006] NOM：中国利率市场真正被低估的风险不是降息预期，而是流动性正常化的传导加速}
{\small\color{refgray}归类：宏观与利率：美联储按兵不动，利率曲线重新定价路径被低估 / 中国信贷与金融：总量超预期但结构疲弱，流动性正常化传导加速}

NOM：中国利率市场真正被低估的风险不是降息预期，而是流动性正常化的传导加速

过去两个月，中国债券市场经历了一轮罕见的“先涨后跌”行情。5月末，各期限利率下行4-6个基点，市场一度沉浸在“宽松延续”的叙事中。但进入6月，利率迅速反弹约4-5个基点，将此前涨幅几乎全部回吐。

市场参与者习惯性地将这一切归因于央行降息或政策信号。但NOM最新发布的这份中国利率策略报告，提供了一个更值得警惕的判断框架：真正驱动利率走向的，不是降息与否，而是一套被市场低估的流动性正常化机制正在加速运转。这份报告的核心主张是，过去几个月支撑债市极度宽松的流动性环境正在系统性逆转，而这一过程才刚刚开始。

NOM明确提出，维持做空3年期NDIRS的头寸，目标看向1.60\%。这个判断并非基于对经济数据的乐观，而是基于对流动性供给结构、债券发行节奏以及机构行为三重变量的拆解。以下是我们从这份研报中提炼出的五个核心洞察。

1. 央行已经在悄悄收紧中期流动性，而市场直到6月才反应过来

市场对央行操作的关注，往往集中在7天逆回购利率或MLF利率是否调整。但NOM指出，真正重要的变化发生在“中期货币政策工具”的净投放上。

数据显示，今年1-2月，央行通过MLF、买断式逆回购、国债净买入等工具，累计净投放流动性约2.05万亿元，远超历史均值。这解释了为什么一季度市场流动性极度充裕。但从3月开始，情况逆转：3月净回笼200亿元，4月净回笼560亿元，5月净回笼850亿元。到了6月5日，央行又通过3个月期买断式逆回购净回笼300亿元。结果就是，截至6月初，今年以来的累计中期净投放已降至仅140亿元，远低于 年约1.1万亿元的平均水平。

这个数字意味着什么？意味着过去两个月支撑债市上涨的“流动性充裕”基础，其实已经在逐步被抽走。但市场对此反应滞后——隔夜利率DR001在5月中旬就开始上行，而7天回购利率DR007和1年期NCD收益率直到6月5日才开始显著反弹。NOM的判断是，这种滞后恰恰说明，市场对流动性正常化的定价还不够充分。

这里有一个关键信号值得注意：6月3日和4日，央行连续两日暂停逆回购操作。上一次暂停还要追溯到2024年8月。NOM的解读是，央行并非要制造紧缩，而是在传递一个清晰的信号——当市场资金已经足够便宜时，它不会再主动注入更多流动性。这个信号的意义，可能比一次降息或加息更大。

2. 政府债券供给将在三季度集中放量，这是流动性最直接的“抽水机”

流动性收紧的另一面，是供给端即将到来的压力。NOM的数据显示，截至5月底，国债和地方债的净融资进度仅完成全年额度的35\%。虽然国债发行节奏略快于历史同期，但地方债的发行明显偏慢——前五个月的完成比例低于 年38\%的平均水平。

5月政府债券净供给尤其低，这恰好与当月流动性极度充裕形成共振，也是推动基金大规模买入长久期债券的重要原因之一。但NOM预计，6月国债和地方债的月净供给将从5月的890亿元上升至约1000亿元，三季度进一步升至约1300亿元，四季度回落至约900亿元。

这个节奏意味着什么？三季度将是供给压力最大的时期。当大量债券集中发行，银行体系需要吸收这些供给，流动性自然收紧。更重要的是，财政支出加速本身也可能改善经济增长预期，从而进一步压低债市的避险需求。

对投资者而言，这里有一个容易被忽视的因果关系：不是经济数据好转导致利率上行，而是债券供给本身通过流动性渠道推高了利率。这个逻辑在二季度可能比基本面更重要。

NCD（同业存单）市场的变化，往往是银行体系流动性状况的“温度计”。NOM观察到，自2025年6月以来，月度NCD净发行多数时间为负值，原因是央行此前持续注入中期流动性，银行无需通过NCD市场主动负债。但5月，这一趋势逆转——月净NCD供给自去年10月以来首次转正

[... middle omitted ...]

和流动性收紧，最终会体现在机构行为上。NOM的持仓数据显示，在4-5月大量积累超长久期国债和政策性银行债后，6月以来基金已经开始减持。第一周，基金还在加仓10年以上品种，同时净卖出前端和中段；到了本周，卖出已经蔓延至后端，基金大量抛售10年期政策性银行债。

但这里有一个重要的对冲力量：保险机构。数据显示，保险机构本周明显增加了对30年期国债的买入。这与它们的历史操作模式一致——在市场抛售时加仓，在上涨时减仓。NOM认为，如果30年期国债收益率升至2.25\%，保险的买入需求将进一步增强，从而对收益率上行形成一定封顶。

此外，目前10年期与30年期国债的利差约为49个基点，处于较宽水平。这一利差本身为30年期品种提供了缓冲，也意味着长端的上行空间可能不如短端那么直接。

NOM的结论是：在信用需求和零售数据等经济指标出现实质性改善之前，或者流动性收紧的幅度远超当前预期之前，利率曲线不太可能出现显著陡峭化。因此，他们更倾向于通过3年期NDIRS而非5年期来押注利率上行。

\subsection{[R007] UBS：汽车半导体复苏已至，但真正的分歧在“中国放缓”与“工业超预期”之间}
{\small\color{refgray}归类：汽车半导体与人形机器人：复苏已至，但结构性分化是关键}

UBS：汽车半导体复苏已至，但真正的分歧在“中国放缓”与“工业超预期”之间

汽车半导体行业正在走出长达两年的去库存周期。UBS在最新发布的研报中确认，2026年Q1汽车半导体收入同比增速达到7.8\%，这是两年多来首次出现明确的同比增长。模拟芯片收入更是同比增长20\%，且这一势头预计将贯穿整个2026年。

但这份报告的真正价值不在于确认复苏，而在于揭示一个多数市场参与者尚未完全定价的结构性分化：中国市场的需求正在放缓，而工业领域的增长正在超预期加速。这两股力量的方向相反，但都在同一个时间窗口内作用于全球汽车半导体公司的业绩与估值。

UBS维持了对2026年汽车半导体市场10\%同比增长的预测，但明确提醒，中国1-5月零售数据已出现15\%-20\%的同比下滑，其中新能源汽车下滑10\%-15\%。如果这一趋势在下半年引发库存修正，全球汽车半导体增长可能面临显著的下行风险。

与此同时，工业领域的半导体需求正在被AI相关产品拉动，UBS将2026年工业半导体收入增速预测从16\%上调至25\%。这意味着，那些同时覆盖汽车和工业市场的模拟芯片公司，其业绩分化将比行业整体更加剧烈。

复苏是真实的，但复苏的结构并不均匀。对于投资者而言，问题不再是“要不要配置汽车半导体”，而是“在哪些公司、哪些下游、哪些区域中配置”。

UBS的分销商追踪数据显示，Q1多家模拟芯片公司已开始提价，包括TI、英飞凌和恩智浦。提价的直接原因是通胀成本传导，但其隐含意义远不止于此。

在去库存周期的尾声阶段，市场普遍预期2026年模拟芯片价格将出现低个位数下降。但UBS指出，目前的价格走势表明，实际定价可能持平甚至略好于预期。这意味着，此前被计入估值模型的利润率假设可能偏低。

对于模拟芯片公司而言，价格稳定或温和上涨带来的利润弹性，在当前的估值水平下尚未被充分反映。UBS覆盖的模拟芯片公司当前交易在约30倍12个月前瞻市盈率，而十年均值仅为19倍。估值溢价已经较高，但若涨价周期持续且利润超预期，这一溢价可能获得基本面支撑。

不过，UBS也谨慎地指出，涨价并非全行业现象。供应紧张主要集中在AI相关产品领域，更广泛的模拟芯片市场仍面临需求分化和局部过剩。这意味着涨价周期的可持续性需要逐季验证。

2. 汽车复苏的结构性差异：中国占全球20\%-30\%的需求，但正在成为拖累

汽车半导体的复苏是真实的，但UBS的模型拆解显示，这一复苏高度依赖非中国市场。2026年，UBS预计中国汽车半导体收入增长5\%，而非中国市场增长8\%。但如果考虑中国本土半导体公司的份额侵蚀，全球龙头公司在中国的汽车半导体收入可能持平。

这一判断的关键支撑来自中国1-5月的零售数据。全球汽车半导体需求的20\%-30\%来自中国，而这一市场正在经历显著的放缓。如果这一趋势在H2引发渠道库存修正，UBS的10\%全年增长预测可能面临下调风险。

这里需要特别关注的是，中国市场的放缓并非周期性，而是结构性。中国新能源汽车渗透率已处于高位，补贴退坡后的需求自然回落，叠加本土芯片公司的替代加速，将共同压缩全球龙头在中国市场的增长空间。

UBS的数据显示，2025年中国市场占英飞凌收入的30\%，占瑞萨的31\%，占ADI的26\%。对于这些公司而言，中国市场的增速放缓将直接拉低整体增长，而这一影响在当前的估值中是否充分反映，是一个值得追问的问题。

3. 工业超预期正在重塑行业增长结构，但AI驱动的需求能否持续仍需验证

工业半导体是本次报告中最大的向上修正项。UBS将2026年工业半导体收入增速从16\%大幅上调至25\%，2027年从10\%上调至16\%。Q1工业收入同比增长26\%，主要驱动力来自模拟芯片公司的强劲表现。

这一增长的核心推手是AI。多家公司在季度报告中上调了AI相关产品的指引，并获得

[... middle omitted ...]

含了过多乐观预期。

UBS的数据显示，模拟芯片公司相对于半导体大盘股的估值折扣已从一年前的-2\%扩大至-17\%。从表面看，这似乎为模拟芯片提供了安全边际。但UBS在报告中的措辞是“sector preferences - remaining positive but more selective”，即保持积极但更加挑剔。

这一态度的转变值得注意。在行业复苏初期，整体配置逻辑成立；但随着估值抬升和基本面分化加剧，选股的重要性开始超过仓位判断。

UBS最偏好的标的包括TI、瑞萨和STMicroelectronics，均为买入评级。而ON Semiconductor、英飞凌和Melexis则被给予中性评级。这一组合透露出两个信号：一是UBS更看好那些在工业领域有更强暴露的公司；二是对中国市场依赖度较高的公司被相对谨慎对待。

30倍市盈率意味着市场已经定价了较为乐观的增长路径。任何低于预期的数据，无论是中国市场的进一步恶化，还是工业AI需求的降温，都可能引发估值修正。

\subsection{[R008] GS：五月信贷超预期，但真正的信号藏在结构里}
{\small\color{refgray}归类：中国信贷与金融：总量超预期但结构疲弱，流动性正常化传导加速}

市场对五月金融数据的初步解读，多数停留在“总量超预期”这个层面。新增社融2.03万亿，新增人民币贷款5200亿，双双高于GS和彭博的预测值。如果只看这些数字，似乎可以得出“宽信用正在见效”的乐观结论。

但GS这份研报的核心判断恰恰相反：总量超预期不是复苏信号，而是结构性疲弱的又一次确认。

原因在于，数据超预期主要来自银行信贷的供给端发力，而非实体融资需求的真实回暖。当一份信贷报告需要用“票据融资大幅增加”和“中长期企业贷款收缩”来解释时，决策者应当警惕的不是流动性不足，而是信用传导机制中的梗阻。

这份研报最有价值的洞察，不是告诉读者信贷数据超预期了，而是拆解了“为什么超预期反而令人担忧”。对于关注中国经济周期和资产定价的投资者而言，理解这种结构性的分化，远比记住一个总量数字重要。

五月新增人民币贷款5200亿，高于GS预测的4000亿和彭博一致预期的4500亿。表面看是银行放贷意愿增强，但细拆结构会发现，这更像是银行在监管和考核压力下的被动行为。

企业贷款是五月信贷的主要支撑，新增6400亿，高于去年同期的5300亿。但关键在于构成：票据融资贡献了5570亿，短期贷款1000亿，而中长期企业贷款反而净减少200亿。去年同期，中长期企业贷款是净增3300亿。

这个对比意味着什么？中长期贷款通常对应企业的资本开支、设备投资和扩张计划，是衡量真实融资需求的核心指标。当这个指标从净增3300亿转为净减200亿，而票据融资大幅飙升时，最合理的解释是：企业没有投资意愿，银行只能用低成本票据融资来“填满”信贷额度。

票据融资本质上是银行间的一种短期资金安排，与实体经济活动的关联度较低。它的增加往往反映的是银行在完成信贷投放指标，而不是企业有真实的融资需求。GS研报虽然没有直接点明这一点，但数据本身已经给出了清晰的信号。

对于投资者而言，这意味着：如果下一阶段票据融资占比继续上升，而中长期贷款持续疲弱，那么信贷总量对经济增长的领先意义将进一步削弱。

五月居民贷款存量净减少1410亿，而去年同期是净增加540亿。这个对比的幅度，比企业贷款的结构变化更令人担忧。

居民贷款的收缩，在逻辑上可以归因于两个因素：一是房地产市场的持续低迷导致按揭需求不足，二是消费信心的疲弱使得信用卡和消费贷款需求下降。GS研报没有展开讨论具体原因，但数据本身已经表明，居民部门正在主动或被动地去杠杆。

从宏观角度看，居民去杠杆与企业的短期融资行为构成了一个相互强化的负反馈循环：居民收入预期不稳，减少消费和购房，导致企业收入下降、投资意愿减弱；企业投资收缩又进一步影响就业和居民收入。信贷数据中的居民贷款收缩，正是这个循环在金融层面的映射。

对于资产定价而言，居民贷款的趋势比企业贷款更具前瞻性。因为居民信贷直接关联消费和房地产市场，而这两个领域是中国经济内循环的核心驱动力。如果居民贷款在接下来的几个月继续收缩，那么市场对经济复苏的预期可能需要进一步下修。

五月M1同比增速从4月的5.0\%回升至5.5\%，M2则持平于8.6\%。乍看之下，M1回升似乎意味着企业活期存款增加，是企业经营现金流改善的信号。

但GS研报提供了一个关键的补充视角：五月财政存款增加7100亿，比去年同期少了约1700亿。这意味着财政支出节奏比去年更快，而这很可能是M1回升的主要驱动力。

这个判断的逻辑链条是：财政支出加快，资金从国库转入企业和居民账户，增加了活期存款，从而推高了M1。如果剔除财政因素，M1的改善幅度可能远没有数据显示的那么显著。

对于投资者而言，这个区分至关重要。由财政驱动的M1回升，其持续性和对经济的支撑力度，远不如由企业自发经营改善驱动的M1回升。一旦财政支出节奏放缓，M1增速可能再次回落。因此，不能将五月M1的回升解读为企业端复苏的信号。

[... middle omitted ...]

会加速。

对于债券投资者而言，这意味着信用利差的走势可能需要重新评估。如果社融增速持续放缓，而信贷结构依然以短期融资为主，那么信用风险的定价逻辑可能需要从“总量宽松”转向“结构分化”。

5. 报告没有完全回答的关键问题：信贷需求何时才会真正改善？

GS这份研报清晰地指出了问题所在，但对于“信贷需求何时才会真正改善”这个问题，并没有给出明确的答案。这并非报告的缺陷，而是当前经济环境下，任何预测都面临高度不确定性的体现。

从逻辑上推演，信贷需求的改善需要满足几个条件：一是房地产市场的企稳，二是居民收入预期的修复，三是企业投资回报率的回升。目前这三个条件都不具备。五月的数据表明，政策层面可能仍在通过银行体系维持信贷供给，但需求端并未响应。

这里有一个值得投资者持续跟踪的观察框架：关注中长期企业贷款和居民贷款的边际变化，而不是总量数据。 只有当这两个指标同时出现连续两个月的环比改善，才能初步判断信贷需求正在回暖。在此之前，任何基于总量数据的乐观解读都可能是误导性的。

\subsection{[R009] NOM：美元多头共识已达峰值，历史规律指向未来三个月大概率走弱}
{\small\color{refgray}归类：外汇与美元：美元多头共识已达峰值，历史规律指向未来三个月走弱}

NOM：美元多头共识已达峰值，历史规律指向未来三个月大概率走弱

当前市场对美元的共识有多强？NOM最新外汇策略报告给出了一个值得警惕的答案：当“美国例外论”成为客户会议和媒体高频词时，美元多头头寸正在逼近一个危险的临界点。

这份研报的核心判断并非否定美元近期的强势逻辑——美国数据持续超预期、利率重新定价、资本流入加速，这些支撑因素仍然存在。但NOM提出了一个反直觉的历史规律：当美国经济惊喜指数突破60这一阈值后，美元在未来三个月下跌的概率高达75\%。当前该指数正处于近三年高位，这意味着市场可能站在美元情绪周期的拐点。

更关键的是，NOM识别出三个可能触发美元逆转的具体风险点：7月美国就业数据的季节性残留效应、新任美联储主席沃什可能释放鸽派信号、以及AI驱动的科技股叙事面临成本质疑。这些因素尚未被市场充分定价，但它们正在积累。

对于持有美元敞口的投资者而言，当前需要关注的不是美元还能涨多少，而是什么会打破这个共识。

1. 经济惊喜指数已达历史极端区间，美元后续回报历史规律指向弱势

NOM的核心论据建立在一个简洁而有力的统计规律上。报告追溯了美国经济惊喜指数过去23年的全部历史，发现该指数突破60水平共出现41次。当考察随后三个月的美元表现时，结果呈现出高度一致的负面特征：美元兑G10等权重指数的平均回报为-1.8\%，而“命中率”——即美元走弱的概率——高达73.2\%。

这个规律在分货币对层面同样显著。美元兑欧元、英镑、澳元、新西兰元、挪威克朗和瑞典克朗在三个月窗口期的平均跌幅均超过1.5\%，命中率在63\%至80\%之间。即使是对美元相对抗跌的加元和日元，三个月平均回报也分别为-1.2\%和-0.7\%。

NOM进一步指出，如果将阈值提高到70，美元在所有表现窗口期的回报都变得更加一致地为负。这意味着当前水平——惊喜指数在60附近徘徊——正处于从正面转向负面的“临界点”附近。

需要强调的是，这个历史规律并非经济基本面的直接反映，而是市场情绪和定价的滞后效应。当数据惊喜达到极端水平时，市场已经充分甚至过度定价了这些好消息，后续边际上的数据改善空间收窄，而任何不及预期的数据都会触发多头平仓。这正是NOM所说的“惊喜的惊喜”——好消息本身可能成为美元反转的催化剂。

2. 美元多头仓位尚未到极端水平，但非美货币空头积累已构成不对称风险

从持仓结构看，当前美元多头并非历史极值。CFTC非商业持仓数据显示，美元净多头占未平仓合约的比例约为8\%，远低于2025年初的27\%。这似乎表明市场尚未过度拥挤。

但NOM提醒，更值得关注的是非美货币的空头积累。英镑、日元、加元、瑞士法郎等G10货币的空头头寸都在显著增加，尤其是英镑和日元，净空头占比分别达到32\%和37\%。这种“全面做空非美”的持仓结构，实际上构成了一个不对称的风险场景：一旦美元出现回调迹象，这些空头头寸的集中平仓将放大美元的下跌幅度。

这与2025年初的情况形成对比——当时美元多头更为极端，但非美空头相对分散，因此美元调整时非美货币的反弹幅度相对温和。当前的结构意味着，如果触发因素出现，美元的下行速度可能比历史均值更快。

报告还提到，中东地缘政治局势的缓和预期——尤其是特朗普关于“周末签署协议”的表态——可能成为美元回调的短期催化剂。在当前的持仓结构下，任何地缘政治风险的下降都会削弱美元的避险溢价，而空头回补压力会进一步放大这一效果。

NOM明确列出了三个可能打破当前美元共识的具体风险，这些风险并非小概率事件，而是有明确时间窗口和触发条件的潜在催化剂。

第一个风险是美国就业数据的季节性残留效应。过去三年，7月非农就业数据均出现显著的下行意外，而2026年7月的数据将于8月7日公布。在连续三次非农超预期之后，如果这一模式重现，将直接冲击“美国劳动力

[... middle omitted ...]

分点，但NOM注意到一个关键变化：越来越多的AI消费企业开始对成本提出质疑，而科技公司正通过股权融资而非现金流来支撑资本支出。Oracle在6月11日宣布超预期资本支出后股价下跌，就是一个信号。如果AI叙事从“增长故事”转向“成本压力故事”，可能通过财富效应和金融条件收紧两条路径对美元形成负面冲击。

这三个风险中，任何一个单独出现都可能只是短期扰动，但如果叠加发生——例如7月就业数据意外下行，叠加沃什鸽派表态，再叠加AI股回调——将构成一个完整的美元逆转叙事。

尽管NOM的逻辑链条清晰，但这份报告留下了一个关键问题：如果美元确实在未来三个月走弱，调整幅度会有多大？持续多久？

从历史数据看，当惊喜指数突破60后，三个月平均跌幅约为1.8\%，但这个数字掩盖了很大的离散度。在2009年和2022年，美元在类似场景下反而大幅走强，分别上涨17.5\%和8.5\%，这与当时全球金融环境（欧债危机、美联储激进加息）密切相关。这说明历史规律并非机械的，需要结合当前的宏观背景判断。

\subsection{[R010] NOM：美联储的“按兵不动”比降息或加息更值得市场重新定价}
{\small\color{refgray}归类：宏观与利率：美联储按兵不动，利率曲线重新定价路径被低估}

市场正在为一个并不存在的降息周期做准备。这并非危言耸听，而是NOM最新一期美国经济周报的核心判断。该报告详细拆解了6月美联储议息会议可能呈现的决策图景，其结论清晰而坚定：联邦基金利率将在当前水平停留至2027年底。对于习惯了“加息-暂停-降息”线性思维的市场参与者而言，这意味着一场对资产定价逻辑的根本性修正。

报告提供的不仅仅是一次会议的预测，而是一整套关于美联储新常态的叙事。其中包含了几个关键信号：鹰派措辞的回归、点阵图中位数的大幅上移、以及新任主席沃尔什对前瞻指引这一传统工具的潜在解构。这些信号叠加在一起，指向一个远比市场当前定价更为严峻的利率环境。

1. 六月议息会议的核心动作不是加息，而是清除所有关于降息的暗示

NOM预期，6月FOMC会议将维持利率不变，这将是连续第四次按兵不动。但真正重要的政策信号并非利率本身，而是会后声明的措辞调整。报告明确指出，声明将删除所有关于“宽松倾向”（easing bias）的前瞻指引语言，且这一变化将以全票通过的方式呈现。

这意味着，美联储将主动关闭市场对未来降息的任何想象空间。近期多位FOMC与会者已经公开讨论加息的可能性，鹰派声音的密集程度在过去几个月中显著上升。尽管委员会内部仍存在一个“鸽派阵营”，对通胀前景保持谨慎乐观，并认为当前政策具有适度限制性，但这一阵营的声音正在被整体转向鹰派的共识所淹没。

对于资产定价而言，这带来的直接含义是：此前市场所消化的任何降息预期，都需要被重新校准。如果美联储连“暗示降息可能”的措辞都不再保留，那么基于降息逻辑的久期交易、成长股估值、以及新兴市场资本流动假设，都将面临根本性的动摇。

2. 点阵图将展示一个“无降息”的长期路径，中性利率的抬升是关键

报告对点阵图的预测是整篇分析中最具冲击力的部分。NOM预期，2026年和2027年的利率中位数预测将上调至3.625\%，意味着在未来一年半内，利率将保持在当前水平不变。而2028年的中位数预测为3.375\%，仅隐含一次降息。更值得关注的是，长期中性利率中位数将从3.125\%上调至3.250\%。

这一调整的深层含义在于：美联储参与者对“政策立场”的描述方式，正在倒逼他们上调对中性利率的估计。如果经济在高利率环境下依然表现出韧性，那么“限制性”的门槛就需要被重新定义。中性利率的上调，意味着未来任何一个降息周期的终点都会更高，利率中枢的系统性抬升正在成为现实。

报告还做出了一项极具洞察力的判断：新任主席沃尔什可能拒绝提交自己的利率预测。这将使点阵图中的“点”从19个减少到18个。沃尔什一贯批评前瞻指引的有效性，他的这一举动将不仅是个人偏好，更可能预示着美联储沟通方式的系统性变革。如果连点阵图这一最核心的市场沟通工具都开始失去其完整性，市场对美联储未来路径的预测将面临更大的不确定性。

报告对5月PCE通胀的预测同样令人警惕。基于CPI和PPI数据，NOM估计5月核心PCE通胀环比将加速至0.345\%，同比则攀升至3.426\%，这将是2023年10月以来的最高水平。

这一加速并非全面性的。报告指出，核心CPI商品通胀在5月转为负值，关税影响持续消退，全球半导体短缺对消费电子的价格压力也在减弱。但问题出在服务业。核心CPI服务通胀依然坚挺，租金相关分项居高不下，超级核心CPI分项持续以可观速度增长。交通运输服务价格（如机票和汽车保险）则表现出剧烈波动。

更值得关注的是上游价格压力。PPI数据显示，美国制造商的投入成本持续上升，这与近期工厂调查中“支付价格”指数的走高是一致的。尽管核心CPI商品价格有所软化，但制造商仍在将更高的成本转嫁给消费者——PPI成品消费品（剔除食品和能源）价格在5月继续上涨。这意味着，商品端的暂时缓和可能只是假象，上游成本压力正在积累，并将在未来几个

[... middle omitted ...]

通缩。他也可能将当前的高通胀定性为伊朗战争带来的暂时现象。如果市场接受了这一叙事，那么当前的高利率环境将被视为“等待战争结束”的过渡期，而非永久性的政策收紧。

但报告也提示了沃尔什可能带来的“破坏性”变化。他可能继续淡化核心PCE通胀的作用，转而提出替代性的通胀衡量指标。他可能讨论缩表的可能性，甚至为更小的资产负债表规模制定路线图。他可能对美联储的沟通策略进行改革。这些变化中的任何一项，都将对市场预期产生深远影响。

5. 报告尚未回答的关键问题：这场“暂停”会持续多久，触发因素是什么

NOM的报告在逻辑上极为严密，但它也留下了一个关键问题没有完全回答：美联储的“按兵不动”状态，其持续时间究竟是由什么决定的？

报告提到，沃尔什可能将暂停归因于伊朗战争带来的不确定性。但战争本身是一个高度不可预测的变量。如果战争迅速结束，通胀压力随之消退，美联储是否有空间提前转向？反之，如果战争持续并导致能源价格进一步飙升，美联储是否会被迫转向加息？报告对此没有给出明确的情景分析。

\subsection{[R011] Bernstein：医疗保健政策的下一个拐点不在2028，而在2026中期选举}
{\small\color{refgray}归类：医疗政策与邮轮：2026中期选举是医疗政策拐点，奢华邮轮品牌定位是真正护城河}

Bernstein：医疗保健政策的下一个拐点不在2028，而在2026中期选举

市场对医疗保健板块的关注，大多集中在2028年总统大选可能带来的政策摇摆上。但Bernstein这份最新研报提出了一个更紧迫的判断：真正影响资产定价的政策拐点，可能出现在2026年中期选举之后，而非四年后的大选周期。

这份报告的核心洞察在于：如果民主党在2026年拿下众议院，医疗保健政策的方向将发生实质性转变——从削减Medicaid和 Marketplace注册，转向恢复覆盖率和扩大保障。而这一转变对Medicaid导向的保险公司和医院的影响，远比市场当前定价所反映的要深远。

研报同时指出，尽管特朗普的民调支持率已降至第二任期内最低点，但众议院的席位争夺并非一边倒。Bernstein的分析框架显示，共和党在众议院仍有约200个席位的“地板”，这意味着民主党的席位增长可能不如历史典型的中期选举那样剧烈。正是这种“有限但关键”的转变，让政策博弈变得更加微妙——妥协的可能性上升，极端政策的概率下降。

对于投资者而言，这不仅仅是一次政治预测。它直接关系到Medicaid管理式医疗组织的风险池变化、坏账趋势、以及垂直整合监管的走向。Bernstein的评级表已经透露出信号：Centene、Molina、Humana等安全网MCO均获得“跑赢大盘”评级，而医院股HCA仅为“与大盘持平”。

以下是我们从这份报告中提炼出的五个关键层次，每一个都指向一个尚未被市场充分定价的变量。

1. 民主党重掌众议院将扭转Medicaid削减，但医院的定价尚未反映这一预期

Bernstein的分析指出，民主党若拿下众议院，首要政策目标将是恢复Medicaid资金和ACA增强补贴。这意味着此前因共和党政策导致的Medicaid注册下降和风险池恶化，有望得到缓解。对于Medicaid MCO而言，这直接关系到医疗损失率的改善和注册人数的稳定。

但研报特别强调了一个市场盲点：医院板块。虽然医院同样是这些政策的潜在受益者——坏账减少、未补偿护理负担下降——但当前医院股的股价表现并未体现出这一预期。Bernstein认为，这种定价偏差可能为投资者提供了机会窗口，前提是中期选举结果符合预测。

这一判断的深层含义在于：市场的注意力过度集中在Medicaid MCO上，而对医院端的影响尚未形成共识。如果民主党获胜，医院板块的估值重估可能比保险板块更具弹性，因为当前的预期差更大。

研报对参议院的预测更加谨慎。虽然民主党有50\%的概率拿下参议院，但需要赢得俄亥俄、得克萨斯等特朗普领先超过10个百分点的州，难度显著高于众议院。Bernstein的基准情景是共和党保持50-51席。

这种“众议院蓝、参议院红”的分裂格局，对政策的影响是决定性的。研报明确指出，无论民主党赢得一院还是两院，最终政策都需要妥协才能通过立法并避免否决。这意味着，极端政策——如对Medicare Advantage或PBM的严厉监管——在 年期间发生的概率较低。

这一判断对投资者的启示是：不要过度炒作“颠覆性政策”的风险。市场可能高估了民主党在医疗保健领域的行动能力，而低估了两党在扩大覆盖率和控制成本方面达成有限共识的可能性。

3. 2028年大选的政策焦点将从“废除”转向“限制”，垂直整合成为最大变量

Bernstein对2028年选举周期的早期分析，揭示了政策焦点的转移。与过去围绕“Medicare for All”的激进辩论不同，下一轮政策讨论将更加务实：覆盖率和降低未参保率、可负担性、以及垂直整合与AI监管。

研报特别指出，虽然“Medicare for All”不太可能成为实际政策，但支持该主张的候选人将在民主党初选中占据重要位置，这可能在2027年引发市场波动。然而，最实

[... middle omitted ...]

改善的幅度与时间线

Bernstein的框架在逻辑上自洽，但有一个关键变量尚未充分展开：政策落地到风险池改善之间的传导机制和时滞。

民主党恢复Medicaid资金和ACA补贴，理论上会改善Medicaid MCO的风险池，因为更多健康个体将重新进入保险池，降低医疗损失率。但这一过程需要时间——注册周期、费率调整、以及医疗利用模式的变化，都可能影响改善的幅度和速度。

研报对Centene和Molina的盈利预测显示， 年的每股收益增长路径相对清晰，但市场对这一改善的定价是否充分，仍取决于具体政策的实施细节。例如，恢复资金是通过直接拨款还是通过延迟OBAMACARE某些条款的实施？不同的政策路径将对风险池产生不同的影响。

这一未解问题，恰恰是投资者需要深入跟踪的核心变量。政策方向明确，但传导路径的细节决定了资产定价的精确性。

Bernstein这份报告最值得借鉴的，不是它对中期选举的预测，而是它的分析框架：将政治概率转化为资产定价的输入变量，而非单纯的背景噪音。

\subsection{[R012] JPM：全球电动车市场的真正拐点不在销量增速，而在中国渗透率的结构性跃迁}
{\small\color{refgray}归类：电动车与供应链：中国渗透率结构性跃迁，电力设备回调是重新定价}

JPM：全球电动车市场的真正拐点不在销量增速，而在中国渗透率的结构性跃迁

5月全球电动车销量数据出来了。环比增长9\%，同比增长8\%，看起来是一个温和复苏的信号。但JPM这份最新研报真正值得关注的点，不是总量数字，而是一个被多数人低估的结构性变化：中国BEV（纯电动车）渗透率在4月达到42\%之后，5月稳稳守住了这一水平，同时总EV（含插混）渗透率攀升至63\%，创下2025年8月以来的新高。

这不是一个简单的季节性波动。在整体乘用车市场同比下降22\%的背景下，BEV销量反而实现了5\%的同比增长。这意味着电动车正在从一个“政策驱动的新兴市场”转变为一个“即使在宏观逆风中也能自主增长的成熟品类”。对于锂矿、电池材料、乃至整个汽车产业链的投资者而言，这个信号比任何短期销量预测都重要。

报告同时揭示了一个更大的叙事：全球三大市场中，只有中国在完成真正的渗透率跃迁。欧盟虽然同比增长强劲（32\%），但渗透率仍停留在24\%的低位；而美国市场BEV渗透率连续五个月纹丝不动地停在6\%。这三条线之间的分歧，正在重新定义全球锂需求的地理分布和定价逻辑。

1. 中国市场的核心变化不是销量数字，而是电动车正在变成“新燃油车”

JPM的数据显示，5月中国BEV销量63.7万辆，环比增长10\%，但同比增速仅为5\%。单纯看这个数字，容易得出“增长放缓”的结论。但真正有价值的信息藏在一个对比中：在总乘用车销量同比下降22\%的恶劣环境下，BEV销量逆势增长。这意味着电动车已经从“可选消费”变成了“刚需替代”。

更关键的是渗透率结构。5月中国BEV渗透率42\%，与4月持平。但总EV渗透率达到63\%，这意味着每卖出三辆车，就有两辆是带插电的。这个数字在2025年5月仅为53\%。一年之内，中国汽车市场的电动化率提高了10个百分点。

这种跃迁的含义是什么？它意味着中国汽车市场正在经历一个不可逆的品类替代过程。当BEV渗透率超过40\%后，传统燃油车的规模效应开始加速瓦解——4S店网络收缩、零部件供应链重组、二手燃油车残值下降，这些因素会反过来进一步推动消费者转向电动车。JPM研报虽然没有直接论述这一点，但从数据趋势可以推断，中国市场的电动化已经从“政策驱动”进入了“自我强化”阶段。

对于投资者而言，这意味着中国锂需求的波动性正在降低。过去市场担忧的是“补贴退坡导致销量断崖”，现在这个风险已经显著下降。即使宏观经济继续承压，电动车的渗透率也不太可能大幅回落。

欧盟-10地区5月BEV销量21.4万辆，同比增长39\%，YTD累计同比增长32\%。这个增速在所有区域中最亮眼。但JPM研报同时指出，欧盟BEV渗透率从4月的24\%提升至5月的32\%。这个跳升幅度很大，需要仔细拆解。

报告给出的解释是“国家层面的补贴方案和可负担的大众电动车进入市场”。但仔细看数据，欧盟总乘用车销量5月仅增长3\%，说明BEV的增长并非来自整体市场的扩张，而是来自燃油车的直接替代。这与美国市场形成鲜明对比——美国总销量持平，但BEV渗透率纹丝不动。

欧盟市场的特殊性在于，它同时受益于两个不同层面的推动力：一是政策端，各国补贴方案正在加码；二是供给端，大众、Stellantis等厂商的平价电动车开始放量。这两个因素叠加，使得欧盟成为当前全球电动车市场中增长确定性最高的区域。JPM团队将欧盟描述为“最强区域”，这个判断有数据支撑。

但这里有一个尚未被充分讨论的问题：欧盟的补贴可持续性如何？部分国家的补贴方案有明确的时间窗口，一旦退坡，渗透率是否会出现类似中国2023年初的短期波动？报告没有展开，但这可能是下半年需要持续跟踪的关键变量。

美国5月BEV销量8.6万辆，环比增长10\%，但YTD同比下降25\%。更值得注意的是，BEV渗透率已经连续五个月停留在6\%。从2025年1

[... middle omitted ...]

者而言，美国市场的停滞意味着全球锂需求的结构正在发生变化。过去市场习惯用“中国+欧洲+美国”的线性外推来预测锂需求，但如果美国长期停留在6\%的渗透率，全球锂需求的天花板就需要重新评估。这可能是这份报告隐含的最大风险信号。

JPM在报告最后做了一个关键的判断：“考虑到储能市场的持续强劲和供给的时间滞后，我们认为锂辉石价格将从这里面临温和的上行压力。”这是一个明确的看多信号，但它的前提条件值得仔细推敲。

这个判断的逻辑链条是这样的：电动车销量虽然增速放缓，但绝对量仍在增长；同时储能市场（ESS）需求强劲，这两个因素叠加，导致锂需求依然坚挺。而供给端由于矿山投产的时间滞后，短期内无法快速响应，因此价格有上行空间。

但这里有一个尚未被充分讨论的变量：中国BEV渗透率在42\%之后，增速是否会自然放缓？如果渗透率从42\%提升到50\%需要更长的时间，那么锂需求的增量主要将来自欧盟和储能市场，而非中国。这意味着，锂价的上行空间将高度依赖于欧盟补贴政策的可持续性和美国市场的突破。

\subsection{[R013] GS：全球资金正在重新定价“安全”，而非追逐风险}
{\small\color{refgray}归类：风险偏好与资金流向：全球资金重新定价“安全”，而非追逐风险}

全球资金流向正在传递一个比表面数据更复杂的信号。截至6月10日当周，全球股票基金净流入310亿美元，固收基金净流入177亿美元，看起来是典型的risk-on格局。但真正值得关注的不是总量，而是结构：资金正在同时涌入美国科技股和短期债券，却在撤出中国A股和欧洲股票。这不是简单的“风险偏好回升”，而是全球投资者在重新定义什么才是当前环境下真正的“安全资产”。

这份GS最新发布的资金流向周报，揭示了一个正在发生的资产再定价过程。它告诉我们，市场并非在无差别地追逐收益，而是在有选择性地为确定性支付溢价。对于产业决策者和资产配置者而言，理解这个结构比盯着大盘涨跌重要得多。

当周全球股票基金净流入310亿美元，连续第二周保持高位。但拆开来看，驱动力量几乎完全来自美国市场。欧洲股票基金持续净流出，新兴市场内部也出现显著分化：台湾和韩国股票基金是EM净流入的主要来源，而全球EM基准基金和中国大陆股票基金则录得净流出。

这意味着什么？全球投资者的权益配置正在向“美国科技叙事”集中，而这一叙事正在与地缘政治风险、供应链重构等议题脱钩。台湾和韩国的资金流入并非简单的EM轮动，而是全球科技产业链资金流向的延伸——它们是美国科技投资的“卫星市场”。

对于中国企业而言，这种资金结构意味着：即便整体EM资金面改善，中国大陆股票仍然面临结构性流出压力。这不是周期性问题，而是全球资本重新评估中国市场风险溢价的结果。报告显示，中国大陆股票基金的资金流出在当周仍在持续，且来自境外投资者的流出尤为明显。

固收基金持续获得强劲资金流入，但最值得关注的是品种结构。短期债券基金和通胀保护债券基金是流入最为稳定的类别。这不是投资者在“配置债券”，而是在“锁定短期确定性和通胀对冲”。

与此同时，新兴市场硬通货债券基金却录得净流出。这形成了一组鲜明对比：全球投资者愿意买入美国短期国债和TIPS，却不愿意持有新兴市场美元债。背后的逻辑是——投资者对主权信用风险的定价正在发生系统性变化。美国债务上限谈判的暂时解决并没有消除市场对长期财政可持续性的担忧，但短期工具提供了可预测的回报，而新兴市场美元债则同时面临汇率、利率和主权信用三重不确定性。

对于资产配置者而言，这组数据提示：固收市场的机会正在向短端和通胀保护端集中，而长端利率债和新兴市场债券可能面临持续的流动性折价。

3. 外汇流动揭示的“安全资产”定义正在从美元扩展至韩元和匈牙利福林

跨境外汇流动方面，美元和韩元是当周净需求最强的货币，而印度卢比、巴西雷亚尔和人民币则录得最大净流出。这一格局与权益和固收资金流向高度一致：资金流向美国科技产业链（韩元）和美国本土资产（美元），远离中国和部分EM国家。

但更值得关注的信号来自匈牙利。报告特别指出，自年初以来，匈牙利债券的流入显著增加。GS的分析认为，这与匈牙利经济政策即将转向——包括采用欧元的可能性——有关，并据此判断匈牙利债券收益率将逐步向欧元区收敛，福林有进一步升值的不对称风险。

这是一个典型的“边缘市场因制度变革而获得重新定价”的案例。匈牙利债券的流入并非简单的EM追逐利差，而是投资者在押注一个具体政策路径的兑现。这也提示我们：在当前的全球资金环境中，能够提供“制度确定性”或“政策锚”的市场，正在获得结构性溢价。

4. 报告尚未完全回答的关键问题：中国大陆股票流出的底部在哪里

尽管GS的数据清晰地展示了中国大陆股票基金的资金流出趋势，但报告并未对流出何时见底给出明确判断。从历史数据看，中国股票基金的资金流出已经持续多个季度，且境外投资者的参与度持续下降。但这是否意味着中国资产已经被“过度抛售”，还是说这一趋势才刚刚进入中期？

报告没有提供答案，因为这取决于几个尚未被充分验证的假设：中国经济增长的内生动力能否在短期内显著改善？地缘政治风

[... middle omitted ...]

金流出的逻辑。当前全球资金正在系统性回避三样东西：欧洲权益（增长预期疲弱）、中国股票（风险溢价上升）以及新兴市场美元债（信用和汇率双重不确定性）。

这种回避不是临时的战术调整，而是全球资产配置者在重新定义“安全”的边界。当美国科技股和短期债券成为资金的双向避风港时，其他资产类别要想吸引资金，必须提供比过去更强的确定性——无论是通过政策改革（如匈牙利）、产业链绑定（如台湾和韩国），还是通过估值跌到足够有吸引力的水平。

对于读者而言，接下来的关键问题是：中国大陆股票需要跌到什么程度，才能重新进入全球配置者的“安全资产”清单？台湾和韩国的科技资金流入是否已经接近饱和？匈牙利模式是否可以被其他边缘市场复制？

这些问题的答案，需要更完整的报告数据和更深入的讨论。我们已经在知识星球和微信群里开始了对这些问题的拆解，欢迎继续加入交流。

Personal reading notes and learning share only. Not investment advice.

\subsection{[R014] JPM：2026世界杯的真正价值不在门票，而在北美服务生态的定价权重估}
{\small\color{refgray}归类：未被主题摘要引用}

JPM：2026世界杯的真正价值不在门票，而在北美服务生态的定价权重估

2026年世界杯将于2026年6月开赛，距现在正好一年。对于全球资本市场，这通常被归类为“事件驱动型主题”——赛事临近时炒作一波，结束后迅速退潮。但JPM最新发布的专题研报提出了一个值得严肃对待的判断：本届世界杯的独特性在于，它是历史上第一次由北美三个经济体联合主办、赛事规模从32队扩至48队、104场比赛覆盖16座城市的超级事件，其经济影响远非“一次性消费脉冲”所能概括。

这份报告的真正洞察，不在于预测世界杯期间美国GDP将增加172亿美元，而在于它揭示了北美服务生态（从住宿、出行到数字广告）正在经历一场结构性压力测试。市场目前的定价远远低估了这种压力测试对相关行业估值中枢的长期含义。

JPM在报告中反复强调一个数字：500万张门票申请对应约700万张可售座位，超额认购倍数约为70倍。这个数字本身是惊人的，但更重要的是它背后的两层含义。

第一，赛事扩编至48队、104场比赛，意味着比赛日从传统的约30天拉长到近40天。这对于主办城市而言，不是一周的狂欢，而是一个完整的高峰运营季。第二，78场比赛和11个主办城市集中在美国，且大量比赛位于所谓的“sprawl/drive markets”——即需要驾车跨城出行的市场。这意味着，世界杯的消费辐射不是点状的，而是网络状的。

报告引用住宿研究团队的估算：世界杯将为美国酒店业带来9.1亿美元的增量客房收入，对应30-40个基点的RevPAR增长，主办城市在2026年6-7月的RevPAR将提升7-25\%。这个数字看起来不大，但需要放在一个背景下来理解：美国酒店行业的RevPAR年增长率通常只有2-3\%。一个单一事件能在一个月内贡献接近全年十分之一的增量，这本身就是对行业定价权的重估信号。

*所以呢？** 市场目前对酒店、OTA、租车等板块的估值，仍然基于“常态需求曲线”的假设。世界杯带来的需求冲击，本质上是一个压力测试——它检验的是这些服务商在需求骤增时，能否将流量转化为真实的定价溢价。如果答案是肯定的，那么这些行业的长期估值锚点就需要上移。

2. 数字广告才是真正的“隐形冠军”，但市场的关注度严重不足

报告提到，JPM的互联网团队预计世界杯将带来50亿美元的全球增量广告支出，其中40亿美元流向数字渠道，数字广告渗透率达到73\%。这个数字值得反复咀嚼。

传统上，大型体育赛事的广告支出主要流向电视转播权。但73\%的数字渗透率意味着，本届世界杯的广告增量将主要被Google、Meta、TikTok等数字平台捕获，而非传统的广播电视网络。报告中的受益者篮子（JPRWC26B）包含了Alphabet和Meta，这并非偶然。

更深一层的逻辑是：世界杯期间的数字广告投放，不仅是品牌曝光，更是对平台广告系统在高并发、高竞争环境下的效率检验。如果这些平台能够证明它们能在全球最高流量的场景下维持甚至提升广告ROI，那么它们对品牌预算的长期吸引力将显著增强。

*但这里有一个报告没有完全展开的问题：** 50亿美元的增量广告支出，有多少是“新增”的，有多少是“转移”的？如果品牌只是将原本计划在2026年其他时段投放的预算集中到世界杯期间，那么对平台而言，这只是一个时间上的再分配，而非总量的结构性提升。报告没有给出这个拆解，而这恰恰是投资者判断广告板块持续性收益的关键。

JPM对出行和住宿板块的覆盖非常细致。从租车到网约车，从豪华酒店到经济型连锁，报告都给出了具体的增量预测。例如，网约车和外卖平台的预订量预计分别增加3.77亿美元和7300万美元。

这些数字本身是合理的，但我们需要问一个更尖锐的问题：这些增量利润，最终会留在平台端，还是被司机、房东或竞争性定价策略侵蚀掉？

以网约车为例。世界杯期

[... middle omitted ...]



JPM同时构建了一个世界杯赞助商篮子（JPRWC26S），并指出在过去两届世界杯中，赞助商股票的表现显著优于受益者篮子。

这个观察非常关键。它暗示了一个可能被市场忽略的逻辑：赞助商（如可口可乐、百威、Visa等）的受益方式，与直接服务提供商完全不同。 赞助商不依赖于世界杯期间的即时消费流量，而是通过品牌资产的全球曝光来获得长期市场份额的提升。这种“无形资产增值”的路径，更难被竞争或定价侵蚀，因此其股价表现的持续性更强。

从图1（Performance Lookback）可以看出，赞助商篮子的相对表现曲线更为平滑，波动率更低。这对于寻求中期主题敞口的投资者而言，可能是一个更优的选择。

*但这里有一个隐含的前提：** 赞助商篮子的超额收益，是否依赖于赛事本身的成功和全球关注度？如果2026年世界杯因为地缘政治或移民政策变化而出现观众人数不及预期，赞助商的品牌曝光效果就会打折。报告明确将“美国移民政策变化”和“地缘政治逆风”列为关键风险，这一点需要投资者保持警惕。

\subsection{[R015] GS：欧洲石油巨头隐含的长期油价远低于市场共识，这才是真正的价值锚点}
{\small\color{refgray}归类：能源与大宗：欧洲石油巨头隐含油价远低于共识，厄尔尼诺精准打击作物窗口}

GS：欧洲石油巨头隐含的长期油价远低于市场共识，这才是真正的价值锚点

当市场目光聚焦于美伊和谈可能带来的霍尔木兹海峡重新开放，并因此推动布伦特油价下跌时，一个更深层的定价问题被忽略了：欧洲上市石油巨头（EU Big Oils）的当前股价，究竟在反映什么样的长期油价预期？GS最新发布的研报给出了一个令人意外的答案——这些公司的估值所隐含的长期油价，远低于当前市场共识和远期期货所暗示的水平。这并非一个简单的“便宜还是贵”的判断，而是指向了一个关于资产定价、股东回报和地缘政治风险重新定价的核心分歧。

这份报告的核心判断是：欧洲石油巨头的估值已经消化了显著低于市场预期的长期油价，这为投资者提供了一个不对称的风险收益窗口。但前提是，投资者需要理解这种“隐含定价”背后的结构性假设，并识别出哪些公司能够在低油价环境中依然保持强劲的现金流和股东回报能力。

GS采用了一个简洁但有力的框架：求解欧洲石油巨头在实现8\%的自由现金流收益率（FCF yield）时，所对应的 年平均布伦特油价。结果令人瞩目——这一隐含价格仅为每桶60美元。这与此期间彭博共识预期的每桶78美元、GS自己的预测每桶77美元，以及远期期货价格每桶79美元相比，存在约20美元的显著折价。

换句话说，市场对欧洲石油公司的定价，并非建立在对当前高油价持续性的乐观假设上，而是已经提前、且相当充分地计入了油价回归到接近页岩油边际成本的水平。这意味着，如果未来油价能够维持在每桶70美元甚至更高的水平，这些公司的实际现金流和股东回报将显著高于当前市场所预期的。

这个发现的核心含义在于：投资者对欧洲石油股的“价值陷阱”担忧，可能恰恰是误解了市场已经完成的风险定价。市场不是在为“高油价红利”买单，而是在为“长期低油价环境下的生存和回报能力”定价。任何超出每桶60美元的油价情景，都可能带来超预期的回报。

GS的压力测试揭示了另一个关键洞察：欧洲石油巨头的自由现金流收益率对油价的敏感性，在油价下行区间表现出不对称的韧性。在每桶90美元、80美元、70美元、60美元和50美元的六大油价情景下，2026年整个板块的自由现金流收益率依次为13.4\%、11.6\%、9.9\%、8.1\%和6.4\%。

即使油价跌至每桶50美元，板块整体的自由现金流收益率仍能维持在6.4\%的水平。这背后的原因是，这些公司经过多年的成本优化和资本纪律，其盈亏平衡点已经大幅下降。更重要的是，在每桶60美元的隐含定价水平上，8.1\%的自由现金流收益率已经为股东回报提供了坚实的底部支撑。这意味着，即使油价在未来几年内持续承压，这些公司依然有能力维持甚至提高分红和回购，除非油价跌至每桶50美元以下。

这一分析对资产定价的启示是：欧洲石油股的估值底线比市场普遍认为的要坚实得多。投资者不应仅关注油价绝对水平的变化，而应关注公司现金流对油价下行的“免疫能力”。那些成本结构更优、资产组合更多元的公司，其估值下行风险实际上已经被充分定价。

与自由现金流收益率相对应，GS也测算了在不同油价情景下，欧洲石油巨头的股东总回报收益率（股息加回购）。在每桶90美元时，总回报收益率可达8.0\%；即使在每桶60美元时，也有7.0\%。值得注意的是，从每桶90美元降至60美元，总回报收益率仅下降了1个百分点，表现出极高的稳定性。

这种稳定性主要来自两个结构性因素：一是这些公司普遍承诺了基于绝对金额或固定比例的分红和回购计划，而非与油价挂钩的浮动政策；二是许多公司正在利用当前的高油价窗口去杠杆，从而在油价下行时拥有更大的财务灵活性来维持股东回报。例如，BP在此次研报中被特别提及，因其拥有强劲的财务去杠杆能力和来自交易利润的潜在上行空间。

对于高净值投资者和产业决策者而言，这意味着欧洲石油股正在演变为一种“类债券”的收益型资产，其现金

[... middle omitted ...]

绪正在压低布伦特油价，但GS的隐含定价分析是基于当前股价倒推的，而非对未来地缘政治事件的预测。如果霍尔木兹海峡真的重新开放，全球石油供给的增加可能会在短期内将油价推向甚至低于每桶60美元的水平。然而，报告的分析表明，即使在这一极端情景下，欧洲石油公司的自由现金流收益率仍能维持在6.4\%左右，这远高于许多其他行业的平均水平。

但这里仍有一个未解的问题：市场对地缘政治风险的定价是一次性完成，还是需要经历一个反复确认的过程？如果油价在未来几个月内快速跌破每桶60美元，这些公司的估值可能会面临短期压力，但其长期股东回报的吸引力反而会进一步增强。这个矛盾本身就构成了一个值得深入探讨的交易主题。

基于GS的研究，投资者应该建立一个新的观察框架，而不是继续纠缠于对油价的短期预测。这个框架包含三个层次：

第一，关注公司的盈亏平衡点。那些拥有更低成本、更长寿命资产的油田项目，以及更灵活的资本开支计划的公司，将拥有更强的下行保护。GS特别提到了Galp，因其拥有最长的石油储量寿命。

\subsection{[R016] GS：日本市场真正值得关注的不是AI反弹，而是股东会投票率背后的治理信号}
{\small\color{refgray}归类：区域市场：韩国居民资产结构性迁移，日本治理改善信号被低估}

GS：日本市场真正值得关注的不是AI反弹，而是股东会投票率背后的治理信号

每年六月的最后一周，日本迎来股东大会的绝对高峰期。根据GS最新发布的日本市场周报，6月25日和26日两天，将有超过1100家3月决算期的公司集中召开年度股东大会。这一现象本身并不新鲜——日本企业长期习惯于在同一时间窗口完成股东大会，以便于统一应对机构投资者的集中质询。

但这份报告真正值得关注的判断，并不是会议日程本身，而是一个被当前市场情绪掩盖的结构性信号：那些在去年股东大会上获得低赞成票的公司，正在系统性地改善治理表现，而这种改善正在转化为可量化的超额收益。

GS筛选出的12家2025年AGM赞成票率处于最低十分位的公司中，已有10家在今年的投票中获得了明显的赞成率提升。更关键的是，这一组合在去年股东大会到今年伊朗冲突爆发前，跑赢TOPIX指数13个百分点。尽管此后因市场风格转向AI主题而暂时落后，但GS认为，低赞成票率仍然是预测企业未来股东友好行为的最强先行指标。

这一判断的隐含含义是：当前市场对日本股市的定价，可能过度集中于AI和科技主题的短期叙事，而低估了治理改善这一慢变量对资产定价的长期重塑力。

1. 低AGM赞成票率是一个被低估的“治理改善催化剂”，而非简单的负面信号

许多投资者的直觉是，一家公司在股东大会上获得低赞成票，说明管理层与股东之间存在矛盾，是负面信号。但GS的研究提供了相反的实证：低赞成票率恰恰是推动企业在下一年度采取股东友好行动的最强动力。

逻辑链条并不复杂。日本交易所集团（JPX）和机构投资者对治理透明度的要求逐年提高。当一家公司的管理层在股东大会上获得低于80\%甚至70\%的赞成票时，这不仅是面子问题，更会触发机构投资者的持续关注和跟进。管理层为了在来年避免再次陷入低票困境，往往会在分红、回购、业务重组、董事会独立性等方面做出实质性让步。

GS在2月发布的治理策略报告中，系统筛选出了2025年AGM赞成票率处于最低十分位的公司组合。截至当前，该组合中已公布今年投票结果的10家公司全部实现了赞成率提升。以AEON为例，其赞成率从77.8\%跃升至94.0\%，提升了16.3个百分点。Coca-Cola Bottlers Japan和SHIMAMURA的升幅也分别达到14.0和13.0个百分点。

这意味着，低赞成票率不是终点，而是一个治理改善周期的起点。 对于投资者而言，这提供了一个低成本的筛选框架：在每年股东大会季结束后，识别那些赞成率偏低但基本面尚可的公司，持有至下一个股东大会季，本质上是在博弈治理改善的确定性溢价。

2. 市场修正中的“防御性因子”正在被重新定价，优待+分红组合跑赢大盘12个百分点

GS在6月3日发布了一份针对“股东优待（yuutai）”策略的专题报告，探讨了在重大市场修正期间，拥有股东优待制度的股票是否具备防御属性。仅一周之后，市场就给出了一个近乎完美的实证窗口。

自6月3日起，日本股市经历了一轮由外部地缘风险引发的修正。在此期间，GS的“优待在+分红”组合（Y+D Screen）跑赢了“无优待、无分红”组合（NYND Screen）12个百分点，同时跑赢TOPIX指数4个百分点。如果聚焦于中小盘股票，Y+D组合相对于NYND组合的超额收益也达到7个百分点。

这一数据的含义远超短期交易策略的范畴。它揭示了一个更深层的市场行为规律：在不确定性上升的环境中，投资者愿意为那些“能感受到的股东回报”支付溢价。 股东优待——无论是商品券、折扣券还是积分——在日本零售投资者中具有极强的心理锚定效应。当市场恐慌时，这些“看得见摸得着”的回报机制，比财务报表上的股息率更能稳定持有人结构。

从行业表现来看，这一逻辑也得到了验证。在修正期间，食品饮料和零售板块分别录得3\%的周度正收益，而

[... middle omitted ...]

出明显的“追涨AI、避险卖出”特征。而内资——尤其是个人投资者——则在市场修正中持续加仓，且其偏好明显偏向于国内需求、高分红和股东优待类股票。

*这种仓位结构的背离，往往意味着下一轮风格切换的前奏。** 当外资主导的AI和科技主题出现获利回吐或基本面边际恶化时，内资重仓的治理改善和防御型资产可能获得相对收益。而从GS的筛选结果来看，低AGM赞成票组合在AI主题主导的市场中表现落后，恰恰说明其估值尚未被充分定价。

4. 报告尚未完全回答的关键问题：治理改善的“边际递减效应”何时出现

GS的报告提供了清晰的实证框架，但一个关键问题尚未得到充分讨论：治理改善带来的超额收益，是否存在边际递减的拐点？

从逻辑上看，当一家公司的赞成率从65\%提升到80\%，市场会给予显著的正面反馈。但当它从80\%提升到90\%时，这一改善的边际信息价值就会下降。GS的筛选组合中，部分公司今年的赞成率已经达到90\%以上。对于这些公司，治理改善的驱动力是否已经耗尽？下一阶段的超额收益来源是什么？

\subsection{[R017] Bernstein：邮轮市场的真正稀缺资产，不是船，是品牌定位的不可替代性}
{\small\color{refgray}归类：医疗政策与邮轮：2026中期选举是医疗政策拐点，奢华邮轮品牌定位是真正护城河}

Bernstein：邮轮市场的真正稀缺资产，不是船，是品牌定位的不可替代性

当大多数投资者还在用“运力增长”和“客单价”来理解邮轮行业时，一份来自Bernstein的深度报告提出了一个更值得追问的问题：在一个供给和需求都在结构性扩张的市场里，什么才是真正不可复制的护城河？

答案不是船队规模，不是航线数量，甚至不是豪华程度。答案是一种近乎反直觉的品牌定位——在奢华邮轮领域，真正跑赢的不是最贵的那个，而是最“不一样”的那个。

这份研报聚焦于一个长期被主流投资者忽视的细分赛道：奢华邮轮。Bernstein的判断很明确：奢华邮轮正站在邮轮行业最强的结构性需求风口上，而在这个细分市场中，只有一家公司具备“纯正标的”属性——Viking。但真正值得深思的，不是Viking的估值，而是它凭什么能在定价上跑赢船型规模相似的竞争对手，以及这种定价权的可持续性。

这不是一篇简单的“推荐买入”报告。它实际上在回答一个更深层的问题：当消费升级的故事不再新鲜，当“有钱人变多”变成共识，企业之间真正的差距，来自对“什么是奢华”的重新定义。

Bernstein在报告中明确指出，奢华邮轮面临的两个核心需求驱动因素，在可预见的未来都不会逆转。

第一个因素是人口结构。邮轮消费本身就有明显的年龄倾向，而奢华邮轮尤甚——平均乘客年龄达到63岁。这一群体在美国人口中的占比正在持续扩大，更重要的是，他们的财富集中度也在上升。目前，美国55岁以上人群持有近75\%的家庭财富。与此同时，美国收入分布的两极分化进一步加剧，收入最高四分位群体的收入增长最快，休闲时间也最多。这个群体是奢华旅游消费的核心客群。

第二个因素来自供给端。奢华酒店的新增供应正在放缓。建设成本、劳动力成本和融资成本的上升，已经拖慢了高端酒店的开发速度。这意味着，一部分原本流向奢华酒店的需求，将溢出到奢华邮轮市场。

这两个因素叠加，构成了一组罕见的“需求加速、供给受限”的定价环境。但问题在于，这种宏观利好是否足以支撑所有奢华邮轮品牌的增长？Bernstein的回答是：未必。宏观只是背景，真正的分化来自品牌能否抓住这个机会窗口。

与主流邮轮市场不同，奢华邮轮内部存在显著的差异化。Bernstein的报告将奢华邮轮分为三个层级：中型船奢华（约1000客位）、小型船奢华（500-700客位）和游艇级（100-400客位）。每晚价格从300-700美元一路攀升至5000美元以上，甚至高达8000美元。

关键变量是船型大小。消费者愿意为更小的船支付指数级更高的价格。乘坐阿曼集团110客位游艇的每晚价格，是Seabourn一艘五倍大小船只的10倍。这背后的逻辑不难理解：更小的船意味着更高的私密性、更多的个人空间和更高的服务密度。

但这份报告真正有意思的发现是：Viking打破了这条定价曲线。在约1000客位的船型中，Viking的定价远高于同级别的Oceania，甚至接近小型船奢华品牌。消费者愿意为Viking支付约两倍于船型规模所“应得”的价格。这意味着，Viking的定价权并非来自硬件规模，而是来自品牌本身。

Bernstein用三个维度解释了Viking如何实现这种定价溢价。

第一，Viking卖的“不是船，是目的地”。其他奢华邮轮品牌强调船上体验——餐饮、服务、全包式享受。Viking则聚焦于“有思想的人”，通过船上历史学家、文化讲座和精心策划的岸上游览，让乘客在抵达目的地之前就已经开始“学习”。这种定位，本质上是在重新定义奢华邮轮的价值主张：不是被服务，而是被启发。

第二，Viking提供的是“放松的奢华”，而非“炫耀的奢华”。它不提供香槟和鱼子酱，不配备管家和白手套，不设赌场，不拍卖艺术品。它的设计语言是斯堪的纳维亚式的简约和克制。这种风格在吸引“非传统邮轮乘客”方

[... middle omitted ...]

测算显示，到2030年，Viking将贡献奢华邮轮市场超过50\%的新增运力。但报告同时指出，Viking的船型设计使其能够以极低的每客位成本增加运力——每客位成本约为50万美元，与主流邮轮品牌相当，但收益显著更高。

这意味着，Viking在扩张时，不会像其他奢华品牌那样面临“规模稀释品质”的困境。它的船型标准化程度高，设计成熟，供应链稳定，因此边际扩张成本可控。而其他品牌，尤其是游艇级品牌，每客位成本可能高出数倍，且船型非标准化，扩张节奏受制于定制化造船周期。

这一成本优势，是Viking在供给端建立壁垒的关键。它不仅能跟上需求增长，还能在增长中保持利润率的稳定性。

5. 报告尚未完全回答的关键问题：游艇级品牌的竞争格局会如何演变

Bernstein的报告对游艇级品牌的定位进行了区分：它们与传统奢华邮轮不构成直接竞争。丽思卡尔顿、四季、东方快车、阿曼等品牌的游艇，更多是在延伸其已有的酒店品牌资产，吸引的是对“邮轮”概念有抵触但对“游艇”概念有兴趣的高净值客户。

\subsection{[R018] Bernstein：AI在放射科的落地，真正被低估的不是技术，而是结构性供需缺口}
{\small\color{refgray}归类：医疗与教育：AI放射科从可选项变必需品，教育市场结构分化加速}

Bernstein：AI在放射科的落地，真正被低估的不是技术，而是结构性供需缺口

医疗AI的故事已经讲了十年。但Bernstein这份最新研报揭示了一个被市场忽视的关键信号：放射科AI的渗透正在从一个“可选项”转变为一个“必需品”。这不是因为技术突然成熟了，而是因为人口结构和劳动力市场的双重挤压，已经让医疗系统别无选择。

这份报告最值得关注的判断是：AI在医学影像领域的价值，不是来自于它能做什么，而是来自于它不得不做什么。老龄化叠加放射科医生供给缺口，正在制造一个结构性需求缺口，而AI是目前唯一能够在不显著增加人力成本的前提下填补这个缺口的方案。

Bernstein的分析师团队指出，美国65岁以上人口将以每年约3\%的速度增长至2030年，而放射科医生的供给缺口将在2034年达到10\%。这两条线的交叉点，就是AI医学影像从“锦上添花”到“雪中送炭”的转折点。

COVID-19疫情是一个关键的催化剂，但它不是问题的根源。疫情只是将原本缓慢积累的压力一次性释放了出来。英国NHS的数据提供了一个清晰的参照：等待超过六周进行CT或MRI检查的患者数量，从疫情前的约8000人飙升至2020年的约80000人，而到了2025年，这个数字仍然维持在约75000人的高位。

这意味着什么？这意味着疫情后的“恢复正常”并没有真正恢复。医疗系统在应对积压需求时，发现常规的扩招医生、增加设备投入已经无法跟上需求的增长速度。这不是一个短期波动，而是一个结构性失衡的持续表现。

Bernstein的报告引用McKinsey的数据指出，AI能够释放医疗设备准备人员高达48\%的工作时间。这个数字放在放射科的具体场景中，意味着同样的设备和人员，可以在AI辅助下处理更多的扫描量。不是替代医生，而是让医生的时间从重复性劳动中释放出来，集中在诊断和决策上。

对于投资者和产业决策者而言，这里的关键洞察是：放射科AI的需求不是由技术供给驱动的，而是由系统性的效率缺口拉动的。只要这个缺口存在，AI工具的采购就不会是“预算充裕时的锦上添花”，而是“不得不做的效率投资”。

市场对AI医学影像的认知，往往停留在“检测病灶的准确率”上。但Bernstein的研报揭示了一个更重要的维度：AI的真正价值在于它如何嵌入放射科的现有工作流。

报告详细分析了几个领先的创业公司，包括Aidoc、Viz.ai和Rad AI。这些公司的共同点不是它们算法的准确率有多高，而是它们如何重构放射科的工作流程。

Aidoc的aiOS平台不仅仅是一个图像分析工具，它是一个“始终在线”的图像分析引擎，能够自动对CT和X光片进行优先级排序，将急性病例提前到放射科医生的阅片列表顶端。Viz.ai则更进一步，它的平台不仅检测异常，还能自动通知相关的临床团队，实现从诊断到治疗的快速衔接。Rad AI则专注于报告生成环节，利用生成式AI自动起草影像报告的结论部分。

这些案例揭示了一个核心洞察：在放射科这个高度结构化的环境中，AI的价值不在于单点突破，而在于全流程的效率提升。一个能够自动识别肺结节的算法固然有用，但如果它不能与PACS系统无缝对接、不能自动调整医生的工作列表、不能生成结构化的报告，它的实际价值就会大打折扣。

对于产业竞争格局而言，这意味着拥有“平台能力”的公司将比单纯卖算法的公司拥有更强的护城河。谁能够成为放射科工作流的“操作系统”，谁就能在长期竞争中占据优势地位。

Bernstein的研报列举了几个值得关注的估值案例。Aidoc在2026年4月的E轮融资中估值达到5.2亿美元，Viz.ai在2022年4月的D轮融资中估值达到12亿美元，Rad AI在2025年5月的C轮融资中估值达到5.25亿美元。

这些估值数字本身并不惊人，但值得注意的是它们背后的逻辑变化。早

[... middle omitted ...]

指标不再是论文发表数量或公开竞赛排名，而是客户留存率、每客户收入贡献、以及与现有PACS和HIS系统的整合深度。

尽管Bernstein的研报提供了丰富的信息，但有一个关键问题它没有完全展开：AI医学影像软件的定价权到底在哪里？

目前的商业模式主要有两种：一种是按扫描次数收费，另一种是年度订阅费。但这两种模式都面临一个共同的挑战——医院的付费意愿。在大多数医疗系统中，影像科本身是一个成本中心而非利润中心。医院愿意为新的MRI设备支付数百万美元，因为设备直接决定了能否提供服务。但AI软件是一种“效率提升工具”，它的价值难以量化，因此也容易在预算紧张时被削减。

这里存在一个尚未被市场充分定价的变量：如果AI工具能够显著减少放射科医生的雇佣需求，那么它的价值就不是“让医生更快”，而是“让医院少招人”。在放射科医生供给缺口持续扩大的背景下，AI工具的定价逻辑可能会从“成本节约”转向“产能替代”。如果这个判断成立，那么AI医学影像的市场规模可能会比当前市场预期的要大得多。

\subsection{[R019] 摩根斯坦利：市场低估的不是消费复苏，而是亚洲工业超级周期对中国出口的锚定效应}
{\small\color{refgray}归类：未被主题摘要引用}

摩根斯坦利：市场低估的不是消费复苏，而是亚洲工业超级周期对中国出口的锚定效应

这份来自摩根斯坦利亚洲团队的投资者演示，标题直接点明了当前中国经济最关键的结构性变量——AI与能源超级周期。但真正值得反复咀嚼的判断，并不在标题本身，而在于报告揭示的一个核心张力：中国经济正在经历一场“双速”分化，而出口驱动的增长韧性，远比市场普遍认知的更持久、更结构化。

市场目前对 年中国经济的讨论，仍然高度集中在房地产何时触底、消费何时回暖、通缩压力何时缓解。这些固然是重要变量，但摩根斯坦利的分析框架提示我们：真正被低估的，是亚洲工业超级周期对中国出口竞争力的重塑。这不是一个短期库存回补的故事，而是一个可能持续到2030年的结构性再定价过程。

报告的核心主张可以概括为一句话：中国正在借助AI与能源相关的资本开支浪潮，以及自身在全球制造业中已经占据的压倒性份额，将出口市场占有率从当前的14.8\%推升至2030年的16.5\%，这一过程将有效对冲内需的疲弱，使GDP增速在4.7\%-5.0\%的区间内保持稳定。但这一判断成立的前提，以及它背后隐藏的风险，需要逐层拆解。

1. 亚洲工业超级周期并非AI主题的衍生品，而是一次广泛且独立的产能扩张

摩根斯坦利在报告中明确指出，亚洲正在经历自2000年代中期以来最强的工业周期。这一判断的关键证据，并非仅来自半导体或AI相关产品的出口数据。报告展示了一张极具说服力的图表：自2025年10月以来，亚洲剔除半导体和燃料的非科技出口出现了急剧且广泛的增长，涵盖中间品、资本品、消费品乃至乘用车等几乎所有产品类别。

这意味着什么？这轮工业周期的驱动力不是单一的AI军备竞赛，而是亚洲制造业整体竞争力的提升和全球供应链重构的深化。中国作为亚洲制造业的核心节点，其出口增长并非依赖某一两个爆款产品，而是建立在全产业链的规模优势和响应速度之上。

对于投资者而言，这一判断的含义是深远的。如果仅仅是AI主题驱动，那么出口增长的可持续性将高度依赖于科技巨头的资本开支节奏，波动性极大。但如果是广泛的工业周期，则意味着中国出口的韧性来自更基础的结构性因素——全球企业对亚洲供应链的依赖度在上升，而中国在其中占据了不可替代的位置。

2. 中国制造业的全球份额优势，正在从“数量”转向“价值”层面的锁定

报告提供了一组令人印象深刻的数据：中国在几乎所有制造业细分领域的全球增加值占比，从2017年到2024年都在提升。非金属矿物制品（水泥、玻璃等）从43\%升至48\%，皮革及鞋类从44\%升至47\%，其他杂项制造从17\%跃升至35\%，基本金属和金属制品从29\%升至35\%，纺织和纺织品从33\%升至40\%。

这些数字合在一起意味着什么？中国制造业的全球主导地位，已经不是简单的“世界工厂”概念所能概括。在大多数领域，中国的份额已经超过了全球其他主要生产国的总和。这种压倒性的规模优势，带来了两个关键的经济含义：第一，任何试图将供应链迁出中国的努力，都将面临极高的成本和效率损失；第二，中国企业在定价权和产业链整合方面，拥有了前所未有的议价能力。

摩根斯坦利预计，中国全球出口市场份额将从2024年的14.8\%升至2030年的16.5\%，在乐观情景下甚至可能达到18.0\%。这个预测的背后逻辑，正是基于中国制造业已经锁定的结构性优势，而非短期周期性的因素。

报告最深刻的洞察之一，在于揭示了出口增长与内需复苏之间的脱节。摩根斯坦利明确指出，“Positive spillover to the broader economy could be milder this time”。原因有两个，且都值得仔细推敲。

第一，就业弹性在下降。报告展示的散点图显示，制造业企业的营收增长与就业增长之间的关联度，在过去十年间显著减弱。这意味着，即使出口订

[... middle omitted ...]

经大不如前。就业市场难以充分受益于出口繁荣，居民收入增长受限，消费自然难以强劲复苏。这正是中国经济“双速”分化的根源——出口强劲，但内需滞后。

报告在讨论内需滞后时，特别强调了劳动力市场的疲软是关键的制约因素。制造业和非制造业的就业景气指数均处于低位，青年失业问题更为突出。但真正值得警惕的，是报告提出的一个前瞻性判断：AI的扩散可能加剧劳动力市场的压力。

摩根斯坦利将AI的影响分为三个阶段：生成式AI、代理式AI、物理AI。每一阶段对不同群体的冲击各不相同。生成式AI可能侵蚀中产阶级的工资稳定性，代理式AI将加剧结构性失业风险，而物理AI（如人形机器人）则可能对蓝领服务业岗位造成更大冲击。

这一判断的So what含义是：即使出口持续增长，如果AI对就业的替代效应超过了出口对就业的拉动效应，那么内需的复苏将更加遥不可及。这不仅是经济问题，也是社会问题。报告虽然没有展开讨论政策应对，但隐含的逻辑是清晰的——政府可能需要更大力度的社会保障和再分配政策来对冲这一趋势。

\subsection{[R020] JPM：中国电力设备在美渗透的拐点比市场预期的更快到来}
{\small\color{refgray}归类：电动车与供应链：中国渗透率结构性跃迁，电力设备回调是重新定价}

这份JPM本周发布的研报，来自分析师团队在上海参加数据中心展和智慧能源展后的现场调研笔记。其核心判断不是关于中国电力设备市场的整体景气度——这已经被充分定价——而是关于一个正在发生但尚未被市场完全消化的结构性变化：中国电气设备供应商正在以超出预期的速度渗透美国数据中心供应链。

报告透露的信号是具体的、可验证的：部分中国企业的变压器交付周期是西方同行的四分之一到六分之一，美国数据中心客户已经开始实地考察中国工厂并下达订单，而更关键的是，这一轮渗透并非简单的价格战，而是建立在交期、定制化能力和逐步积累的认证基础之上。对于关注电力设备板块的投资者而言，真正需要回答的问题不是“中国设备能否进入美国”，而是“哪些企业能在这轮渗透中建立可持续的竞争优势，而不是仅仅享受一轮周期红利”。

1. 美国电气设备供应链的瓶颈正在为中国供应商打开一扇以前紧闭的门

研报中反复出现的核心事实是：美国电力设备市场的供需缺口正在为中国企业创造历史性的窗口期。一家输配电设备供应商的变压器交付周期为6至12个月，而西方同行需要2至3年。这不是微小的差距，这是决定数据中心能否按时上线的关键变量。

更值得关注的是，这并非个例。JPM的渠道调研显示，中国供应商在高压/低压设备领域的渗透率正在上升，驱动因素正是供应紧张和交期优势。数据中心的电力设备需求具有高度的时间敏感性——项目延迟意味着算力无法按时交付，这对云服务提供商（CSP）和大型互联网公司而言是巨大的机会成本。因此，当西方供应商的交付周期拉长至两年以上时，中国供应商的6个月交期就变得具有压倒性的吸引力。

但交期优势只是入场券。真正让这轮渗透不同于以往的是：美国客户正在主动评估中国供应商。报告明确提到，一些CSP今年早些时候已经访问了中国工厂并下达了订单。这意味着市场准入正在从被动接单转向主动认可。对于已经在美国市场有所布局的企业，如伊戈尔电气、安科瑞智能、上海电气和山东泰开，这轮窗口期的价值可能远超当前股价的反映。

如果只是把中国市场的低价产品卖到美国，这轮渗透的可持续性将大打折扣。但报告揭示了一个更有利的经济学结构：美国市场的定价能力远强于中国。

一家供应商提到的数据是：美国市场的变压器毛利率可能达到30\%至40\%以上，而中国市场的毛利率仅在十几个百分点。同时，美国市场的平均售价（ASP）是中国市场的数倍，即使中国供应商的报价低于西方同行，其绝对利润水平仍然远高于国内。这种利润率差异意味着中国供应商有足够的空间进行产品升级、认证投入和本地化服务建设，从而形成正向循环。

但需要强调的是，这种利润率优势并非无条件的。报告指出，美国公用事业领域的渗透仍然有限，因为该领域对业绩记录、认证和采购名单准入有更严格的要求。目前中国供应商的主要切入点仍是数据中心，通过本地销售团队、EPC承包商和分销商渠道进入。这意味着，能够建立完整的美国本地化销售和服务体系的企业，将比单纯依赖出口的企业享有更持久的竞争优势。

3. 固态变压器（SST）正在成为下一代数据中心电力架构的关键变量，但商业化节奏仍有分歧

研报中关于固态变压器的内容值得单独拆解。伊顿在中国发布了面向AI数据中心的中压SST 2.0，采用10kV中压交流输入至800V直流输出的架构，单系统容量2.5MW，端到端效率98.5\%，可用性超过99.999\%。更重要的是，伊顿表示在中国有一个已运行超过一年的参考项目，这为产品验证提供了支撑。

但真正引发行业讨论的是特锐德与伊顿的战略合作。双方计划联合开发全预制数据中心电力解决方案，将伊顿的SST和中压直流能力与特锐德的高压预制舱系统结合，目标是为下一代AI数据中心构建220kV至800V直流的全链路架构。特锐德同时发布了“AI PowerHouse”系统级解决方案，宣称150天全站交付、2

[... middle omitted ...]

累的企业将占据先发优势。

4. 2027年美国数据中心需求将更加强劲，但供应链瓶颈可能成为制约因素

报告中最具前瞻性的判断来自一家企业提供的信息：其今年迄今为止已获得约10亿元人民币的美国数据中心相关输配电订单，但更大的AI数据中心采购周期尚未完全启动，许多项目仍处于管道中，受限于电力可用性、电网连接和长设备交期。

这意味着当前看到的订单可能只是冰山一角。一些示范项目已经在提前订购变压器，因为交期超过一年。展望未来，更多的数据中心并网项目和表后发电项目将推动更强劲的电气设备需求。但这里有一个隐含的风险：如果中国供应商的渗透速度跟不上需求爆发，美国市场可能面临更严重的设备短缺，从而推高价格并延长项目周期。

另一个值得注意的细节是，数据中心客户在关键组件上仍然倾向于选择全球领先的电气品牌，如施耐德、西门子和ABB，尤其是在高规格项目中。中国供应商的机会更多存在于成本敏感型项目或客户允许本地替代的场景。这意味着，中国企业的渗透将是一个渐进的过程，而非一夜之间的替代。

\subsection{[R021] UBS：人形机器人供应链的“确定性”正在从概念走向订单，但市场仍低估了传统业务的反转弹性}
{\small\color{refgray}归类：汽车半导体与人形机器人：复苏已至，但结构性分化是关键}

UBS：人形机器人供应链的“确定性”正在从概念走向订单，但市场仍低估了传统业务的反转弹性

2026年6月10日，UBS在嘉兴举办了一场面向机构投资者的实地调研，走访了双环传动和荣泰电工两家公司，并与核心管理层进行了深度交流。这份调研纪要看似聚焦于两家零部件企业的季度经营更新，但其背后揭示了一个更深层的产业信号：中国汽车零部件供应链正在经历一轮“需求结构切换”的关键拐点。

市场目前的关注点高度集中在人形机器人这一新兴赛道上。这当然重要。但UBS这份报告最有价值的判断，并非机器人业务本身，而是“传统NEV业务正在环比改善，而欧洲市场渗透率提升正在为这些企业提供新的增长动量”。换句话说，人形机器人的故事是估值弹性，而传统业务的复苏才是安全边际的基石。

这份报告真正值得细读的地方在于：它提供了三家核心标的（科达利、双环传动、荣泰电工）在“新旧动能切换”这个时点上的微观证据。这不是一份宏观展望，而是一次带着订单和产能数据回来的实地验证。

1. 双环传动的NEV齿轮产能利用率突破90\%，这是比机器人更重要的信号

双环传动在交流中披露了一个关键数字：Q2 2026 NEV齿轮出货量环比显著增长，产能利用率已超过90\%。在汽车零部件行业，产能利用率一旦突破85\%-90\%的临界点，就意味着固定成本摊薄进入加速区间，毛利率将出现非线性改善。

这个数字的意义在于：它验证了UBS此前对国内NEV需求在二季度企稳回升的判断。过去几个季度，市场对NEV零部件企业的担忧主要集中在“价格战传导”和“产能过剩”上。双环传动90\%以上的产能利用率，至少说明在齿轮这个细分领域，头部企业的产能消化能力远好于市场悲观假设。

更值得关注的是，双环传动正在将齿轮技术外溢到AIDC发电机领域，核心客户包括卡特彼勒和玉柴。管理层指引该业务将在2026年开始贡献收入。这意味着双环传动正在从单一汽车零部件供应商，向“汽车+数据中心+机器人”的多元驱动结构演变。这种技术外溢能力，是传统零部件企业估值体系重构的关键变量。

UBS报告没有完全展开的是：AIDC发电机齿轮的技术壁垒、毛利率水平、以及卡特彼勒的订单规模，目前仍缺乏公开数据支撑。这部分业务的能见度，将是判断双环传动长期估值中枢能否上移的核心。

2. 荣泰电工的精密丝杠已获PPAP批准，人形机器人从“样品”进入“订单”阶段

荣泰电工CEO在交流中给出的信息，可能是这份报告中最具“从0到1”意义的数据点：公司用于灵巧手和本体的精密丝杠，已获得某北美头部人形机器人制造商的PPAP批准，并在6月观察到订单开始爬坡。同时，泰国制造基地正在加速建设，为H2 2026的量产做准备。

PPAP（Production Part Approval Process）是汽车和高端制造业中极为严格的供应商准入流程。获得PPAP批准，意味着荣泰电工的丝杠产品在尺寸、材料、工艺、耐久性等方面都通过了客户的全套验证。这不再是“实验室样品”或“小批量试制”，而是拿到了进入量产供应链的正式门票。

UBS报告同时指出，荣泰电工正在向该北美客户送样更多产品，包括结构件。管理层对电机和减速器的进展也表示“按计划推进”。这意味着荣泰电工在机器人关节总成中的“零件渗透率”正在扩大，从单一丝杠向更多环节延伸。

但这里有一个需要持续验证的关键问题：订单爬坡的速度和规模。PPAP批准是必要条件，但不是充分条件。量产阶段的良率、成本控制、产能爬坡速度，才是决定荣泰电工能否真正从机器人业务中兑现利润的核心变量。UBS报告对此着墨不多，这部分信息只能通过后续季报和客户订单公告来持续跟踪。

3. 欧洲NEV渗透率提升，正在成为传统业务的“第二增长曲线”

UBS报告提到，欧洲市场NEV渗透率上升，有望为传统业务带来新的增长动能。荣泰电工管理

[... middle omitted ...]

意味着中国零部件企业的“出海”逻辑，正在从“规避关税”升级为“贴近客户、参与当地供应链建设”。这种从“出口”到“本地化”的转变，将显著改变海外业务的盈利结构和估值逻辑。

UBS报告没有完全展开的是：欧洲NEV渗透率提升的可持续性，以及中国零部件企业在欧洲的产能扩张是否会引发新的竞争格局变化。这些问题的答案，将决定这一“第二增长曲线”的天花板在哪里。

4. 报告尚未完全回答的关键问题：机器人业务的利润兑现路径与竞争格局演变

尽管UBS报告提供了大量正面信号，但有两个关键问题目前仍处于“待验证”状态。

第一，机器人业务的利润兑现路径。无论是双环传动的定制化减速器，还是荣泰电工的精密丝杠，目前都处于“小批量爬坡”或“PPAP刚获批”的阶段。从PPAP批准到大规模量产，再到贡献正现金流，中间存在多个不确定性节点：良率、成本、产能建设进度、客户订单稳定性。UBS报告对这些环节的量化预测非常有限，这意味着市场对人形机器人业务的估值，目前更多是基于“可能性”而非“确定性”。

\subsection{[R022] GS：中国教育市场最值得跟踪的不是增长，而是结构分化的加速}
{\small\color{refgray}归类：医疗与教育：AI放射科从可选项变必需品，教育市场结构分化加速}

当多数人还在关注教育行业“复苏”或“监管松动”时，一份来自GS的5月跟踪报告揭示了一个更值得深究的信号：市场的真实故事不是整体回暖，而是三条独立逻辑线正在同时加速——AI原生产品的用户争夺、硬件定价权的重新分配、以及海外需求的触底节奏。这三条线并不平行，它们的交汇点将决定未来12个月哪些公司能获得估值重评。

这份5月26日更新的教育行业追踪报告，覆盖了AI教育应用、AI学习平板、移动端用户活跃度、非学科培训牌照、留学趋势以及东方甄选等多个维度。数据密度极高，但去掉噪音后，真正值得决策者关注的只有三个问题。

豆包爱学的DAU同比增速在5月底加速至146\%，这已经是它连续第二个月加速增长（4月为60\%）。更重要的是，它在中国AI原生应用DAU排名中稳居第六位。考虑到它是一款垂直教育产品，这个排名的含金量远超表面数字。

横向对比更说明问题。Gauth（好未来旗下出海AI产品）5月月流水约110万美元，同比增长24\%，近12个月累计流水已达1320万美元。Gauth在全球AI教育应用中MAU排名第一，但它的增长斜率正在放缓——4月同比增速还有35\%，5月降至24\%。豆包爱学则在加速追赶。

这两款产品的反差揭示了一个关键判断：AI教育应用的竞争已经从“谁先推出产品”进入“谁能持续获取用户注意力”的阶段。豆包爱学的DAU增长加速，叠加它在K-12学习工具类应用中的使用时长份额持续扩大（报告中有详细图表），说明字节跳动在教育场景的AI能力正在从“工具”向“习惯”迁移。而Gauth虽然流水仍在增长，但用户增长放缓意味着它可能需要更强的变现能力来维持估值逻辑。

这里尚未完全回答的问题是：豆包爱学的用户增长能否转化为可持续的付费？GS报告提供了DAU和时长数据，但没有披露豆包爱学的变现模型。如果字节跳动选择以广告而非订阅为主要变现方式，那么它对好未来等传统教育公司的竞争压力将不是直接替代，而是挤压用户时间，进而影响后者的获客成本。

这是本期报告中最容易被误读的数据。6大学习平板品牌5月线上GMV同比下降35\%，相比4月的-2\%大幅恶化。好未来5月GMV约3.55亿元，同比下降24\%，1QFY26（5月季度）GMV同比下降1\%，低于GS此前预期的同比增长11\%。

表面看，这是一个负面信号。但拆开数据看，真正值得关注的不是总量下滑，而是ASP（平均售价）的同比降幅在持续收窄。好未来学习平板的ASP在5月同比下降仅4\%，环比下降7\%至约3700元。这个降幅收窄的趋势已经持续了数月，与此前动辄双位数下降形成鲜明对比。

这意味着什么？学习平板市场正在经历一场“量价分离”的结构性调整。去年同期的高基数效应导致销量同比下滑，但ASP的企稳表明产品结构在升级，消费者愿意为更高配置、更强AI功能的产品支付溢价。GS报告没有直接说，但这个数据隐含的逻辑是：如果ASP降幅继续收窄甚至转正，那么即使销量增速放缓，头部品牌的总收入仍有可能保持正增长。

好未来1QFY26的GMV低于预期，但ASP降幅收窄给了市场一个观察窗口：如果2QFY26的ASP同比能回到平盘甚至转正，那么当前市场对好未来硬件业务的悲观预期可能需要修正。报告没有展开的是，ASP企稳是否来自产品组合中高端机型占比提升，还是来自提价能力恢复？这需要后续跟踪每款机型的SKU数据才能验证。

中国学生赴加拿大、澳大利亚和英国的留学签证获批数量在1Q26同比下降9\%，相比1Q25的-22\%和2025全年的-16\%明显收窄。这是一个明确的触底信号。

但这里有一个容易被忽略的结构性问题：三个目的国的分化。报告给出了合计数据，但没有逐一拆解加拿大、澳大利亚和英国各自的趋势。从政策环境看，加拿Daiwa澳大利亚在 年都收紧了留学生签证政策，而英国相对稳定。如果合计降幅收窄主要是由英国拉

[... middle omitted ...]

续下降，好未来的“双轨”策略正在接受检验

报告披露了一个持续但未被充分讨论的趋势：非学科培训牌照数量在5月继续下降。这不是新消息，但它的持续性值得关注。自2023年以来，非学科培训牌照的发放一直处于收紧通道，这与市场此前预期的“监管放松后牌照放开”完全相反。

同时，好未来的两个主要移动端产品表现出分化：学而思网校MAU同比增速在5月放缓至10\%（4月为22\%），而培优MAU维持在19\%的同比增长。培优是面向学有余力学生的“培优”课程，网校则是更广泛的在线学习平台。两者的增速分化说明，在监管环境收紧和需求结构变化的背景下，好未来的用户增长正在向高客单价、高粘性的“培优”业务集中。

这对好未来的估值逻辑有直接影响。市场给好未来的估值中，通常将学习平板和培优业务视为增长引擎，而网校更多被视为流量入口和品牌阵地。如果网校增速持续放缓，而培优保持稳健，那么好未来的整体用户增长天花板可能比市场预期的要低。但同时，如果培优用户的ARPU值能持续提升，那么收入质量反而可能改善。

\subsection{[R023] GS：资本开支的“结构性集中”正在重塑全球工业投资逻辑}
{\small\color{refgray}归类：AI与科技：资本开支结构性集中，光互连与PCB成为下一波赢家}

全球资本开支正在经历一场罕见的双重分化。一方面，总量数据持续向好，2026年预期增速高达13\%，远超历史中位数2\%的水平；另一方面，增长几乎完全由数据中心、半导体和电力公用事业三个领域驱动，这三个领域在总资本开支中的占比已从2022年的22\%跃升至2026年的39\%。这份GS最新发布的跨行业资本开支追踪报告（Capex Tracker）揭示了一个核心判断：市场当前低估的不是资本开支总量的韧性，而是增长结构高度集中所带来的供应链定价权转移和周期性风险错配。

这份报告覆盖了约4000家公司、34个终端和子终端市场，追踪规模高达3.4万亿欧元的资本开支计划。其结论并非简单的“资本开支强劲”，而是指向一个更深刻的产业格局变化：少数几个结构性增长领域正在吸收越来越多的投资资源，而传统周期性行业正在被挤出。这种集中度达到历史极值，意味着投资者需要重新评估哪些环节真正享有议价能力，哪些资产可能面临需求透支后的回调风险。

1. 资本开支增长的结构性集中已达到历史极值，传统周期行业正被系统性挤出

报告最值得关注的数字不是总量增速，而是增长分布的变化。2026年预期资本开支增速为13\%，2027年预期为12\%，均显著高于历史趋势。但拆解来看，增长贡献高度集中：数据中心 年复合增速约为35\%，半导体约为23\%，油气约为4\%。与此同时，化学品和工程机械领域的资本开支预期分别下调了4.1和4.0个百分点，机场领域的复合增速甚至为-4\%。

这意味着，全球资本开支的增量几乎全部来自少数几个与AI、能源转型相关的领域。传统制造业、化工、建筑设备等行业的投资预期正在收缩。这种结构性集中并非短期现象——从2022年到2026年，数据中心、半导体和电力公用事业三个领域在总资本开支中的占比翻了一倍，从22\%升至39\%。换句话说，全球工业投资的“马太效应”正在加速：赢家通吃的逻辑从消费互联网蔓延到了重资产投资领域。

对于投资者而言，这意味着不能再将资本开支视为一个整体性指标。总量的乐观掩盖了结构性的分化。那些处于资本开支集中领域的上游供应商——尤其是提供电力设备、冷却系统、特种气体、精密零部件的公司——可能享有持续的需求溢价和定价权。而那些依赖传统周期行业订单的企业，则面临需求增量不足、竞争加剧的风险。

2. 数据中心和半导体是资本开支增长的双引擎，但供应链的瓶颈正在从芯片转向电力与散热

报告明确指出，数据中心和半导体是资本开支增长最快的两个领域，复合增速分别达到35\%和23\%。GS在报告中直接点名了施耐德电气（Schneider Electric）和VAT Group作为买入标的，前者是数据中心电力基础设施的核心供应商，后者是半导体真空阀门领域的全球龙头。

但这份报告更深层的含义在于，资本开支的流向正在揭示供应链瓶颈的迁移。过去两年，市场关注的焦点是AI芯片的产能和供应。然而，随着数据中心建设规模的指数级扩张，瓶颈正在从芯片本身转向配套的电力、散热和基础设施。一个数据中心的生命周期中，电力成本占比高达30\%-40\%，而冷却系统的能耗又在其中占据相当比例。这意味着，电力设备和热管理领域的资本开支增长可能比半导体设备更具持续性。

报告还提到，电力公用事业领域的资本开支增长同样强劲，特别是电网建设。这并非偶然。数据中心的高密度用电需求正在倒逼全球电网升级，而新能源并网同样需要大规模的电网投资。Prysmian和Nexans这两家电缆制造商被GS列为买入标的，正是这一逻辑的直接体现。

对于产业决策者而言，这意味着供应链的“卡脖子”环节正在从芯片制造向电力基础设施迁移。那些能够提供高效电力转换、散热和电网互联技术的公司，可能在未来几年享有结构性增长红利。

3. 结构性增长领域与周期性行业的背离正在扩大，两者面临完全不同的风险类型

报告中

[... middle omitted ...]

给予更高的估值溢价，因为其增长确定性更强。但对于周期性行业的公司，即使估值看起来便宜，也可能面临“价值陷阱”——需求持续低迷导致盈利无法兑现。报告的数据显示，2026年资本开支预期增速高达13\%，但这一数字主要由少数领域拉动，对于大部分传统工业企业而言，实际可获得的订单增量可能远低于这一数字。

4. 报告尚未完全回答的关键问题：资本开支集中度达到极值后，是否存在回调风险

尽管报告呈现了清晰的结构性增长逻辑，但有一个关键问题尚未得到充分讨论：当资本开支的集中度达到历史极值时，是否存在系统性回调的风险？

从历史经验看，资本开支集中度的快速上升往往伴随着过度投资的风险。当大量资本涌入少数几个领域时，产能建设速度可能超过实际需求增速，导致供给过剩、回报率下降。数据中心和半导体领域目前正处于这一阶段。虽然AI需求增长迅速，但数据中心建设的规模已经引发了对电力供应和碳排放的担忧。如果AI应用落地速度低于预期，或者能源成本上升超出预期，那么当前激进的资本开支计划可能面临修正。

\subsection{[R024] JPM：中国房地产并非整体回暖，而是一场K型分化下的资产再定价}
{\small\color{refgray}归类：地产与消费：中国楼市K型分化，香港写字楼迎结构性拐点}

JPM：中国房地产并非整体回暖，而是一场K型分化下的资产再定价

市场对“中国房地产复苏”的讨论，大多仍停留在总量层面：销售面积是否企稳、价格是否见底、政策是否继续加码。但JPM最新发布的实地调研报告，提供了一个远比总量数据更尖锐的视角——深圳和上海的地产市场正在经历一场清晰的“K型复苏”。

所谓K型，即市场的不同层级正在走向截然不同的方向：高端住宅需求旺盛、去化迅速，入门级住房交易量有所恢复，而中间价位段则持续承压。这不是一个“复苏”的故事，而是一个“分化加剧、资产重新定价”的故事。

这份报告的独特价值不在于其结论本身——分化是行业共识——而在于它通过实地调研、专家访谈和项目走访，给出了分化的具体形态、背后的驱动因素，以及对不同资产类别的含义。对于持有地产信用债、关注商业地产REIT化进程，或者正在思考中国资产配置框架的读者而言，这份报告提供了一组不可忽视的微观信号。

1. K型分化的本质不是消费降级，而是财富效应与收入预期的结构性裂口

JPM在深圳和上海走访了多个高端住宅项目，包括深圳宝安的观潮、上海徐汇的安澜上海以及杨浦的翠湖滨江。这些项目总价从2000万元至超过1亿元不等，但销售情况出奇一致：每批次开盘即售罄，买家多为30-40岁、从事科技、金融或AI/半导体行业的本地精英。报告引用销售经理的描述，有客户年收入达1亿元。

与此同时，入门级住房（约300万元总价）的交易量在经历了过去五年的价格大幅下跌后有所恢复，但中间价位段（500万至1000万元）表现最弱。

这不仅仅是“有钱人更有钱”的老生常谈。报告揭示了一个更关键的结构性变化：高端市场的购买力来源于AI和机器人产业带来的正向财富效应，而大众市场受制于就业安全感和工资预期。换句话说，这不是消费降级或升级的问题，而是不同收入群体的“收入预期”正在走向截然不同的路径。对于以中间价位段为主的开发商和持有相关信用债的投资者而言，这个裂口意味着需求端的结构性支撑不足，定价压力将持续存在。

2. 深圳的“政策底”已经出现，但供给端的新变量可能压制大众市场恢复

报告详细梳理了深圳2026年上半年的三轮政策放松：增值税减免、商业地产首付降至30\%、核心区购房限制放松。这些政策确实带来了效果——5月份总成交量突破1万套，创14个月新高。专家也确认，深圳二手房均价从2020年高点下跌43\%后，自2026年1月以来已反弹超过2\%，价格底部信号明显。

但JPM并没有停留在“政策有效”的叙事上。报告特别指出一个被市场忽视的供给端变量：深圳计划在2026至2030年间推出20万至30万套保障性住房，定价为可比商品房的50\%至70\%。2025年深圳已推出4.3万套保障房，是当年新增商品房供应量（3.5万套）的1.2倍。

这个数字的含义是什么？大众市场的需求恢复，正面临一个规模空前的替代性供给冲击。即使政策放松带动了短期成交量回升，中低价位商品房的价格恢复空间可能被保障房供给显著压缩。对于布局在低线城市的开发商而言，这是一个需要持续跟踪的风险因素。

3. 商业地产的分化比住宅更精细：地段、租户组合、运营策略决定一切

报告不仅走访了住宅项目，还考察了多个商业地产项目，包括深圳南山的K11 ECOAST、上海静安的恒隆广场66和龙湖闵行天街。结果同样是分化：恒隆广场66的高端定位带来了超过10\%的季度销售额同比增长，而K11 ECOAST在工作日午间时段人流稀疏。

龙湖闵行天街的销售经理将万达视为主要竞争对手，并认为万达更关注短期利润而非购物体验。这看似是一句运营层面的评论，但背后反映了一个更深刻的竞争逻辑：在低线城市，商业地产的护城河来自于对当地消费习惯的长期培育，而非简单的品牌招商。

对于持有商业地产相关信用债的投资者，这份报告提供了一个重要的筛选维度：不

[... middle omitted ...]

收入（租金加物业管理费）在2026年前四个月同比增长3\%，达到47亿元。管理层预计2026年全年商业运营收入为145亿元，同比增长4\%。

更值得关注的是REIT作为新融资渠道的探索。新城已以上海青浦吾悦广场为底层资产发行了私募REIT，通过出售66\%的权益募集了4亿元，并计划在2026年内发行公募REIT，规模约20亿元。管理层认为，REIT渠道有助于改善债务结构，以股权替换债务。

这个故事的逻辑是清晰的：开发业务收缩、IP业务稳定、REIT打开融资新通道。但报告也留下了两个关键问题：第一，20亿元的公募REIT发行相对于120亿元的IP资产总量而言，规模仍然有限；第二，低线城市商业地产的租金增长能否持续支撑更高的估值倍数，尚需时间验证。新城2028年到期债券的收益率已从19\%压缩至13\%，JPM对此维持中性评级——估值已经反映了部分预期，但尚未完全兑现。

报告在深圳专家访谈中明确指出，大规模保障房供给可能在未来几年压制大众市场恢复。但这一判断仍有几个未解之处。

\subsection{[R025] 摩根斯坦利：市场对楼市“小阳春”的可持续性存在显著误判}
{\small\color{refgray}归类：地产与消费：中国楼市K型分化，香港写字楼迎结构性拐点}

在经历了3-4月超预期的二手房成交回暖后，市场情绪一度转向乐观。然而，摩根斯坦利最新发布的6月高频数据显示，这一轮复苏的动能正在迅速衰减。截至6月10日，25城二手房实时成交同比增速已从4月的30\%骤降至9\%，低能级二线城市甚至出现同比负增长。这并非一次简单的季节性回调，而是政策脉冲效应消退与积压需求释放完毕后的结构性回落。市场真正需要重新定价的，不是需求何时反弹，而是供给侧的出清速度与房价下行压力将持续多久。

为何这个判断在当前时点尤为重要？因为市场对“小阳春”的解读存在根本分歧。部分投资者将其视为周期拐点的前兆，而摩根斯坦利在5月发布的报告《An Inflection or Another False Start?》中已明确指出，3-4月的反弹更可能是积压需求的集中释放，而非基本面改善。如今6月数据进一步印证了这一判断，但股价自5月中旬以来已回调18\%，而恒生指数同期仅下跌8\%。这意味着，市场对房地产板块的风险定价仍不充分，尤其是在政策效果递减、二手房挂牌量未见显著收缩的背景下。

这份报告的核心贡献在于，它通过高频数据拆解了25个城市的成交结构，揭示了三个关键信号：第一，政策放松的边际效用正在递减，且城市间分化加剧；第二，低能级城市的二手房市场已率先转弱，这可能是一线城市的前瞻指标；第三，即便在一线城市内部，成交量与价格走势也出现背离。这些信号合在一起指向一个判断：2026年下半年，二手房成交同比增速可能转负，而房价将在 年持续承压，仅极少数库存去化较好的核心城市可能出现温和回升。

摩根斯坦利的数据显示，25城二手房实时成交同比增速从4月的30\%降至5月的26\%，再到6月上旬的9\%，下滑速度远超季节性因素所能解释。尤其值得关注的是，低能级二线城市——如南昌、佛山、南通、东莞、宁波、西安——的减速幅度最为明显。而天津、长沙、成都三城甚至已出现同比负增长。

这意味着什么？3-4月的成交放量，本质上是由政策放松（包括公积金贷款额度提升、房贷利率下调）与居民情绪短期修复共同驱动的。但政策对需求的刺激效果具有一次性特征：它能够加速积压需求的释放，却无法创造新的购买力。当这批“存量需求”消化完毕后，市场便回到了由收入预期、就业前景和房价预期决定的基本面轨道上。低能级城市之所以率先转弱，恰恰是因为这些城市的基本面支撑更为脆弱——人口流入放缓、产业基础薄弱、库存压力更大。

对投资者的含义在于：不能再用3-4月的高增速作为线性外推的依据。政策脉冲的衰减速度比市场预期的更快，而低能级城市的转弱可能是全市场成交下行的领先指标。

报告特别指出，4月底政策放松的城市出现了显著分化：苏州和武汉的成交依然强劲，而广州、深圳、佛山则出现了两位数的环比减速。这并非随机分布，而是与各城市的库存水平和去化能力密切相关。

苏州和武汉的共同特征是什么？两者在过去两年经历了较大幅度的价格调整，二手房性价比已相对合理，且库存去化周期处于中等偏低水平。而广州、深圳虽然基本面更强，但二手房挂牌量持续攀升，导致买方议价能力增强，成交量难以持续放大。佛山则属于典型的低能级城市，政策放松难以对冲人口和产业基本面的弱势。

这一分化的核心启示是：在需求侧政策工具接近用尽的背景下，供给侧的去库存进度将成为决定城市房价走势的核心变量。摩根斯坦利在报告中也提到，一线城市可能因更快的去库存进度而出现温和的价格回升，但这一判断仍需验证——因为一线城市的二手房挂牌量仍在增加，价格下行压力并未解除。

报告虽然没有直接给出25城的挂牌量数据，但通过成交增速的快速下滑，可以合理推断：在成交量收缩的同时，挂牌量并未同步减少。事实上，许多城市的二手房挂牌量仍在历史高位附近。这意味着供需关系正在进一步恶化——供给端没有收缩，而需求端已经疲软。

这是一个关键但容易被忽视的

[... middle omitted ...]

4. 报告尚未完全回答的关键问题：二手房价的下跌是否已经传导至新房市场

摩根斯坦利在报告中强调了二手房成交转弱对房价的传导机制，但并未充分展开一个关键问题：二手房价格的下跌是否已经系统性地传导至新房市场？从历史经验看，二手房价格通常是新房价格的领先指标。当二手房市场出现量价齐跌时，新房市场的去化压力会进一步加大，尤其是在开发商仍面临资金链压力的背景下。

目前新房市场的表现如何？报告没有提供直接数据，但我们可以从逻辑上推导：如果二手房成交转弱，意味着置换链条中的“卖旧买新”需求将受到抑制。同时，二手房价格下跌会压低新房的价格锚，使得开发商不得不以更大的折扣促销，从而进一步压缩利润率。这对高杠杆开发商而言是双重打击——既面临销售回款放缓，又面临资产减值压力。

这一问题之所以重要，是因为它直接关系到房地产板块的投资逻辑：如果新房市场尚未充分反映二手房的下行压力，那么当前板块的估值修复可能并不稳固。这也是摩根斯坦利在报告中维持“行业风险回报偏向下行”判断的核心原因之一。

\subsection{[R026] 摩根斯坦利：重卡销量超预期背后，真正的分化信号才刚刚开始}
{\small\color{refgray}归类：未被主题摘要引用}

2026年5月，中国重卡市场交出了一份超出预期的成绩单。中国汽车工业协会数据显示，当月重卡销售10.95万辆，同比增长23\%，比行业研究机构CV World此前预测的10.3万辆高出约6个百分点。这份来自摩根斯坦利的月度数据追踪，表面上只是一次常规的行业销量更新，但如果我们把视线从单月数字上移开，转而审视前五个月的结构性变化，会发现一个更值得关注的判断正在浮现：市场对重卡行业的关注点，正在从“总量增长”转向“结构分化”，而后者对投资决策的含义，远比前者深刻。

2026年前五个月，中国重卡累计销售54.43万辆，同比增长23\%。这个增速本身并不令人震惊——它更多反映的是去年同期低基数下的恢复性增长。真正值得追问的问题是：谁在增长？谁在丢失份额？增长的动力来自哪里？这些问题的答案，将直接决定未来12至18个月重卡行业竞争格局的走向。

这份报告的真正价值，不在于它告诉我们重卡卖了多少辆，而在于它揭示了行业内部正在发生的、投资者尚未充分定价的再分配过程。

5月重卡销量超预期，表面上是一个好消息。但摩根斯坦利的数据显示，超预期的背后是几家欢喜几家愁。

中国重汽是5月最大的赢家。单月销售3.45万辆，同比增长41\%，市场份额达到31.5\%，同比提升3.9个百分点。这一表现背后有两个支撑：一是出口持续强劲，单月出口超过1.8万辆；二是在新能源重卡领域也处于领先位置。出口和新能源的双轮驱动，让中国重汽在行业整体复苏中跑出了加速度。

对比之下，陕汽的表现令人担忧。5月销售约1.5万辆，同比下滑3\%，市场份额13.7\%，同比下降3.7个百分点。前五个月累计市场份额也下降了1.2个百分点。在行业整体增长23\%的背景下，陕汽的份额流失格外显眼。

北汽福田和徐工则实现了小幅的份额扩张。福田5月销售1.6万辆，同比增长37\%，前五个月累计市场份额提升1.4个百分点。

这些数字放在一起，意味着什么？行业总量的增长是均匀的吗？不是。增长的果实正在被少数玩家集中收割。当市场从低谷恢复时，往往不是所有参与者都能分享复苏红利，而是强者恒强的格局在加速固化。这是投资者需要建立的第一层认知。

2. 出口已成为区分赢家与输家的关键变量，且这一趋势远未到头

中国重汽5月出口超过1.8万辆，占其单月销量的比重超过52\%。这是一个值得反复咀嚼的数字。这意味着中国重汽已经不再是一家单纯的“国内重卡制造商”，而是一家业务重心正在向海外转移的全球化企业。

出口对于重卡企业的意义，不仅仅是多卖几辆车那么简单。海外市场的利润率通常高于国内市场，因为竞争格局更分散、定价权更强、品牌溢价空间更大。更重要的是，出口提供了国内周期波动的对冲。当国内基建投资放缓或排放标准切换导致需求波动时，海外市场可以充当稳定器。

摩根斯坦利的报告没有展开分析出口的结构，但这里有一个关键问题值得追问：中国重汽的出口是去了哪些市场？是一带一路沿线的基础设施建设需求，还是成熟市场的替换需求？不同市场的增长可持续性和利润率差异很大。如果出口主要流向新兴市场，那么地缘政治风险和汇率风险就是不可忽视的变量。报告没有给出答案，但这个问题本身就是投资者需要持续追踪的线索。

3. 新能源重卡的渗透率正在成为下一个区分维度，但数据缺口需要填补

报告提到中国重汽在新能源重卡领域处于领先位置，但没有给出具体的新能源渗透率数据。这恰恰是目前市场分析中最容易被忽略的盲区。

新能源重卡与传统燃油重卡的竞争逻辑完全不同。新能源重卡的核心竞争力在于电池成本、充电基础设施和运营效率，而非发动机技术或经销商网络。这意味着，今天在传统重卡领域领先的企业，未必能在新能源赛道上保持优势；反之亦然。

从行业趋势看，新能源重卡的渗透率正在从政策驱动转向经济性驱动。随着电池成本下降和换电站网络的完善，新能源重卡在

[... middle omitted ...]

在54.43万辆的总量中意味着超过1600辆的增量份额。与此同时，陕汽、东风等企业的份额在下降。一汽的份额也略有下滑。

行业集中度的变化往往是一个缓慢但不可逆的过程。当头部企业通过出口和新能源建立新的竞争壁垒时，二线企业的追赶难度会越来越大。规模效应、供应链议价能力、研发投入强度，这些因素会形成正反馈循环，让领先者的优势自我强化。

对于投资者而言，这意味着重卡行业的投资逻辑正在从“行业贝塔”转向“个股阿尔法”。过去几年，投资重卡行业更多是赌行业周期的方向和幅度。但未来，行业内部的差异化将越来越大。选对企业的回报，可能远超行业平均水平的增长。

摩根斯坦利这份报告的核心价值在于数据的及时性和颗粒度，但它本质上是一份“快照”，而不是“深度透视”。有几个关键问题，报告没有回答，但投资者必须自己寻找答案。

第一个问题是增长的可持续性。5月销量超预期，有多少是政策刺激的短期效果，有多少是真实需求的恢复？如果下半年基数效应消退，增速是否会显著回落？报告没有给出前瞻性的预测。

\subsection{[R027] 摩根斯坦利：香港地产市场真正被低估的不是住宅，而是写字楼的结构性拐点}
{\small\color{refgray}归类：地产与消费：中国楼市K型分化，香港写字楼迎结构性拐点}

摩根斯坦利：香港地产市场真正被低估的不是住宅，而是写字楼的结构性拐点

这份由摩根斯坦利发布的香港地产研报，表面上是一次常规的行业专家电话会纪要。但如果我们把它放在2026年这个时间节点来读，会发现它传递了一个比“住宅销售回暖”远更重要的信号：香港甲级写字楼市场正在经历一个由供给端驱动的结构性拐点，而市场情绪对此的定价远未充分。

报告的核心价值不在于确认住宅市场“量价齐升”的已知事实，而在于揭示了两个被当前市场叙事忽略的关键变量：一是写字楼供给的未来路径已经锁定，二是学生公寓正在从边缘话题变成可规模化的新资产类别。理解这两个变量，才能理解摩根斯坦利为何在行业近期表现不佳、政策不确定性犹存的背景下，依然维持“Attractive”的行业评级。

这不仅仅是一份关于地产周期的报告。它是一份关于“资产稀缺性如何被重新定价”的底层逻辑推演。

报告中最直观的数据是：截至2026年第一季度，香港已完工住宅库存从2025年第一季度的28,000套峰值下降至20,000套，降幅达30\%。与此同时，5月份住宅销售量同比增长44\%，连续第15个月维持在5,000套以上。

这两个数据放在一起，揭示的不仅仅是需求回暖。它意味着开发商正在从“去库存压力下的价格竞争”转向“库存可控后的定价主动”。当库存从28,000套降至20,000套，市场已经从买方市场重新回到了卖方市场的边缘。

摩根斯坦利预测2026年全年住宅价格将上涨7\%-10\%，而截至目前的官方数据显示，年初至今已上涨约6\%。如果这个预测成立，那么上半年6\%的涨幅意味着下半年仍有约4\%的上涨空间。但更值得关注的不是价格本身，而是这个价格走势背后的驱动因素从“政策刺激”切换到了“供需基本面”。

报告没有明确展开的是：库存下降30\%的背后，有多少是主动去化，有多少是因为开发商刻意放缓了推盘节奏？这里仍需验证。但无论如何，库存水平的快速下降已经为开发商提供了更从容的定价空间，而这恰恰是过去两年市场最稀缺的东西。

2. 甲级写字楼正在经历一个罕见的“需求修复+供给封顶”双重窗口

如果说住宅市场是“已知的好转”，那么写字楼市场才是这份报告真正的洞察所在。

报告指出，中环和西九龙地区的甲级写字楼正在引领复苏，银行和金融业是需求的主力——2026年第一季度新租约中，银行与金融业占比高达61\%。更关键的是，2026年至2029年间，全港甲级写字楼的新增供应总量仅为260万平方英尺。对于一个年净吸纳量连续保持正值的市场来说，这个供应量意味着什么？

意味着未来三年，核心区域的写字楼供给将进入一个事实上的“封顶期”。

摩根斯坦利援引戴德梁行的预测：大中环和核心中环的写字楼租金将在2026年分别上涨6\%-8\%。这个预测的底气就来自供给断档。当需求端的银行和金融业仍在扩张（无论是由于家族办公室涌入还是中资机构扩大布点），而供给端的新增面积几乎可以忽略不计时，租金上涨就从一个“可能”变成了“大概率”。

这里有一个容易被忽略的细节：报告中提到写字楼净吸纳量在2026年第一季度为21.7万平方英尺，仍然为正。在传统认知中，香港写字楼市场长期受困于高空置率，但这份数据表明，最坏的时刻可能已经过去。真正的问题不是需求不够，而是过去几年的供给释放过于集中。一旦供给端进入收缩期，租金弹性可能会超预期。

报告在总结部分轻描淡写地提到了一句话：“学生公寓是一个新的机会，香港八所主要大学的学生与床位比例为3.4:1，床位数量大约相当于非本地学生数量的一半。”

3.4:1的学生床位比，意味着每3.4个学生只有1个床位可用。考虑到香港高校的非本地学生比例仍在上升通道中，这个缺口只会扩大而非缩小。报告给出的数据是“床位数量大约相当于非本地学生数量的一半”，这意味着至少有50\%的非本地学生无法在校内住宿，必

[... middle omitted ...]

从“边缘话题”变成“核心赛道”。

4. 报告尚未完全回答的关键问题：政策与利率的不确定性如何落地？

这份报告最坦诚的地方在于，它明确承认了三个不确定性来源：新的境外投资规则、股市波动以及利率路径的不确定性。报告的原话是“目前尚无实质影响”，但这恰恰是最需要警惕的表述。

“尚无实质影响”不等于“不会产生影响”。新的境外投资规则如果收紧，将直接影响内地资金流入香港住宅市场——而内地买家在过去几年一直是高端住宅市场的重要支撑。股市波动则会影响财富效应，进而抑制高端需求的释放。利率路径的不确定性更是直接决定了按揭成本和开发商的融资成本。

报告对这些变量采取了“观察但不预判”的姿态。这背后隐含的判断可能是：即便在最悲观的情景下，香港地产的估值（相对于NAV折让50\%）也提供了足够的安全边际。但安全边际不等于没有下行风险。对于决策者来说，真正需要追问的是：如果利率维持高位的时间比预期更长，或者政策收紧的幅度超出预期，哪些公司的资产负债表最脆弱？哪些资产类别会最先受到冲击？

\subsection{[R028] 摩根斯坦利：韩国居民资产迁移到股票市场的趋势具有结构性，而非周期性}
{\small\color{refgray}归类：区域市场：韩国居民资产结构性迁移，日本治理改善信号被低估}

摩根斯坦利：韩国居民资产迁移到股票市场的趋势具有结构性，而非周期性

韩国正在经历一场深刻的居民资产负债表重构。过去二十五年间，韩国居民资产增长了91\%，增速在主要经济体中位居前列。但真正值得关注的，不是总量增长，而是资产配置方向的根本性转变——从长达数十年的房地产与现金偏好，转向权益资产。摩根斯坦利在最新发布的研报中明确指出，这一转变具有可持续性，并将KOSPI的12个月目标从8500点上调至9000点，牛市情景上调至10500点。

这份报告的核心判断是：韩国居民资产向权益市场的迁移，不是一次短期的交易性行为，而是一个由政策、人口结构与市场改革共同驱动的结构性拐点。对于关注亚太资产配置的投资者而言，这可能是未来三到五年最重要的宏观叙事之一。

1. 居民资产配置的结构性拐点已经到来，但市场尚未充分定价其持续性

韩国居民长期以来将财富高度集中于房地产。截至2025年，主要住宅与其他房产合计占比超过70\%，而股票、债券与基金的合计占比仅为3.6\%。相比之下，近一半的金融资产仍以现金和存款形式持有。这种配置结构在过去二十年中几乎没有变化，直到2025年出现了显著拐点。

数据显示，2025年韩国居民权益资产持有量同比增长48\%，而此前十年的平均增速仅为7.4\%。这一跳跃并非偶然。报告指出，房地产政策的持续收紧——包括信贷限制与对多套房持有的税收惩罚——正在系统性地降低房地产作为财富储存工具的吸引力。与此同时，股票市场的改革措施、企业治理改善以及股东回报提升，正在为权益资产创造更具吸引力的制度环境。

这意味着，居民资金从房地产向权益市场的迁移，其驱动力不是短期的市场情绪，而是持续三到四年的政策方向。报告认为，这种“错失恐惧”已经从房地产转向股票，并且这种转变具有自我强化的特征：随着更多居民参与权益市场，财富效应的正向反馈将进一步巩固这一趋势。

市场对韩国零售投资者的传统印象是“追逐波动、短线交易、主题炒作”。摩根斯坦利在2020年的一份深度分析中确实得出了这一结论。但五年后的今天，报告认为情况已经发生了实质性变化。

变化来自三个层面。第一，信息的获取质量显著提高。第二，人工智能工具正在缩小信息不对称，为零售投资者提供了更平等的分析能力。第三，海外投资的实践经验积累，使得韩国零售投资者的交易行为更加成熟。

从数据上看，零售投资者在 年周期中的净买入高度集中，而当前周期的净买入分布更加分散。这意味着投资者不再盲目追逐单一主题，而是采取了更为战术性和控制性的投资策略。报告还注意到，零售投资者的数量已经增长到1450万人，占韩国成年人口的50\%，而2019年这一比例仅为21\%。

这些变化的意义在于：零售投资者不再只是市场的波动源，而是正在成为市场的稳定器。一个更加成熟、更具信息优势、更分散的零售投资者群体，将减少市场的极端波动，并为优质公司的估值提供更持续的支撑。

3. 财富效应的传导机制正在放大，但这也意味着更高的宏观风险

报告对财富效应的分析值得特别关注。历史数据显示，KOSPI每上涨1\%，大约可以拉动商品消费增长 个百分点。但报告认为，当前的财富效应传导机制正在发生变化，其影响范围更广、强度更大。

原因有三。第一，参与者的基数已经大幅扩大——1450万零售投资者意味着财富效应对消费的渗透率远高于过往周期。第二，权益资产的持仓规模从2019年的17.6万亿韩元跃升至2025年的662万亿韩元，增长了13.8倍，这意味着股价变动对居民财富的影响力度显著提升。第三，新增投资者中，20-40岁年龄段的群体占零售投资者的43\%，这一群体的边际消费倾向更高，财富效应对消费的传导效率因此更强。

但硬币的另一面是风险。报告明确指出，韩国居民正在进入一个风险资产敞口升高的时期，这带来了新的宏观脆弱性。到2027年，

[... middle omitted ...]

的转变则更加迅速和激进，投资者表现出更高的风险偏好。

这种差异体现在数据上：过去十年，韩国的金融资产增长了94\%，而日本仅为30\%。韩国的零售投资者不仅更早地拥抱了权益资产，而且在交易行为上更加主动。

这一对比对于理解两个市场的未来走势至关重要。日本的经验表明，政策驱动的资产配置转变可以持续多年，但速度较慢，对市场的冲击相对温和。韩国的路径则意味着，市场可能会经历更剧烈的资金流入，但也面临更大的波动风险。对于投资者而言，关键问题不是“韩国是否会复制日本”，而是“韩国的更快速度是否意味着更早面临调整压力”。

5. 报告尚未完全回答的关键问题：零售投资者行为的可持续性如何验证？

尽管报告对零售投资者的行为进化持积极态度，但一个关键问题仍然悬而未决：这种“更聪明”的投资行为是否能够在市场下跌周期中得到验证？目前的数据主要来自上涨周期，零售投资者的分散化买入和战术性操作在牛市中更容易保持理性。但一旦市场转向， 年周期中零售投资者表现出的“追逐波动”特征是否会重新出现？

\subsection{[R029] JPM：电力设备板块的20-30\%回调，不是风险出清，而是结构性机会的重新定价}
{\small\color{refgray}归类：电动车与供应链：中国渗透率结构性跃迁，电力设备回调是重新定价}

JPM：电力设备板块的20-30\%回调，不是风险出清，而是结构性机会的重新定价

过去一个季度，亚洲电力设备公司股价从四月高点普遍回撤20\%至30\%。如果只看价格信号，这很容易被解读为行业景气见顶的信号。但JPM最新发布的研报给出了一个相反的判断：这轮回调不是因为基本面恶化，而是市场在多重噪音中暂时失去了定价焦点。真正的结构性逻辑——美国电网资本开支的结构性增长、数据中心之外的重型电气化需求、以及中国电力设备在欧盟市场的低渗透率——并未被破坏，反而因为回调变得更清晰。

这份报告最值得关注的核心判断是：回调后的估值水平，为那些订单结构偏向公用事业和电网、而非纯数据中心概念的公司，提供了更有吸引力的入场窗口。JPM特别点名了晓星重工（Hyosung Heavy）和威胜控股（Wasion Holdings），认为它们在当前估值下提供了明显的安全边际。

1. 二季度回调的真正驱动力不是基本面，而是估值、资金流和催化剂真空三重叠加

JPM的分析显示，这轮回调并非源于订单或盈利的实质性恶化。报告明确指出，一季度新订单数据创下历史新高后，市场面临的是一个“高基数下的真空期”——缺乏高频前瞻数据来确认订单增速能否持续。对于电力设备这类季度数据驱动的行业，这种信息真空很容易放大短期情绪波动。

更关键的是资金流层面的结构性变化。JPM观察到，在韩国市场，部分对冲基金正在将电力设备股作为空头仓位，以融资买入存储器等AI主题股票。这种“做空电力设备、做多AI”的配对交易，使得电力设备板块的走势与美股科技股出现了一定程度的联动。这意味着，这轮回调中有相当一部分是资金流驱动的，而非基本面驱动的。

对于中国电力设备公司，JPM认为资金流向AI主题板块是另一个拖累因素。数据显示，二季度中国AI主题板块资金流入高达56\%，而电力设备板块则净流出11\%。这种资金轮动造成了板块间的相对表现分化，但并未改变电力设备公司的订单和盈利前景。

2. 数据中心连接延迟的担忧被过度放大，真正支撑订单的是公用事业电网的结构性资本开支

市场最担心的风险点是数据中心连接延迟，认为这可能拖累新订单增长。但JPM的分析揭示了一个容易被忽视的结构性事实：对于韩国重型电气设备公司（如晓星重工和现代电气），数据中心在新订单中的占比不足10\%，订单主体来自公用事业和电网客户。这意味着，即使数据中心建设出现短期延迟，对这些公司的订单影响也非常有限。

真正的增长引擎是美国电网的资本开支。JPM判断，美国的输配电资本开支将继续保持结构性增长，这背后是电力系统老化、可再生能源并网需求、以及电气化趋势带来的多重驱动。晓星重工超过90\%的新订单来自美国公用事业和电网客户，这种订单结构使其成为本轮回调中防御性最强的标的之一。

这里需要指出的是，JPM的判断隐含了一个尚未完全验证的假设：美国电网资本开支的增速能否在未来2-3年维持当前水平。如果美国经济进入衰退周期，公用事业资本开支可能面临预算压力，这一点报告没有充分展开。

3. 欧盟对中国电力设备加征关税的风险被高估，市场占有率不足10\%是核心安全垫

针对市场对地缘政治风险的担忧，JPM给出了一个数据驱动的不对称判断：欧盟对中国电力设备加征关税的概率很低。核心逻辑在于，中国电力设备在欧盟的市场份额仍低于10\%，远未达到威胁本土供应商（如西门子能源）的程度。这与汽车行业的情况有本质区别——中国电动汽车在欧盟的市场份额和增长势头已经引发了明显的竞争焦虑。

此外，中国对欧盟出口的主要是初级电气设备，而非软件和系统。这意味着中国产品的竞争力更多体现在成本端，而非技术替代层面。在欧盟能源转型承诺推动下，电气设备需求正在系统性上升，而本土产能存在明显缺口。在这种供需失衡的背景下，欧盟各国大概率会继续采购中国产品。

JPM的这个判

[... middle omitted ...]

清晰的区分。对于韩国电气设备公司，晓星重工当前交易在2028年预期市盈率约20倍，是韩国电气设备中最便宜的标的。相比之下，LS Electric的2027年预期市盈率高达53倍，现代电气为29倍。这种估值差异反映了市场对不同公司订单结构、增长确定性和风险敞口的定价差异。

对于中国电力设备公司，JPM重点推荐威胜控股。该标的当前交易在约10倍一年期前瞻市盈率，同时保持年化超过20\%的盈利增长。这种“低估值+高增长”的组合在当前的电力设备板块中相当罕见。报告还提到了国电南瑞、许继电气等覆盖标的，但威胜控股显然是JPM在当前位置最看好的中国电力设备公司。

不过，这里有一个值得追问的问题：威胜控股的低估值是否反映了某些特定的风险折价，比如客户集中度、区域市场风险，或者治理结构问题？报告没有详细展开这些可能的估值压制因素，对于追求深度研究的读者来说，这些细节恰恰是判断“价值陷阱”还是“价值洼地”的关键。

5. 报告尚未完全回答的关键问题：订单增速能否在二三季度实现环比改善

\subsection{[R030] 摩根斯坦利：市场真正低估的不是厄尔尼诺本身，而是它对作物窗口的精准打击}
{\small\color{refgray}归类：能源与大宗：欧洲石油巨头隐含油价远低于共识，厄尔尼诺精准打击作物窗口}

摩根斯坦利：市场真正低估的不是厄尔尼诺本身，而是它对作物窗口的精准打击

NOAA刚刚确认厄尔尼诺已正式形成，并将“非常强”事件的概率从37\%大幅上调至63\%。这个数字正在快速从气象部门的概率表渗透到大宗商品交易员的头寸决策中。但摩根斯坦利这份最新研报的核心判断是：市场可能已经定价了“厄尔尼诺会发生”，但远未定价“它会在什么时间、什么地点、以什么强度击中最脆弱的作物窗口”。

这份由拉丁美洲农业研究团队Julia Rizzo主导的报告，以罕见的紧迫感回应了投资者最关心的问题——如果今年出现超级厄尔尼诺，大宗商品、通胀、汇率、股票将面临怎样的重置？报告没有给出简单的“涨或跌”答案，而是构建了一个基于作物窗口、区域差异和时序节奏的分析框架。以下是我们从中提炼的五个核心洞察。

1. 概率的大幅跳升只是表层，真正的变量是“海洋-大气耦合”能否在夏季完成确认

63\%这个数字本身已经足够引人注目。但摩根斯坦利报告真正有价值的判断，在于它区分了“概率”和“确认”之间的关键差距。NOAA的最新数据确实指向了一个可能跻身历史前三的强厄尔尼诺事件——与1982/83、1997/98和2015/16年类似。但报告同时明确指出，官方指引仍然是概率加权的结果，尚未确认这是一次创纪录事件。

关键在于“春季可预测性障碍”。当前模型的不确定性仍然较大，部分气象专家警告不要过早外推至“超级厄尔尼诺”情景。真正的确认信号需要观察：夏季海洋-大气耦合是否如期加强、Niño 3.4区域的升温是否持续超过+2.5°C阈值、以及不同模型在越过春季预测障碍后是否收敛。

这意味着，市场目前定价的更多是“厄尔尼诺已来”这个事实，而不是“它可能有多强”。如果后续确认信号逐一兑现，大宗商品和农业股票的定价将面临系统性重估。反之，如果事件强度不及预期，当前部分资产的价格可能已经包含了过高的风险溢价。

在所有大宗商品中，糖对厄尔尼诺的敏感度最高。报告的逻辑链条非常清晰：厄尔尼诺通常导致亚洲产区（尤其是印度和东南亚）干旱，而印度是全球第二大糖生产国。印度季风降雨低于正常水平的风险，直接对应全球糖供应的收紧。摩根斯坦利将糖列为“最清晰的上行风险”。

但报告同时强调了一个容易被忽视的时间维度：短期内，巴西中南部正在进行的压榨季可能带来充足供应，对价格形成压制。真正的影响可能要到2027/28作物年度才会充分显现。这意味着，糖的多头头寸需要耐心，也需要对库存消耗节奏有更精细的判断。

棕榈油和棉花也处于类似的逻辑中——干旱收紧供应，但价格反应的时点和幅度取决于厄尔尼诺是否在关键生长期造成实质性损害。咖啡的情况则更加路径依赖，不同产区的降水模式差异可能导致价格走势分化。

3. 大豆和玉米的复杂性在于：厄尔尼诺对巴西不同产区的作用方向截然相反

这是报告中最具洞察力的部分之一。对于大豆和玉米，简单的“看多或看空”框架完全失效。厄尔尼诺对巴西农业的影响是高度区域化的：在马托格罗索和MATOPIBA产区，降水模式异常可能对2026/27年度的播种和单产构成下行风险；而在南里奥格兰德州和阿根廷，厄尔尼诺通常带来有利的降雨改善。

玉米的情况更为复杂。巴西的二季玉米（safrinha）占全国产量的80\%，其种植窗口高度依赖于大豆收获的时间节点。如果厄尔尼诺导致大豆播种延迟，二季玉米的种植窗口将被压缩，进而影响单产潜力。这是一个典型的“时序风险”——不是厄尔尼诺本身有多强，而是它是否恰好击中这个脆弱的日历窗口。

对于投资者而言，这意味着不能将巴西农业视为一个同质化的“风险暴露”。仓位选择必须精细到产区、作物和时段的组合。

摩根斯坦利在报告中给出了明确的股票影响判断，这是机构研报中最具可操作性的部分。

受益方包括：圣马丁（SMTO3.SA），作为巴西最大的甘蔗加工商之一，

[... middle omitted ...]

及二季玉米的种植窗口。这些微观指标将成为判断股票走势的前瞻信号。

5. 报告尚未完全回答的关键问题：厄尔尼诺对通胀和航运的传导路径是否被低估

研报在讨论厄尔尼诺的宏观影响时，触及了食品通胀和巴拿马运河航运两个领域，但并未展开详细的量化分析。这正是值得投资者深入思考的方向。

食品通胀方面，报告指出厄尔尼诺将在本已高企的全球食品价格之上叠加供应侧压力。但这一传导链条的强度取决于：厄尔尼诺是否在多个主产区同时造成减产、各国贸易政策（如出口限制）是否会进一步放大价格波动、以及从大宗商品价格到终端零售价格的传导时滞。如果这些因素同时发酵，2027年的食品通胀压力可能超出当前市场的温和预期。

巴拿马运河的航运约束则是另一个值得关注的二阶影响。报告提到，厄尔尼诺可能再次导致运河水位下降，影响通行能力和运费。这对于全球农产品贸易流和集装箱运价都构成尾部风险。但航运市场的反应往往滞后于气象信号，且受到船舶运力供应侧的干扰。这一传导路径的量化评估，可能是完整报告中更有价值的部分。

\subsection{[R031] UBS：酒店业真正的拐点不在RevPAR，而在定价策略的深层分化}
{\small\color{refgray}归类：地产与消费：中国楼市K型分化，香港写字楼迎结构性拐点}

UBS：酒店业真正的拐点不在RevPAR，而在定价策略的深层分化

过去一周，中国酒店行业的整体RevPAR同比增长了3\%。单看这个数字，很容易得出“行业温和复苏”的结论。但如果拆开来看，UBS这份基于STR和Flight Master数据的周报揭示了一个更值得警惕的信号：行业内部的分化正在从“档次差异”演变为“定价逻辑的彻底分裂”——中端市场正在用降价换入住率，而高端市场则在靠供给侧约束和入境需求维持价格刚性。

这不是一个简单的“复苏强弱”问题，而是一个结构性拐点：当一部分玩家选择主动降价以争夺存量客源时，整个行业的利润池分布和竞争规则正在被重写。市场真正需要关注的，不是RevPAR的均值，而是这个均值背后的定价权转移。

UBS数据显示，5月31日至6月6日这一周，中档和经济型酒店的RevPAR表现出现明显分化。中档酒店RevPAR同比下降3.4\%，而经济型酒店同比上升2.8\%。更关键的是，中档和上中档酒店的RevPAR改善，主要来自降低平均房价（ADR）来拉动入住率（OCC）。具体来看，中档ADR同比下降约2.5\%，上中档ADR同比下降约2.5\%，而入住率则分别提升了约-1.0和2.8个百分点。

这意味着什么？中端酒店正在主动牺牲价格弹性来换取流量。这与一季度的情况形成了鲜明对比——当时行业更多是在价格相对稳定的情况下自然恢复。UBS报告明确指出，这标志着酒店集团正在“积极适应商务和休闲旅游市场的变化”。翻译过来就是：需求端的结构性疲软已经迫使中端酒店放弃价格纪律。

但问题在于，这种策略能否持续？降价换量的本质是争夺存量市场，而存量市场的总盘子并没有显著扩大。当所有中端玩家都开始降价，最终的结果可能是入住率没有根本性改善，但行业整体ADR被压低，利润被侵蚀。这类似于航空业的价格战逻辑——降价可以暂时拉升客座率，但无法解决供需失衡的根本问题。

一个尚未被充分讨论的风险是：如果中端酒店持续降价，是否会向下挤压经济型酒店的客源？从数据上看，经济型酒店本周RevPAR反而实现了2.8\%的增长，说明目前两个档次之间的替代效应还不明显。但这一动态值得持续跟踪。

与中端市场形成鲜明对比的是，豪华和高端酒店本周RevPAR分别增长了3.5\%和3.1\%，表现优于有限服务酒店。UBS给出了两个解释：一是高端和豪华板块新增供给相对较少；二是入境旅游增加拉动了该板块的需求。

这两个因素都很关键，但含义不同。供给侧约束是一个结构性利好——过去几年高端酒店的投资节奏放缓，新开业项目减少，这意味着存量高端酒店面临的竞争压力小于中端市场。而入境旅游的增加则是一个周期性变量，受签证政策、汇率和地缘政治等多重因素影响，其持续性需要进一步验证。

从数据细节看，豪华和上高端酒店的ADR同比仅增长0.5\%，但入住率提升了2.0个百分点；高端酒店的ADR甚至同比下降了3.0\%，但入住率提升了3.0个百分点。这说明高端市场同样在承受价格压力，只是由于入住率的强劲反弹，整体RevPAR得以转正。换句话说，高端酒店的价格韧性并非来自定价能力的提升，而是来自需求结构的变化——入境游客对价格的敏感度较低，且更倾向于选择高端品牌。

这里有一个值得深思的问题：如果入境旅游增速放缓，高端酒店是否也会被迫进入降价模式？报告没有给出明确答案，但这一假设值得投资者纳入风险情景。

UBS同时发布了航空业的数据，这为酒店行业的观察提供了重要的交叉验证。6月1日至7日，国内机票价格同比上涨23\%，但旅客量同比下降13\%，降幅比前一周的7\%进一步扩大。国内客运收入增速放缓至7\%，远低于前一周的18\%。

这是一个典型的“涨价抑制需求”案例。机票价格上涨23\%直接导致旅客量下滑13\%，价格弹性系数约为-0.57，说明国内航空需求对价格高度敏感。这与中端酒店降价换

[... middle omitted ...]

P活跃用户数据背后的品牌分化

UBS还引用了UBS Evidence Lab的酒店预订APP周活跃用户数据，涵盖首旅如家、华住会、锦江酒店和亚朵。这张图表的趋势值得单独审视：华住会的活跃用户规模明显领先，且增长曲线最为陡峭，而亚朵虽然基数最小，但增速也较为可观。

这一数据与RevPAR的分化之间存在怎样的关联？报告没有展开分析，但这是一个重要的观察线索。APP活跃用户反映了品牌直销能力和会员粘性，而直销渠道的占比直接影响酒店集团的定价权和佣金成本。如果华住会等头部品牌能够通过会员体系锁定更多直客，它们在面对需求波动时的定价灵活性就会更强，也更有能力在不降价的情况下维持入住率。

报告没有回答的问题是：APP活跃用户的增长是否直接转化为更高的RevPAR和利润率？以及，不同品牌之间的用户粘性差异是否正在放大行业的分化？这些都是需要进一步拆解的二阶影响。

基于以上分析，我们建议将观察框架从“行业整体复苏进度”调整为“定价权竞争格局”。具体而言，可以关注以下几个维度：

\subsection{[R032] 摩根斯坦利：AI光互连的下一波赢家不在光模块，而在电路板}
{\small\color{refgray}归类：AI与科技：资本开支结构性集中，光互连与PCB成为下一波赢家}

当市场将全部注意力集中在GPU、网络芯片和光模块供应商时，一份来自摩根斯坦利的深度报告揭示了一个被低估的结构性机会：光收发器PCB。这并非简单的配套逻辑，而是一个正在经历“量价齐升”式重估的细分市场。

这份报告的核心判断是：AI投资周期正进入新阶段，光互连的重要性正与算力本身持平。当超大规模云厂商将AI集群从数万GPU扩展至数十万GPU时，连接这些系统的光收发器数量正在指数级增长。但摩根斯坦利团队看到的，不仅是收发器数量的增长，更是每块电路板价值的跃迁。

这意味着，传统意义上被视为“被动配套”的PCB制造商，正在成为AI基础设施支出的结构性受益者。而其中最值得关注的，并非当前市场份额最大的玩家，而是一家正在加速追赶的挑战者。

1. 光收发器PCB市场正在以83\%的复合增速成为全球PCB行业最强劲的增长引擎

摩根斯坦利预测，全球AI光收发器PCB市场规模将从2025年的约6.2亿美元增长至2028年的约37.7亿美元，复合年增长率高达83\%。这一增速远超同期光收发器出货量的60\%复合增速，背后是“量”与“价”双重驱动。

从绝对规模来看，这一市场正在快速逼近全球最大的HDI PCB需求方——苹果。报告估算，苹果每年对HDI PCB的需求约为30至35亿美元。而到2027年，AI光收发器PCB市场将成长至与苹果需求相当的规模，并在2028年超越它。

这一对比的意义不仅在于市场体量。多年来，苹果一直是HDI PCB行业技术进步、产能投资和盈利能力的核心驱动力。而现在，AI光互连正在成为下一个主要需求引擎，这可能重塑PCB行业的竞争格局。

根据Prismark的估算，全球HDI PCB市场约为162亿美元。到2028年，仅AI光互连一项就能创造相当于当前全球HDI市场约23\%的增量需求。这不是一个边缘机会，而是PCB行业未来几年最清晰的结构性增长主线。

2. 从400G到1.6T的迁移正在推动每块PCB的价值量翻倍以上

光收发器PCB市场的增长并非仅由出货量驱动，更关键的是单模块PCB价值的显著提升。摩根斯坦利详细拆解了三个驱动因素：

第一，CCL材料升级。从400G模块常用的M6级CCL，向800G和1.6T模块所需的M7或M8级CCL迁移，直接推高了PCB的单价。

第二，层数增加。400G模块的PCB层数为10至12层，800G升至12至14层，1.6T则达到14至16层，设计复杂度的提升带来了更高的附加值。

第三，制造工艺升级。从传统的HDI工艺向mSAP工艺迁移，是PCB价值跃升的最关键因素。mSAP工艺能够实现更密集的导线布局，原本主要应用于智能手机主板和摄像头模组。苹果从2017年的iPhone 8/X开始率先采用mSAP技术，而如今，这一技术正在成为1.6T光收发器PCB的标配。

综合这些因素，摩根斯坦利估算，单块PCB的平均售价将从400G时代的约10美元，提升至1.6T时代的约25美元，增长约2.5倍。更重要的是，随着技术壁垒的提高，毛利率将从20\%至30\%跃升至40\%至50\%以上，接近ABF载板的盈利水平。

这一变化意味着，PCB供应商正在从“配套角色”转变为“技术驱动型参与者”，其盈利能力将随着产品迭代而结构性改善。

3. 现有供应商将继续受益，但真正的超额收益来自市场份额的再分配

在现有供应商中，摩根斯坦利维持对欣兴电子和深南电路的看好评级，但认为两者已部分反映在估值中。真正具有超额收益潜力的，是一家正在加速追赶的公司——臻鼎科技。

报告指出，光收发器PCB市场正在经历供应格局的变化。随着市场快速扩张，现有产能难以满足需求，客户需要寻找更多合格的供应商。臻鼎科技凭借其在苹果iPhone供应链中积累的mSAP工艺经验，正处于切入这一市场的绝佳位置。

摩根斯坦利估算

[... middle omitted ...]

路的25倍，提供了更具吸引力的风险回报比。

4. 报告尚未完全回答的关键问题：CPO的长期威胁与产能扩张的节奏

摩根斯坦利在报告中坦率地讨论了CPO技术的潜在影响，认为其是“根本性的架构转变”和“对可插拔收发器的长期风险”。但报告给出的判断是，CPO的大规模采用在2027至2028年之前可能性有限，主要受制于制造良率、热管理复杂性、成本、生态系统和可维护性等因素。

这里存在一个尚未完全解答的问题：如果CPO的进展快于预期，当前对PCB供应商的投资逻辑将如何调整？报告认为，在3.2T过渡期（预计2028年），可插拔收发器和CPO将共存，但这一判断依赖于对技术成熟度的假设。这些假设仍需持续验证。

另一个值得关注的问题是产能扩张的节奏。报告指出，现有行业产能不足以满足预期需求，这正是新进入者获得份额的机会。但产能扩张本身需要大量资本支出，且存在良率爬坡的不确定性。臻鼎科技和台光电子等公司正在积极扩产，但能否以预期的速度和效率完成产能建设，将直接影响其份额提升的节奏。

\subsection{[R033] 摩根斯坦利：市场低估的不是AI投资放缓，而是日本线缆企业供给侧的差异化定价权}
{\small\color{refgray}归类：AI与科技：资本开支结构性集中，光互连与PCB成为下一波赢家}

摩根斯坦利：市场低估的不是AI投资放缓，而是日本线缆企业供给侧的差异化定价权

当市场目光聚焦于宏观不确定性和AI半导体板块的剧烈回调时，一份来自摩根斯坦利的日本线缆行业研报，给出了一个与近期股价表现截然相反的判断：基本面依然强劲，且真正的结构性机会可能被市场误读了。

这份由分析师Yu Shirakawa主笔的报告，覆盖了住友电工、藤仓和古河电工三家日本线缆巨头。报告发布之际，三家公司股价均经历了不同程度的回调——藤仓一个月内下跌近40\%，住友电工跌超8\%，即便是涨幅最为凌厉的古河电工也回落了6\%。然而，摩根斯坦利却在此刻上调了住友电工的目标价，同时维持对古河电工的“超配”评级，并重申整个行业的“吸引力”评级。

这并非简单的逆势看多。其核心逻辑在于：市场当前对股价回调的解读，可能过度归因于宏观风险或AI投资周期见顶，而真正决定这三家公司未来两年业绩走向的关键变量——供给侧的产能瓶颈、产品结构的升级以及由此带来的议价权重塑——尚未被充分定价。

本文基于这份研报，提炼出三个层次的核心洞察：为什么AI投资不是问题、三家公司的差异化究竟在哪、以及市场尚未完全消化的风险与机会。

1. 超大规模云厂商的资本开支正在加速，而非减速，这是行业最坚实的底部

市场近期对AI相关股票的担忧，很大程度上源于对资本开支周期见顶的恐惧。然而，摩根斯坦利援引其美国IT团队的更新预测指出，2026年11家超大规模云厂商的资本开支同比增速预期已从此前的63\%大幅上修至87\%。

这一数字的意义在于，它直接回答了“AI投资是否已经放缓”这一关键问题。答案是否定的。不仅如此，报告明确指出，约75\%的资本开支将直接投向AI基础设施，包括数据中心、GPU集群和电力设施。这意味着，对于处于产业链上游的线缆和光通信企业而言，订单的可见度和持续性远高于市场普遍预期。

更重要的是，推动这一轮投资的结构性力量正在从训练转向推理。随着生成式AI的普及，低延迟和数据本地化的需求正在催生分布式AI基础设施的扩张——包括区域数据中心和边缘节点。这种变化直接拉动了对高速通信基础设施，特别是光连接器的需求。报告中提到，这一趋势已经体现在日本对亚洲出口数据的持续增长中。

所以，市场对AI投资放缓的担忧，至少在 财年的时间维度上，缺乏数据支撑。真正的风险不在于需求端，而在于供给端——企业能否及时扩张产能以满足这一波需求。

2. 三家公司的分化在于：产能弹性、产品组合和与共识预期的差距

摩根斯坦利对三家公司给出了截然不同的评级和估值判断，其筛选框架聚焦于三个维度：盈利增长能力、产能扩张支撑、以及与市场共识的差距。这构成了理解其投资排序的核心逻辑。

*古河电工是唯一获得“超配”评级的标的。** 报告预计其未来三年的营业利润复合年增长率高达42.1\%，是三家中最高的。增长的主要驱动力来自超高芯数光缆、MT插芯和液冷模块。值得注意的是，摩根斯坦利对古河电工F3/27的营业利润预测为1000亿日元，高于公司自身950亿日元的计划。这种“超越管理层指引”的信心，来源于对液冷模块放量的乐观预期——预计该业务收入将从F3/27的500亿日元跃升至F3/29的3000亿日元。如果这一预测成立，古河电工的盈利结构将发生质变，从传统的线缆制造商转型为AI基础设施的关键组件供应商。

*住友电工被上调目标价，但维持“等权重”评级。** 上调目标价的理由在于，分析师对光通信相关产品的扩张速度有了更强信心，特别是在连续波激光二极管领域，住友电工凭借磷化铟平台的垂直整合生产拥有高市场份额。然而，评级未上调的原因在于，市场共识在近期财报后已经大幅上修，股价进一步上行的空间有限。这揭示了一个微妙的投资逻辑：好公司不等于好股票，当预期已经充分反映在价格中时，超额收益的来源就变得稀薄。

*藤仓

[... middle omitted ...]

O技术的推广。摩根斯坦利预计，从F3/27开始，英伟达下一代的Rubin GPU很可能将采用CPO技术，即将光学组件与半导体芯片集成在同一封装内。

这一技术路线的改变，将带来单位光学组件用量的指数级增长。在传统架构下，一个交换机最多配置64个端口。而采用CPO后，端口密度可以翻倍至128个800Gb/s端口，甚至在大型配置中实现512个800Gb/s端口。关键在于，CPO架构取消了可插拔模块，光纤芯线将直接作为外部布线出现。这意味着，每台交换机所需的光纤芯数和连接点数量将大幅增加。

报告判断，住友电工很可能是CPO趋势的最大潜在受益者，因为其开发的CW-LD器件是CPO方案的核心组件。这一判断目前仍处于早期验证阶段，但如果CPO的渗透速度超出预期，住友电工的估值逻辑将需要被重估。

4. 报告尚未完全回答的关键问题：涨价带来的利润率改善何时兑现？

尽管摩根斯坦利对行业基本面持乐观态度，但报告本身也坦诚地指出了几个尚未被验证的关键变量，其中最重要的一项是价格传导。

\subsection{[R034] Bernstein：AI购物离“自主决策”还差一个支付环节}
{\small\color{refgray}归类：未被主题摘要引用}

过去六个月，从沃尔玛宣布与ChatGPT合作开始，“Agentic Shopping”成为美国零售业最热的关键词。但喧嚣之下，一个根本问题始终悬而未决：AI到底能在多大程度上替代人类的购物决策？

Bernstein（Bernstein）最近发布了一份极具实操性的研报。研究团队没有停留在概念推演，而是直接上手测试了五款AI工具——ChatGPT、Gemini、Claude三大基础模型，以及亚马逊Alexa和沃尔玛Sparky两个零售原生系统。测试任务简单到近乎朴素：为一家人买一盒牛奶，为一周备餐搭建购物篮。

结果却揭示了一个被市场严重低估的核心矛盾。AI在“发现商品”环节已经展现出令人印象深刻的潜力，但在“完成交易”这一最终步骤上，所有工具无一例外地败下阵来。这份报告的核心判断并不在于AI购物何时到来，而在于：真正决定AI购物能否从“玩具”变成“工具”的关键，并非模型能力，而是支付环节的深度嵌入。

这不是一个技术问题，而是一个商业模式和竞争格局的问题。谁先打通支付闭环，谁就掌握了下一代零售的入口。

1. 从“发现”到“购买”的断层，是所有AI购物工具的阿克琉斯之踵

Bernstein的测试设计精准地抓住了购物漏斗的六个关键环节：品类识别、SKU识别、定价、总成本估算、加入购物车、支付。结果一目了然：在“加入购物车”和“支付”这两个环节，所有工具都亮起了红灯。

基础模型（ChatGPT、Gemini、Claude）在信息检索和推荐上表现尚可，但一旦涉及交易执行，它们立刻暴露出“手不够长”的致命缺陷。它们可以告诉你“这盒牛奶在Whole Foods卖4.79美元”，但无法帮你下单。它们依赖网络爬虫和第三方文章来获取信息，导致SKU识别和实时定价经常出现偏差。最典型的是，当被问及具体价格时，Gemini需要用户进一步追问才能给出精确的SKU级信息，而Claude则直接承认“无法直接搜索购物和配送信息”。

零售原生系统（Alexa和Sparky）则呈现出截然不同的面貌。它们拥有实时库存和定价数据的直接访问权，因此在SKU识别和定价准确性上显著优于基础模型。Alexa能够精准推荐“365 by Whole Foods Market Organic Whole Milk, 64 oz”，并给出“4.79美元，今天下午1点到3点送达”的完整信息。Sparky甚至可以直接将商品加入购物车。

但问题出在了最后一步。Alexa无法处理Amazon Fresh和Whole Foods的生鲜产品，当用户要求将商品加入购物车时，它只能提供一个外部链接。Sparky虽然能完成从发现到建篮的全流程，但支付环节仍需用户手动操作。没有一个工具能够实现“无监督”的自主购买。

这意味着，当前所有AI购物工具的“自主性”都停留在“推荐”层面，而非“决策”层面。它们可以帮你找到商品，但无法替你完成交易。这个断层的背后，是AI系统与零售生态之间尚未弥合的鸿沟。

Bernstein的测试揭示了一个核心张力：基础模型和零售原生系统各自拥有对方无法企及的优势，但也面临着对方不存在的短板。

基础模型（ChatGPT、Gemini、Claude）的“自由”是它们最大的武器。它们可以跨越不同零售商、不同平台，进行横向比较和跨生态推荐。当用户需要“一盒64盎司的有机全脂牛奶”时，ChatGPT能同时推荐沃尔玛的Great Value、Whole Foods的365、以及Organic Valley等多个选项，并给出价格区间和配送时间估算。这种“游牧式”的搜索能力，让用户能够获得最全面的市场信息。

然而，这种自由是有代价的。由于缺乏对零售商结构化目录数据的直接访问，基础模型的信息来源不可靠，定价和库存信息常常滞后或错误。它们本质上是一个“高级搜索引

[... middle omitted ...]

个能够同时解决“自由”和“精准”问题的玩家。** 这需要基础模型与零售商之间建立更深的合作关系，比如ChatGPT与Uber Eats的整合，或者Claude与Instacart的潜在合作。但目前的合作深度，还远不足以打通支付闭环。

3. 搭建购物篮：设计哲学的分野暴露了AI的“理解力”与“执行力”之间的鸿沟

如果说购买单一商品测试的是AI的“执行力”，那么搭建一周购物篮测试的则是AI的“理解力”。Bernstein设计的任务包含一个五天的家庭餐单、一个四口之家的需求、一个10万美元的年收入约束、以及一个坚果过敏的限制条件。

在这个任务中，基础模型的表现令人惊讶。ChatGPT、Gemini和Claude都能生成一份结构完整、逻辑自洽的购物清单。它们能够理解“家庭四口人”意味着需要估算食材数量，能够识别“坚果过敏”并自动排除含有坚果的产品，甚至能够根据“年收入10万美元”的家庭预算，推荐性价比较高的商品。这份清单，更像是一份“购物计划”，而非简单的“购物列表”。

\section{Disclaimer}
{\scriptsize\color{refgray}
This community is a paid learning and information-sharing community created for general educational purposes only. The membership fee is charged solely for access to community discussions, educational content, organization, and operational costs. It is not an advisory fee, management fee, performance fee, brokerage fee, or compensation for personalized financial, investment, legal, tax, or professional advice.\par
I participate and share information solely as an individual and for educational discussion among community members. I am not acting as a financial adviser, investment adviser, broker, dealer, portfolio manager, legal adviser, tax adviser, accountant, fiduciary, or any other licensed professional adviser. Nothing shared in this community constitutes, and should not be interpreted as, financial advice, investment advice, legal advice, tax advice, a recommendation, solicitation, offer, endorsement, instruction, or guarantee to buy, sell, hold, short, trade, or invest in any security, cryptocurrency, fund, derivative, strategy, or other financial product.\par
All content shared in this community, including messages, discussions, links, files, charts, examples, market commentary, watchlists, AI-generated or AI-polished summaries, and third-party materials, is general, impersonal, and educational in nature. It does not take into account any individual's financial situation, investment objectives, risk tolerance, investment horizon, liquidity needs, tax status, legal restrictions, or suitability requirements. No content should be relied upon as suitable for any specific person, account, portfolio, or financial decision.\par
Information shared in this community may be incomplete, inaccurate, outdated, speculative, AI-generated, AI-edited, or based on third-party internet sources that have not been independently verified. I make no representation or warranty as to the accuracy, completeness, timeliness, reliability, or suitability of any information shared.\par
Financial markets involve substantial risk, including the possible loss of principal. Past performance, historical data, hypothetical examples, simulated results, backtests, personal opinions, or educational examples do not guarantee future results. Each member is solely responsible for conducting their own research, exercising independent judgment, and making their own decisions. Before making any financial, investment, legal, or tax decision, members should consult qualified and licensed professionals.\par
Participation in this community, payment of any membership fee, or communication with me or other members does not create any adviser-client, fiduciary, professional, contractual, agency, partnership, employment, or other advisory relationship. By joining or participating in this community, you acknowledge and agree that you are solely responsible for your own decisions, actions, transactions, profits, losses, risks, and consequences, and that I shall not be liable for any direct, indirect, incidental, consequential, financial, legal, tax, or other loss or damage arising from or related to any content, discussion, information, or material shared in this community.\par
}
\end{document}
